\chapter{Linearization of symbols, K-theory and Bott periodicity}\label{ch:bott}
\chaptermark{Linearization, K-theory, Bott periodicity}
\providecommand{\Ginf}{G^{-\infty}}
\providecommand{\GinfT}{G^{-\infty}_{\mathcal{T}}}
\providecommand{\GT}{G^{0}_{\mathcal{T}}}
\providecommand{\PsiT}[1]{\Psi^{#1}_{\mathcal{T}}}
\providecommand{\Cp}{\mathcal{C}^\infty_+}
\providecommand{\Kc}{K_{\mathrm{c}}}
\providecommand{\wind}{\operatorname{wind}}
\providecommand{\Iso}{\operatorname{Iso}}
\providecommand{\rank}{\operatorname{rank}}
\providecommand{\Gres}{G^{-\infty}_{\mathrm{iso}}}
\providecommand{\Psiisoph}[1]{\Psi^{#1}_{\mathrm{iso,ph}}}
\providecommand{\OpW}{\operatorname{Op}^{\mathrm{W}}}
\providecommand{\indiso}{\operatorname{ind}_{\mathrm{iso}}}
\providecommand{\Indsl}[1]{\operatorname{Ind}^{#1}_{\mathrm{sl}}}
\providecommand{\ann}{\mathfrak{a}}

The index of an elliptic operator depends only on the homotopy class of its
symbol. To compute it we need to understand such homotopy classes, and the
simplest non-trivial case already contains the key idea: loops of invertible
matrices, which are the symbols of Toeplitz operators on the circle. This chapter
shows that every such loop can be deformed into an explicit \emph{normal form},
reads off the index from the normal form, and deduces Bott periodicity, the
basic theorem of K-theory, from the resulting index formula.

In \cref{ch:index} the index of an elliptic operator was shown to be a
homotopy invariant, the group $\Ginf(Y;E)$ of invertible perturbations of the
identity by smoothing operators was introduced as a model for odd K-theory, and
the Toeplitz algebra on the circle was set up. This chapter combines these
ingredients. The central technical result is a normal form, due to Atiyah, for
loops of invertible matrices: after stabilization and deformation every
(pointed) loop can be brought to the form
$\pi_-e^{-i\theta}+\pi_0+\pi_+e^{i\theta}$ with commuting idempotents
$\pi_\pm,\pi_0$. Applied to the symbols of Toeplitz operators this
\emph{linearization} reduces index computations to an explicit model.
As consequences we obtain:
\begin{itemize}
\item the index formula for Toeplitz and pseudodifferential operators on the circle;
\item the families index of elliptic Toeplitz families, with values in $K^0$, and the
  fact that it gives an isomorphism $K^{-2}(X)\cong K^0(X)$ (Bott periodicity in its
  first form);
\item the weak contractibility of the stabilized group $\GT$ of invertible Toeplitz
  operators, hence a \emph{classifying sequence}
  $\GinfT\to\GT\to\GT/\GinfT$ for K-theory;
\item Bott periodicity in its usual form, $\Kc^0(\R^2\times X)\cong K^0(X)$, and the
  homotopy groups of $\Ginf$: $\pi_k(\Ginf)=\Z$ for $k$ odd and $0$ for $k$ even;
\item a second, independent, realization of the index isomorphism through the
  isotropic calculus of \cref{ch:isotropic}: the circle at infinity of the phase
  plane $\R^2$ replaces the circle, the eigenprojections of the harmonic oscillator
  replace the projections onto Fourier modes, and the annihilation operator replaces
  the Toeplitz shift;
\item the \emph{isotropic--semiclassical quantization maps}
  $\Kc^{-1}(\R^{2n}\times X)\to\Kc^{-1}(X)$ and $\Kc^0(\R^{2n}\times X)\to\Kc^0(X)$,
  obtained by quantizing compactly supported symbols with values in invertible
  matrices or in projections, which realize Bott periodicity and are used in
  \cref{ch:semiclassical-index};
\item complex and symplectic vector bundles and the Thom isomorphism;
\item the analytic index map $\Kc^0(T^*Z)\to\Z$, the starting point of the index
  theorem.
\end{itemize}
K-theory is defined from scratch in \cref{bott:sec:K}; we follow
\cite{Atiyah-K} for the topological background. The main line uses Toeplitz operators
on the circle; compare \cite{Pressley-Segal} for the relation between loop groups and
Toeplitz operators. The isotropic version (\cref{bott:sec:isotropic}) runs parallel to
it and the two index maps are compared in \cref{bott:prop:iso-toeplitz}.

\section{Conventions and background}\label{bott:sec:background}

Throughout, $X$ denotes a compact manifold, possibly with boundary (we need
$[0,1]\times X$), and all maps and homotopies are smooth. The circle is
$\Sphere=\R/2\pi\Z$ with variable $\theta$; we write $z=e^{i\theta}$, so that
$\Sphere$ is also the unit circle in $\C$, and we use the point $1\in\Sphere$
($\theta=0$) as base point. $M(N,\C)$ is the algebra of $N\times N$ matrices.
For maps from a product we use the usual notion of smoothness; a smooth map
into a Fr\'echet space $F$ (such as $\Psido{0}(\Sphere)$ with its standard topology,
see \cref{ch:conormal,ch:manifolds}) is one all of whose derivatives exist and
are continuous.

Two smooth homotopies $h,h'\colon[0,1]\times X\to W$ with $h_1=h'_0$ can be
concatenated to a smooth homotopy after reparametrizing each of them by a smooth
non-decreasing surjection $[0,1]\to[0,1]$ which is constant near $0$ and $1$.
We use this tacitly whenever homotopies are composed.

\subsection{Facts recalled from \cref{ch:index}}
Let $Y$ be a compact manifold of positive dimension, $E\to Y$ a Hermitian vector
bundle, and fix a smooth positive density on $Y$; write $L^2(Y;E)$ for the
corresponding Hilbert space.
\begin{enumerate}[label=(B\arabic*),leftmargin=*]
\item\label{bott:B:Ginf} $\Psido{-\infty}(Y;E)$, the smoothing operators (kernels in
  $\CI(Y^2;\Hom(E)\otimes\Omega_R)$), is a Fr\'echet algebra and
  \begin{equation*}
    \Ginf(Y;E)=\{\Id+A;\ A\in\Psido{-\infty}(Y;E),\ \Id+A \text{ invertible on } L^2(Y;E)\}.
  \end{equation*}
  The inverse of $\Id+A\in\Ginf(Y;E)$ is of the form $\Id+B$ with
  $B\in\Psido{-\infty}(Y;E)$; $\Ginf(Y;E)$ is an open subset of
  $\Id+\Psido{-\infty}(Y;E)$, a topological group with smooth multiplication and
  inversion, and it is connected.
\item\label{bott:B:finite} (Finite rank approximation.) Given a finite dimensional
  subspace $V_0\subset\CI(Y;E)$ there is an increasing sequence of finite dimensional
  subspaces $V_0\subset V_1\subset V_2\subset\dots\subset\CI(Y;E)$ (for instance
  $V_0$ plus the span of the first $j$ eigensections of a Laplacian) whose orthogonal
  projections $\Pi_j$ satisfy $\Pi_jA\Pi_j\to A$ in $\Psido{-\infty}(Y;E)$,
  uniformly with all derivatives on smooth families $A$ parametrized by a compact
  manifold. For a finite dimensional subspace $V$ with orthogonal projection $\Pi_V$
  we identify $\GL(V)$, and so $\GL(N,\C)$ after choosing an orthonormal basis, with
  the subgroup $\{\Id-\Pi_V+g\Pi_V;\ g\in\GL(V)\}\subset\Ginf(Y;E)$. Enlarging $V$ to
  $V'\supset V$ corresponds to the stabilization $g\mapsto g\oplus\Id$.
\item\label{bott:B:K1} $K^{-1}(X)=[X;\Ginf(Y;E)]$, the set of smooth homotopy classes of
  smooth maps, is an abelian group under pointwise multiplication, independent of
  the choice of $(Y,E)$.
\item\label{bott:B:wind} For a smooth loop $b\colon\Sphere\to\Ginf(Y;E)$ the number
  \begin{equation*}
    \wind(b)=\frac1{2\pi i}\int_0^{2\pi}\Tr\Big(b^{-1}\frac{db}{d\theta}\Big)\,d\theta
  \end{equation*}
  is an integer, the winding number of the Fredholm determinant along $b$. For
  $b$ with values in $\GL(N,\C)$ it is the winding number of $\det b$, since
  $d\log\det b=\Tr(b^{-1}db)$. It is invariant under homotopy of loops.
\item\label{bott:B:toeplitz} The Szeg\H{o} projector $S$, which removes the negative Fourier
  modes, is a self-adjoint projection in $\Psido{0}(\Sphere)$ with principal symbol
  equal to $1$ on $\Sphere_+=\{\xi>0\}$ and $0$ on $\Sphere_-=\{\xi<0\}$ in
  $S^*\Sphere=\Sphere_+\sqcup\Sphere_-$. Its range on $\CI(\Sphere)$ is
  $\Cp(\Sphere)$, the smooth functions with only non-negative Fourier modes,
  and on $L^2$ it is $L^2_+(\Sphere)$. The Toeplitz algebra is
  \begin{equation*}
    \PsiT{0}(\Sphere;\C^N)=\{A\in\Psido{0}(\Sphere;\C^N);\ A=SAS\},
  \end{equation*}
  acting on $\Cp(\Sphere;\C^N)$; $\PsiT{-\infty}(\Sphere;\C^N)=S\Psido{-\infty}(\Sphere;\C^N)S$.
  The (Toeplitz) symbol $\sigma_0(A)\in\CI(\Sphere;M(N,\C))$ is the restriction of
  the principal symbol to $\Sphere_+\cong\Sphere$; it is multiplicative and for
  $a\in\CI(\Sphere;M(N,\C))$ the operator $T_a=SaS$ (multiplication by $a$, then
  projection) lies in $\PsiT{0}(\Sphere;\C^N)$ with $\sigma_0(T_a)=a$. A constant
  matrix commutes with $S$, so $T_c=c$. $A$ is \emph{elliptic} if $\sigma_0(A)$ is
  invertible; then $A$ has a parametrix in $\PsiT{0}$ modulo $\PsiT{-\infty}$, is
  Fredholm on $\Cp(\Sphere;\C^N)$ and on $L^2_+(\Sphere;\C^N)$ with null spaces
  in $\Cp$, and its index depends only on the homotopy class of $\sigma_0(A)$ in
  $\CI(\Sphere;\GL(N,\C))$. If an elliptic $A$ is invertible, then
  $A^{-1}\in\PsiT{0}(\Sphere;\C^N)$.
\item\label{bott:B:shifts} The shifts $U=T_{z}$ and $L=T_{z^{-1}}$ act by
  $U(e^{ik\theta})=e^{i(k+1)\theta}$, $L(e^{ik\theta})=e^{i(k-1)\theta}$ for $k\ge1$ and
  $L(1)=0$. Thus $LU=\Id$, $UL=\Id-P_0$, $\ind U=-1$ and $\ind L=1$. Here, and below,
  $P_k$ is the orthogonal projection onto the span of $e^{ij\theta}$, $0\le j\le k$
  (acting componentwise on vector valued functions). For a smooth family $K$ of
  elements of $\PsiT{-\infty}$ parametrized by a compact manifold, $KP_k\to K$ and
  $P_kK\to K$ in $\PsiT{-\infty}$ uniformly as $k\to\infty$.
\item\label{bott:B:lift} (Lifting homotopies.) If $A\colon[0,1]\times X\to
  \PsiT{0}(\Sphere;\C^N)$ is a smooth family of elliptic operators and $A_0(x)$ is
  invertible for all $x\in X$, then there is a smooth family
  $B\colon[0,1]\times X\to\PsiT{-\infty}(\Sphere;\C^N)$ with $B_0=0$ such that
  $A_t(x)+B_t(x)$ is invertible for all $(t,x)$.
\end{enumerate}

A smooth family of invertible elliptic Toeplitz operators has smoothly varying
inverses; we prove this in \cref{bott:lem:inverse} below.

\subsection{Elementary deformations of loops}

A \emph{pointed loop} (over $X$) is a smooth map
$a\colon X\times\Sphere\to\GL(N,\C)$ with $a(x,1)=\Id$ for all $x$; a pointed
homotopy is a smooth family of pointed loops. Loops are \emph{stabilized} by
$a\mapsto a\oplus\Id_{M-N}$, and $a,b$ are \emph{stably pointed homotopic},
$a\sim b$, if $a\oplus\Id$ and $b\oplus\Id$ are pointed homotopic in some
$\GL(M,\C)$. The following three lemmas are used constantly.

\begin{discussion}
  Stabilization is what makes loops manageable. In a fixed $\GL(N,\C)$ two loops
  $a$, $b$ need not commute, and there is no direct reason for $ab$ and $ba$
  to be homotopic. After adding a block of identity matrices there is room to rotate
  one loop into the new block (\cref{bott:lem:rotation}), where it commutes with
  the other. All the deformations below use this freedom to add blocks.
\end{discussion}

\begin{lemma}\label{bott:lem:normalize}
  Let $h\colon[0,1]\times X\times\Sphere\to\GL(N,\C)$ be smooth and set
  \begin{equation*}
    \hat h_t(x,\theta)=h_t(x,1)^{-1}h_t(x,\theta).
  \end{equation*}
  Then $\hat h$ is a pointed homotopy; $\hat h_t=h_t$ whenever $h_t$ is pointed, and $\hat h_t=\Id$ whenever
  $h_t(x,\theta)$ is independent of $\theta$.
\end{lemma}
\begin{proof}
  Immediate from the definition.
\end{proof}

\begin{lemma}[Rotation]\label{bott:lem:rotation}
  For pointed loops $a,b$ with values in $\GL(N,\C)$, $\diag(ab,\Id_N)$ and
  $\diag(a,b)$ are pointed homotopic in $\GL(2N,\C)$. Consequently
  $ab\sim\diag(a,b)\sim ba$, $\diag(a_1\cdots a_r,\Id,\dots,\Id)\sim
  \diag(a_1,\dots,a_r)$, and block sums of pointed loops are commutative up to
  $\sim$.
\end{lemma}
\begin{proof}
  Let $R_\tau=\begin{pmatrix}\cos\tau\Id_N&-\sin\tau\Id_N\\\sin\tau\Id_N&\cos\tau\Id_N\end{pmatrix}$
  and $h_\tau=\diag(a,\Id)R_\tau\diag(b,\Id)R_\tau^{-1}$, $\tau\in[0,\pi/2]$. Then
  $h_0=\diag(ab,\Id)$, $h_{\pi/2}=\diag(a,\Id)\diag(\Id,b)=\diag(a,b)$, and
  $h_\tau(x,1)=\Id$. The remaining statements follow by iteration and by applying
  the same argument with the roles of the blocks exchanged (the constant
  permutation matrix exchanging two blocks is joined to $\Id$ by $R_\tau$).
\end{proof}

\begin{lemma}[Conjugation]\label{bott:lem:conjugation}
  Let $a$ be a pointed loop in $\GL(N,\C)$ and $c\colon X\to\GL(N,\C)$ smooth
  (independent of $\theta$). Then $cac^{-1}\sim a$; if $c$ is constant then
  $cac^{-1}$ and $a$ are pointed homotopic without stabilization.
\end{lemma}
\begin{proof}
  If $c$ is constant, join it to $\Id$ by a path $c_t$ in the connected group
  $\GL(N,\C)$; then $c_tac_t^{-1}$ is a pointed homotopy. In general,
  $h_\tau=\diag(c,\Id)R_\tau\diag(c^{-1},\Id)R_\tau^{-1}$ is a $\theta$-independent
  path from $\Id$ to $\diag(c,c^{-1})$, and
  $h_\tau\diag(a,\Id)h_\tau^{-1}$ is a pointed homotopy from $\diag(a,\Id)$ to
  $\diag(c,c^{-1})\diag(a,\Id)\diag(c^{-1},c)=\diag(cac^{-1},\Id)$.
\end{proof}

\section{Vector bundles and K-theory}\label{bott:sec:K}

All vector bundles are smooth complex vector bundles of finite rank; $\C^N$ also
denotes the trivial bundle. For bundles $E,F$ over the same base we write
$E\cong F$ if there is a smooth bundle isomorphism, i.e.\ a smooth map of total spaces
which is a linear isomorphism $E_x\to F_x$ on each fibre.

\subsection{Bundles}
\begin{lemma}\label{bott:lem:bundles}
  Let $X$ be compact.
  \begin{enumerate}[label=(\roman*)]
  \item Every vector bundle $E$ over $X$ is isomorphic to the range of a smooth
    family of orthogonal projections $p\colon X\to M(N,\C)$; in particular there
    is a bundle $F$ with $E\oplus F\cong\C^N$.
  \item If $\pi\colon X\to M(N,\C)$ is a smooth family of idempotents, $\pi^2=\pi$,
    then $\Ran\pi=\{(x,v);\ v\in\Ran\pi(x)\}$ is a smooth vector bundle.
  \end{enumerate}
\end{lemma}
\begin{proof}
  (ii) The rank of $\pi(x)$ equals $\tr\pi(x)$, so is locally constant. If
  $v_1,\dots,v_r$ is a basis of $\Ran\pi(x_0)$ then $\pi(x)v_1,\dots,\pi(x)v_r$ are
  linearly independent for $x$ near $x_0$, hence a smooth local frame of $\Ran\pi$.

  (i) Choose open sets $U_1,\dots,U_s$ covering $X$ over which $E$ has trivializations
  $\tau_i\colon E|_{U_i}\to U_i\times\C^n$, and smooth $\psi_i\ge0$ with
  $\supp\psi_i\subset U_i$ and $\sum_i\psi_i>0$. Then
  \begin{equation*}
    \iota(v)=(\psi_1(x)\tau_1(v),\dots,\psi_s(x)\tau_s(v)),\quad v\in E_x,
  \end{equation*}
  is a smooth fibrewise injective bundle map $E\to\C^{N}$, $N=sn$. The orthogonal projection
  $p(x)$ onto $\iota(E_x)$ depends smoothly on $x$ (locally it is given by the
  Gram--Schmidt process applied to a smooth frame), so $E\cong\Ran p$, and
  $F=\Ran(\Id-p)$ is a complement.
\end{proof}

\begin{lemma}[Homotopy invariance of bundles]\label{bott:lem:bundle-homotopy}
  Let $W$ be a manifold.
  \begin{enumerate}[label=(\roman*)]
  \item If $E$ is a vector bundle over $[0,1]\times W$, then
    $E|_{\{0\}\times W}\cong E|_{\{1\}\times W}$ as bundles over $W$.
  \item If $\varpi\colon V\to W$ is a real vector bundle (for instance $V=\R^n\times W$
    or $V=T^*Z$) and $E$ is a vector bundle over the total space $V$, then
    $E\cong\varpi^*(E|_{W})$, where $W$ is identified with the zero section.
  \end{enumerate}
\end{lemma}
\begin{proof}
  Choose a connection on $E$ (patching local trivial connections with a partition
  of unity). In (i), parallel transport along the curves $t\mapsto(t,w)$ gives linear
  isomorphisms $E_{(0,w)}\to E_{(1,w)}$ which depend smoothly on $w$ by smooth
  dependence of solutions of linear ordinary differential equations on parameters.
  In (ii), parallel transport along $t\mapsto t\zeta$, $t\in[0,1]$, from the zero
  section to $\zeta\in V$ gives a smooth isomorphism $\varpi^*(E|_W)\to E$.
\end{proof}

\subsection{The group \texorpdfstring{$K^0$}{K0}}

\begin{definition}\label{bott:def:K0}
  For compact $X$, $K^0(X)$ is the set of equivalence classes of pairs $(E_+,E_-)$ of
  vector bundles over $X$, where $(E_+,E_-)\sim(F_+,F_-)$ if there is a bundle $H$
  with
  \begin{equation*}
    E_+\oplus F_-\oplus H\cong F_+\oplus E_-\oplus H.
  \end{equation*}
  The class of $(E_+,E_-)$ is written $[E_+]-[E_-]$, and $[E]=[E]-[0]$. Addition is
  $([E_+]-[E_-])+([F_+]-[F_-])=[E_+\oplus F_+]-[E_-\oplus F_-]$.
\end{definition}

The relation is reflexive and symmetric; it is transitive because if
$E_+\oplus F_-\oplus H\cong F_+\oplus E_-\oplus H$ and
$F_+\oplus G_-\oplus H'\cong G_+\oplus F_-\oplus H'$ then
$E_+\oplus G_-\oplus(F_+\oplus F_-\oplus H\oplus H')\cong
G_+\oplus E_-\oplus(F_+\oplus F_-\oplus H\oplus H')$. Addition is well defined,
associative and commutative; $[E]-[E]$ is the zero element and $[E_-]-[E_+]$ is the
inverse of $[E_+]-[E_-]$. Thus $K^0(X)$ is an abelian group. By
\cref{bott:lem:bundles} every element can be written $[E]-[\C^N]$ (add a complement
of $E_-$ to both bundles), and $[E_+]-[E_-]=0$ if and only if
$E_+\oplus\C^L\cong E_-\oplus\C^L$ for some $L$ (replace $H$ by $H\oplus H^\perp\cong\C^L$).
A smooth map $f\colon X'\to X$ induces a homomorphism $f^*\colon K^0(X)\to K^0(X')$ by
pull-back of bundles. By \cref{bott:lem:bundle-homotopy}, $K^0([0,1]\times X)\to K^0(X)$,
restriction to $\{t\}\times X$, is independent of $t$. The difference of ranks at
a point identifies $K^0(\mathrm{pt})=\Z$; the \emph{rank} of
$[E_+]-[E_-]$ at $x$ is $\dim E_{+,x}-\dim E_{-,x}$. If $x_0\in X$ is a base point,
$\tilde K^0(X)=\ker\big(K^0(X)\to K^0(\{x_0\})\big)$ is the \emph{reduced} group, and
$K^0(X)=\Z\oplus\tilde K^0(X)$, the splitting being given by pull-back along
$X\to\{x_0\}$.

\begin{question}
  Why use formal differences $[E_+]-[E_-]$ rather than bundles themselves?
\end{question}
\begin{answer}
  Because the index is a difference. For one operator, $\dim\Nul A-\dim\Nul A^*$
  can be negative, so it lives in $\Z$, not in $\N_0$. For a family the null
  spaces and cokernels become bundles, and the index becomes ``null bundle minus
  cokernel bundle''. $K^0(X)$ is built from bundles exactly as $\Z$ is built from
  $\N_0$. The extra bundle $H$ in \cref{bott:def:K0} is needed because
  cancellation can fail for bundles: $E\oplus H\cong F\oplus H$ does not imply
  $E\cong F$.
\end{answer}

\subsection{Odd K-theory and pointed loops}

Odd K-theory is $K^{-1}(X)=[X;\Ginf(Y;E)]$, see \ref{bott:B:K1}. The loops in
$\Ginf$ give a second model of even K-theory.

\begin{definition}\label{bott:def:K2}
  For compact $X$, $K^{-2}(X)$ is the set of smooth homotopy classes of smooth maps
  $f\colon X\times\Sphere\to\Ginf(Y;E)$ with $f(x,1)=\Id$ for all $x$, where
  homotopies are required to satisfy the same condition. It is a group under
  pointwise multiplication.
\end{definition}

Equivalently, $K^{-2}(X)$ is the set of homotopy classes of smooth maps from $X$
into the pointed loop group
\begin{equation*}
  \Omega\Ginf(Y;E)=\{b\in\CI(\Sphere;\Ginf(Y;E));\ b(1)=\Id\}
\end{equation*}
(this is homotopy equivalent to the flat loop group $\Ginf_{(1)}(Y;E)$ of
\cref{ch:index}, by the flattening argument used there). The number
$\wind(f(x,\cdot))$ of \ref{bott:B:wind} is locally constant in $x$; we denote by
$\Omega_0\Ginf(Y;E)$ the loops with $\wind=0$.

\begin{lemma}\label{bott:lem:stable}
  Let $X'\subset X$ be closed (e.g.\ $X'=X\times\{1\}\subset X\times\Sphere$).
  \begin{enumerate}[label=(\roman*)]
  \item Every smooth $f\colon X\to\Ginf(Y;E)$ with $f=\Id$ on $X'$ is homotopic,
    through maps equal to $\Id$ on $X'$, to a map into a subgroup $\GL(V)\cong\GL(N,\C)$
    as in \ref{bott:B:finite}.
  \item If two such maps into $\GL(V)$ are homotopic in $\Ginf(Y;E)$ through maps
    equal to $\Id$ on $X'$, they are homotopic through such maps in $\GL(V')$ for some
    finite dimensional $V'\supset V$.
  \end{enumerate}
  In particular $K^{-2}(X)$ is the set of stable pointed homotopy classes of pointed
  loops in the groups $\GL(N,\C)$; in it $[a]+[b]=[ab]=[\diag(a,b)]$, and
  $K^{-2}(X)$ is abelian. The same holds for $K^{-1}(X)$, with free maps $X\to\GL(N,\C)$.
\end{lemma}
\begin{proof}
  (i) Write $f=\Id+A$. With $\Pi_j$ as in \ref{bott:B:finite}, the maps
  $f_t=\Id+(1-t)A+t\Pi_jA\Pi_j$ converge to $f$ uniformly on $[0,1]\times X$ as
  $j\to\infty$; since $f(X)$ is compact and $\Ginf$ is open, $f_t\in\Ginf$ for $j$
  large. Where $A=0$ also $f_t=\Id$, and $f_1=\Id-\Pi_j+(\Id+\Pi_jA\Pi_j)\Pi_j$ has
  values in $\GL(V_j)$.
  (ii) Apply (i) to the homotopy, a map on $[0,1]\times X$, with $V_0=V$ and closed
  set $[0,1]\times X'\cup\{0,1\}\times X$; on $\{0,1\}\times X$ the operator
  $A$ satisfies $A=\Pi_VA\Pi_V=\Pi_jA\Pi_j$, so the approximation does not change
  the endpoints. The final statements follow, using \cref{bott:lem:rotation}.
\end{proof}

\begin{remark}[The isotropic model]\label{bott:rem:iso-model}
  The proof of \cref{bott:lem:stable} uses only the finite rank approximations of
  \ref{bott:B:finite}. The residual isotropic group
  $\Gres(\R^n;\C^N)=\{\Id+A;\ A\in\Psiiso{-\infty}(\R^n;\C^N),\ \Id+A\text{ invertible}\}$
  of \cref{ch:isotropic} has the same property, with $\Pi_k$ the orthogonal projection
  onto the span of the Hermite functions $u_\alpha$, $|\alpha|\le k$ (tensored with
  $\C^N$): for every $A\in\Psiiso{-\infty}(\R^n;\C^N)$, $\Pi_kA\Pi_k\to A$ in
  $\Psiiso{-\infty}$, uniformly on compact families, because the Hermite expansion of a
  Schwartz kernel converges in $\cS$. Hence smooth homotopy classes of maps into
  $\Gres(\R^n;\C^N)$ are again the stable classes of maps into the groups $\GL(M,\C)$,
  and $K^{-1}(X)=[X;\Gres(\R^n;\C^N)]$, $K^{-2}(X)$ can equally be defined with
  $\Gres(\R^n;\C^N)$ in place of $\Ginf(Y;E)$. We use this model freely below.
\end{remark}

\begin{lemma}\label{bott:lem:circle-split}
  For compact $X$ let $\varpi\colon\Sphere\times X\to X$ be the projection and
  $\iota_1\colon X\to\Sphere\times X$, $x\mapsto(1,x)$. Then $\iota_1^*\varpi^*=\Id$ and
  forgetting the base point condition identifies $K^{-2}(X)$ with $\ker\iota_1^*$, so
  \begin{equation*}
    K^{-1}(\Sphere\times X)\cong K^{-2}(X)\oplus K^{-1}(X),\qquad
    [f]\longmapsto\bigl([f\cdot\varpi^*\iota_1^*f^{-1}],[\iota_1^*f]\bigr).
  \end{equation*}
\end{lemma}
\begin{proof}
  If $f\colon\Sphere\times X\to\Ginf$, then $f\cdot\varpi^*\iota_1^*f^{-1}$, i.e.\
  $(\theta,x)\mapsto f(\theta,x)f(1,x)^{-1}$, is pointed, and it is pointed homotopic
  to $f$ when $f$ is pointed. A free homotopy $h_t$ between pointed maps gives the
  pointed homotopy $h_t\cdot\varpi^*\iota_1^*h_t^{-1}$, so the forgetful map
  $K^{-2}(X)\to\ker\iota_1^*$ is injective, and by the first remark it is surjective.
  Since $K^{-1}(\Sphere\times X)$ is abelian, $[f]=[f\cdot\varpi^*\iota_1^*f^{-1}]+
  \varpi^*[\iota_1^*f]$, which gives the splitting.
\end{proof}

Via \cref{bott:lem:suspension} below, $K^{-2}(X)\cong\Kc^{-1}(\R\times X)$, so the lemma
says that $K^{-1}(\Sphere\times X)$ combines the odd K-theory of $X$ and the odd
K-theory of $\R\times X$; the latter will be identified with $K^0(X)$.

\section{Linearization of loops}\label{bott:sec:linearization}

A loop of invertible matrices can be complicated. We now show that, up to
stabilization and pointed homotopy, it can be replaced by the simplest possible
loop: on each of three complementary subspaces it is $z^{-1}$, $1$ or $z$. For
such a loop the index of the Toeplitz operator can be read off by counting
dimensions. The following normal form is due to Atiyah \cite{Atiyah-K}.

\begin{theorem}\label{bott:thm:linearization}
  Let $a\colon X\times\Sphere\to\GL(N,\C)$ be a pointed loop, $X$ compact. Then for
  some $M\ge N$ the stabilized loop $a\oplus\Id_{M-N}$ is pointed homotopic in
  $\GL(M,\C)$ to a loop of the form
  \begin{equation}\label{bott:eq:normal-form}
    \pi_-(x)z^{-1}+\pi_0(x)+\pi_+(x)z,
  \end{equation}
  where $\pi_-,\pi_0,\pi_+\colon X\to M(M,\C)$ are smooth families of idempotents
  which commute with each other and satisfy $\pi_-+\pi_0+\pi_+=\Id$.
\end{theorem}

Note that \eqref{bott:eq:normal-form} is invertible, with inverse
$\pi_-z+\pi_0+\pi_+z^{-1}$, and is automatically pointed. We call it a loop in
\emph{normal form}. The proof occupies the rest of this section and proceeds through
Laurent polynomials, polynomials and linear loops (\cref{bott:fig:linearization}).

\begin{figure}[ht]
\centering
\begin{tikzpicture}[
  box/.style={draw,rounded corners,align=center,font=\small,inner sep=4pt,
    text width=3.1cm},
  lab/.style={font=\footnotesize,align=center},
  >=Stealth]
\node[box] (a) at (0,0) {pointed loop\\ $a(x,z)$};
\node[box] (b) at (5.3,0) {Laurent polynomial\\ $z^{-q}\sum_{j=0}^{2q}b_jz^j$};
\node[box] (c) at (10.6,0) {polynomial loop\\ $b\oplus z^{-1}\Id_{qN}$};
\node[box] (d) at (10.6,-2.6) {linear loop\\ $\Id+(z-1)A$};
\node[box] (e) at (5.3,-2.6) {$(\Id-\pi)+z\pi$,\\ $\pi$ a Riesz projection};
\node[box] (f) at (0,-2.6) {normal form\\ $\pi_-z^{-1}+\pi_0+\pi_+z$};
\draw[->] (a)--node[above,lab]{Fourier}node[below,lab]{\cref{bott:lem:laurent}}(b);
\draw[->] (b)--node[above,lab]{rotation}node[below,lab]{\cref{bott:lem:rotation}}(c);
\draw[->] (c)--node[left,lab]{block matrices,\\ \cref{bott:lem:poly-linear}}(d);
\draw[->] (d)--node[above,lab]{spectrum}node[below,lab]{\cref{bott:lem:linear}}(e);
\draw[->] (e)--node[above,lab]{collect}node[below,lab]{blocks}(f);
\end{tikzpicture}
\caption{The steps in the linearization of loops
(\cref{bott:thm:linearization}). Each arrow is a stable pointed homotopy. The
degree in $z$ is first made finite, then reduced to one by enlarging the
matrices, and finally the linear loop is straightened by its spectral projection.}
\label{bott:fig:linearization}
\end{figure}

\begin{lemma}\label{bott:lem:laurent}
  A pointed loop $a$ is pointed homotopic to a pointed Laurent polynomial
  $\sum_{|j|\le q}a'_j(x)z^j$ with $a'_j\in\CI(X;M(N,\C))$.
\end{lemma}
\begin{proof}
  The Fourier coefficients $a_j(x)=\frac1{2\pi}\int_0^{2\pi}e^{-ij\theta}a(x,\theta)\,d\theta$
  satisfy $\sup_X|a_j|\le C_m(1+|j|)^{-m}$ for every $m$, by integration by parts
  in $\theta$, so $a_{(q)}=\sum_{|j|\le q}a_jz^j\to a$ uniformly on $X\times\Sphere$.
  Let $\delta=\min_{X\times\Sphere}\norm{a^{-1}}^{-1}>0$ and choose $q$ with
  $\sup\norm{a_{(q)}-a}<\delta$. Then
  $h_t=(1-t)a+ta_{(q)}=a\big(\Id+ta^{-1}(a_{(q)}-a)\big)$ is invertible for
  $t\in[0,1]$, and by \cref{bott:lem:normalize} $\hat h_t$ is a pointed homotopy
  from $a$ to $a_{(q)}(x,1)^{-1}a_{(q)}$, which is a pointed Laurent polynomial.
\end{proof}

So we may assume $a=z^{-q}b$ where $b(x,z)=\sum_{j=0}^{p}b_j(x)z^j$, $p=2q$, is a
polynomial in $z$ which is invertible for $|z|=1$ and satisfies $b(x,1)=\Id$. Since
$z^{-q}$ is a scalar, \cref{bott:lem:rotation} gives
\begin{gather*}
  \diag(a,\Id_N)=\diag(b\,z^{-q}\Id_N,\Id_N)\sim\diag(b,z^{-q}\Id_N),\\
  \diag(z^{-q}\Id_N,\Id_{(q-1)N})\sim z^{-1}\Id_{qN},
\end{gather*}
the second by writing $z^{-q}=z^{-1}\cdots z^{-1}$. It therefore suffices to deal
with the polynomial loop $b$.

\begin{lemma}\label{bott:lem:poly-linear}
  Let $b=\sum_{j=0}^pb_jz^j$ be a pointed polynomial loop in $\GL(N,\C)$. Then
  $\diag(b,\Id_{pN})$ is pointed homotopic in $\GL((p+1)N,\C)$ to a loop
  $\Id+(z-1)A(x)$, with $A\in\CI(X;M((p+1)N,\C))$.
\end{lemma}
\begin{proof}
  Set $c_j(x,z)=\sum_{l=j}^pb_l(x)z^{l-j}$ for $0\le j\le p$, so $c_0=b$, $c_p=b_p$
  and $c_j-zc_{j+1}=b_j$. In $(p+1)\times(p+1)$ block form with $N\times N$ blocks
  consider
  \begin{equation*}
    T(t)=\begin{pmatrix}b&tc_1&\cdots&tc_p\\0&\Id&&\\\vdots&&\ddots&\\0&&&\Id\end{pmatrix},\quad
    \Lambda(s)=\begin{pmatrix}\Id&&&\\-sz&\Id&&\\&\ddots&\ddots&\\&&-sz&\Id\end{pmatrix},
    \quad t,s\in[0,1].
  \end{equation*}
  $T(t)$ is block upper triangular with invertible diagonal blocks when $|z|=1$, and
  $\Lambda(s)$ is block lower triangular with identity diagonal, so both are invertible.
  The homotopy $T(t)$, $t\in[0,1]$, followed by $T(1)\Lambda(s)$, $s\in[0,1]$, starts
  at $\diag(b,\Id)$ and ends at
  \begin{equation*}
    T(1)\Lambda(1)=\begin{pmatrix}b_0&b_1&\cdots&b_{p-1}&b_p\\-z&\Id&&&\\&-z&\Id&&\\&&\ddots&\ddots&\\&&&-z&\Id\end{pmatrix}
    =B_0(x)+zB_1(x),
  \end{equation*}
  since the $(1,j)$ entry of the product is $c_j-zc_{j+1}=b_j$ for $j<p$ and $c_p=b_p$
  for $j=p$. Applying \cref{bott:lem:normalize} gives a pointed homotopy (the initial
  loop is pointed) ending at $(B_0+B_1)^{-1}(B_0+zB_1)=\Id+(z-1)A$ with
  $A=(B_0+B_1)^{-1}B_1$.
\end{proof}

\begin{lemma}\label{bott:lem:linear}
  Let $A\in\CI(X;M(M,\C))$. The loop $\ell(x,z)=\Id+(z-1)A(x)$ is invertible for all
  $|z|=1$ if and only if no eigenvalue of any $A(x)$ lies on the line
  $\{\Re\lambda=\frac12\}$. In that case $\ell$ is pointed homotopic, through loops
  of the same form, to $(\Id-\pi)+z\pi$ where $\pi(x)$ is the spectral projection of
  $A(x)$ onto the generalized eigenspaces with $\Re\lambda>\frac12$.
\end{lemma}
\begin{proof}
  For $z\ne1$, $\ell(z)=(z-1)\big(A-w\big)$ with $w=(1-z)^{-1}$, and as $z=e^{i\theta}$
  runs over $\Sphere\setminus\{1\}$,
  \begin{equation*}
    \frac1{1-e^{i\theta}}=\frac12+\frac i2\cot\frac\theta2
  \end{equation*}
  runs over the whole line $\Re w=\frac12$; at $z=1$, $\ell=\Id$. This proves the first
  statement (see \cref{bott:fig:linear}). Now the set $\{(x,\lambda);\ \det(\lambda-A(x))=0\}$ is compact and
  disjoint from $X\times\{\Re\lambda=\frac12\}$, so there is $\delta>0$ such that all
  eigenvalues satisfy $|\Re\lambda-\frac12|\ge\delta$. Let $\Gamma$ be the positively
  oriented boundary of the rectangle
  $\{\frac12+\frac\delta2\le\Re\lambda\le R,\ |\Im\lambda|\le R\}$ with
  $R>1+\sup_X\norm{A}$. Then
  \begin{equation*}
    \pi(x)=\frac1{2\pi i}\int_\Gamma(\lambda-A(x))^{-1}\,d\lambda
  \end{equation*}
  is a smooth family of idempotents commuting with $A(x)$, with range the sum of the
  generalized eigenspaces of $A(x)$ for eigenvalues with $\Re\lambda>\frac12$ and null
  space the sum of the others (the Riesz projection). Set $A_t=(1-t)A+t\pi$. On
  $\Ran\pi$, which is $A$-invariant, $A_t=(1-t)A+t$ has eigenvalues $(1-t)\lambda+t$
  with $\Re\lambda>\frac12$, so real part $>\frac12$; on $\Nul\pi$, $A_t=(1-t)A$ has
  eigenvalues $(1-t)\lambda$ with $\Re\lambda<\frac12$, so real part $<\frac12$. By the
  first part $\Id+(z-1)A_t$ is a pointed homotopy, and it ends at
  $\Id+(z-1)\pi=(\Id-\pi)+z\pi$.
\end{proof}

\begin{figure}[ht]
\centering
\begin{tikzpicture}[scale=1.1,>=Stealth]
\fill[gray!12] (1.5,-1.6) rectangle (3.4,1.6);
\draw[->,gray] (-1.6,0)--(3.6,0) node[right]{\small$\Re\lambda$};
\draw[->,gray] (0,-1.7)--(0,1.8) node[above]{\small$\Im\lambda$};
\draw[very thick] (1.5,-1.6)--(1.5,1.6) node[above,font=\footnotesize]{$\Re\lambda=\frac12$};
\fill (0,0) circle (1.3pt) node[above left,font=\footnotesize]{$0$};
\fill (3,0) circle (1.3pt) node[below,font=\footnotesize]{$1$};
\foreach \x/\y in {-1.1/0.9,0.6/-1.1,-0.5/-0.4}{
  \fill (\x,\y) circle (2pt);
  \draw[->,dashed,shorten >=2pt] (\x,\y)--(0,0);}
\foreach \x/\y in {2.2/1.2,3.1/-1.0,2.5/0.4}{
  \fill (\x,\y) circle (2pt);
  \draw[->,dashed,shorten >=2pt] (\x,\y)--(3,0);}
\node[font=\footnotesize,align=center] at (-0.9,1.45) {$\Nul\pi$};
\node[font=\footnotesize,align=center] at (2.75,1.4) {$\Ran\pi$};
\end{tikzpicture}
\caption{The proof of \cref{bott:lem:linear}. As $z$ runs over the circle, the
loop $\Id+(z-1)A$ fails to be invertible exactly when an eigenvalue of $A$ lies
on the line $\Re\lambda=\frac12$. The homotopy $A_t=(1-t)A+t\pi$ moves the
eigenvalues on the right of the line to $1$ and those on the left to $0$,
without crossing the line.}
\label{bott:fig:linear}
\end{figure}

\begin{proof}[Proof of \cref{bott:thm:linearization}]
  By \cref{bott:lem:laurent} and the reduction following it, $a\oplus\Id$ is stably
  pointed homotopic to $\diag(b,z^{-1}\Id_{qN})$, where $b$ is a pointed polynomial loop
  of degree at most $p=2q$. By \cref{bott:lem:poly-linear,bott:lem:linear},
  $\diag(b,\Id_{pN})\sim(\Id-\pi)+z\pi$. Hence, with $M=(3q+1)N$,
  $a\oplus\Id_{M-N}$ is pointed homotopic (after a constant permutation of blocks,
  \cref{bott:lem:conjugation}) to \eqref{bott:eq:normal-form} with
  $\pi_+=\pi\oplus0$, $\pi_0=(\Id-\pi)\oplus0$, $\pi_-=0\oplus\Id_{qN}$.
\end{proof}

\begin{remark}\label{bott:rem:linearization}
  The proof gives more: $\pi_-$ may be taken constant. We shall not use this. The
  spectral projection in \cref{bott:lem:linear} can also be written as a contour integral
  of the logarithmic derivative of the linear loop over $|z|=1$; the formulation above
  does not require the coefficient of $z$ to be invertible.
\end{remark}

The decisive point for K-theory is that a loop in normal form is trivial as soon as
$\Ran\pi_-$ and $\Ran\pi_+$ are isomorphic.

\begin{lemma}\label{bott:lem:normal-trivial}
  Let $\pi_-,\pi_0,\pi_+$ be commuting smooth idempotents on $\C^M$ over $X$ with
  sum $\Id$, and suppose there is a smooth bundle isomorphism
  $F\colon\Ran\pi_-\to\Ran\pi_+$. Then $\pi_-z^{-1}+\pi_0+\pi_+z$ is pointed homotopic
  to $\Id$ in $\GL(M,\C)$.
\end{lemma}
\begin{proof}
  The fibres split as $\C^M=V_-\oplus V_0\oplus V_+$, $V_\bullet=\Ran\pi_\bullet$.
  Regard $F\pi_-$ and $F^{-1}\pi_+$ as smooth maps into $M(M,\C)$ and set, for
  $\tau\in[0,\pi/2]$,
  \begin{equation*}
    h_\tau=\pi_0+\cos\tau\,(z^{-1}\pi_-+z\pi_+)+\sin\tau\,(F^{-1}\pi_+-F\pi_-).
  \end{equation*}
  $h_\tau$ is the identity on $V_0$ and preserves $V_-\oplus V_+$; identifying $V_+$ with
  $V_-$ by $F^{-1}$ it acts there as
  $\begin{pmatrix}\cos\tau\,z^{-1}&\sin\tau\\-\sin\tau&\cos\tau\,z\end{pmatrix}\otimes\Id_{V_-}$,
  which has determinant $1$. So $h_\tau$ is a loop in $\GL(M,\C)$, $h_0$ is the given
  loop, and $h_{\pi/2}=\pi_0+F^{-1}\pi_+-F\pi_-$ is independent of $z$. By
  \cref{bott:lem:normalize}, $\hat h_\tau$ is a pointed homotopy from $h_0$ to $\Id$.
\end{proof}

\begin{discussion}
  The scalar case explains the lemma. The loop $\diag(z^{-1},z)$ has
  $\det=1$, so winding number zero, and the rotation in the proof joins it to a
  constant matrix. In general the factor $z^{-1}$ on $\Ran\pi_-$ and the factor $z$
  on $\Ran\pi_+$ cancel as soon as the two bundles can be matched by an
  isomorphism $F$. So a normal form loop carries exactly one piece of
  information: the pair of bundles $(\Ran\pi_-,\Ran\pi_+)$, up to isomorphism.
  This is why the index of loops will take values in $K^0$.
\end{discussion}

\section{The index of Toeplitz operators on the circle}\label{bott:sec:circle}

We now quantize the normal form: replace $z^{-1}$ by the backward shift $L$ and
$z$ by the forward shift $U$. The resulting operator acts on each Fourier mode by
an explicit matrix, so its null space and range can be computed directly. Together
with the linearization this gives a second proof of
\cref{index:thm:toeplitz-index}, and one that extends to families.

For commuting smooth idempotents $\pi_-,\pi_0,\pi_+$ on $\C^M$ with sum $\Id$
consider the family of Toeplitz operators on $\Cp(\Sphere;\C^M)$
\begin{equation}\label{bott:eq:normal-op}
  Q(x)=\pi_-(x)L+\pi_0(x)+\pi_+(x)U,\qquad
  \sigma_0(Q)=\pi_-z^{-1}+\pi_0+\pi_+z.
\end{equation}
The matrices act pointwise and commute with $S$, $U$, $L$ and the $P_k$. For
$c\in\C^M$ write $c^\bullet=\pi_\bullet c$.

\begin{lemma}\label{bott:lem:normal-op}
  For every $k\ge0$ the null space of $Q(x)(\Id-P_k)$ is $\Ran P_k$, and
  \begin{equation*}
    W_k(x)=\operatorname{span}\{e^{ij\theta}\C^M\}_{0\le j\le k-1}
    \oplus e^{ik\theta}\Ran(\pi_0+\pi_+)(x)\oplus e^{i(k+1)\theta}\Ran\pi_+(x)
  \end{equation*}
  is a complement to the range of $Q(x)(\Id-P_k)$ in $\Cp(\Sphere;\C^M)$. In
  particular $\Nul Q(x)=\Ran\pi_-(x)$ (constant functions),
  $\Nul Q(x)^*=\Ran\pi_+(x)^*$ and $\ind Q(x)=\rank\pi_--\rank\pi_+$.
\end{lemma}
\begin{proof}
  For $u=\sum_{j\ge k+1}c_je^{ij\theta}$ (rapidly decreasing coefficients),
  \begin{equation*}
    Q(\Id-P_k)u=\sum_{j\ge k+1}\big(c_j^-e^{i(j-1)\theta}+c_j^0e^{ij\theta}+c_j^+e^{i(j+1)\theta}\big)
    =\sum_md_me^{im\theta}
  \end{equation*}
  with $d^-_m=c^-_{m+1}$, $d^0_m=c^0_m$, $d^+_m=c^+_{m-1}$. These relations determine $u$
  from the image, so $Q(\Id-P_k)$ is injective on $\Ran(\Id-P_k)$, and its range
  consists exactly of the $\sum d_me^{im\theta}\in\Cp$ with $d^-_m=0$ for $m<k$,
  $d^0_m=0$ for $m<k+1$ and $d^+_m=0$ for $m<k+2$. The complement $W_k$ follows. For
  a constant $c_0$, $Qc_0=c_0^0+c_0^+e^{i\theta}$, so $\Nul Q=\Ran\pi_-$. The
  index is $\dim\Ran P_0-\dim W_0=M-\rank\pi_0-2\rank\pi_+=\rank\pi_--\rank\pi_+$.
  Finally $Q^*=\pi_-^*U+\pi_0^*+\pi_+^*L$ is of the same form, with the commuting
  idempotents $\pi_\pm^*$ in exchanged roles, so $\Nul Q^*=\Ran\pi_+^*$.
\end{proof}

\begin{corollary}[Index formula for Toeplitz operators]\label{bott:cor:index-circle}
  If $P\in\PsiT{0}(\Sphere;\C^N)$ is elliptic, then
  \begin{equation}\label{bott:eq:index-circle}
    \ind P=\dim\Nul P-\dim\Nul P^*=\frac{i}{2\pi}\int_0^{2\pi}\Tr\Big(\sigma_0(P)^{-1}\frac{d\sigma_0(P)}{d\theta}\Big)d\theta
    =-\wind\sigma_0(P).
  \end{equation}
\end{corollary}
\begin{proof}
  The last two expressions agree by \ref{bott:B:wind}. Let $c=\sigma_0(P)(1)\in\GL(N,\C)$.
  Then $c^{-1}P$ is elliptic with pointed symbol $c^{-1}\sigma_0(P)$; it has the same
  index as $P$ since $c$ is an invertible operator, and the same winding number since
  $\det c$ is constant. So we may assume $\sigma_0(P)$ pointed. Both sides of
  \eqref{bott:eq:index-circle} depend only on the stable pointed homotopy class of the
  symbol: the index by \ref{bott:B:toeplitz}, and because stabilizing $P$ by $\Id$ adds
  nothing to either side; the winding number by \ref{bott:B:wind}. By
  \cref{bott:thm:linearization} (with $X$ a point) we may thus assume
  $\sigma_0(P)=\pi_-z^{-1}+\pi_0+\pi_+z$, and then $P$ may be replaced by $Q$ of
  \eqref{bott:eq:normal-op}. Then $\ind Q=\rank\pi_--\rank\pi_+$ by
  \cref{bott:lem:normal-op}, while $\det\sigma_0(Q)=z^{\rank\pi_+-\rank\pi_-}$, since $\sigma_0(Q)$
  acts as $z^{-1}$, $1$, $z$ on the summands of $\C^M=\Ran\pi_-\oplus\Ran\pi_0\oplus\Ran\pi_+$, so $\wind\sigma_0(Q)=\rank\pi_+-\rank\pi_-$.
\end{proof}

The same holds for pseudodifferential operators on the circle, except that now both
halves $\Sphere_\pm$ of $S^*\Sphere$ contribute. For $P\in\Psido{0}(\Sphere;\C^N)$
write $\sigma_\pm(\theta)=\sigma_0(P)(\theta,\pm1)$.

\begin{proposition}\label{bott:prop:index-psido-circle}
  If $P\in\Psido{0}(\Sphere;\C^N)$ is elliptic, then
  \begin{equation*}
    \ind P=\frac{i}{2\pi}\int_0^{2\pi}\Tr\big(\sigma_+^{-1}\partial_\theta\sigma_+\big)d\theta
    -\frac{i}{2\pi}\int_0^{2\pi}\Tr\big(\sigma_-^{-1}\partial_\theta\sigma_-\big)d\theta
    =\wind\sigma_--\wind\sigma_+ .
  \end{equation*}
  Equivalently $\ind P=\frac{i}{2\pi}\int_{S^*\Sphere}\Tr(\sigma_0(P)^{-1}d\sigma_0(P))$
  with $\Sphere_+$ oriented by $d\theta$ and $\Sphere_-$ by $-d\theta$.
\end{proposition}
\begin{proof}
  Let $S'=\Id-S$, the projection onto $\CI_-(\Sphere)$, the span of the negative
  Fourier modes. By \ref{bott:B:toeplitz}, $SPS'$ and $S'PS$ have vanishing principal
  symbol, so $P_t=P-t(SPS'+S'PS)$ is a family of elliptic operators with constant
  principal symbol, and $\ind P=\ind P_1$. Now $P_1=SPS+S'PS'$ preserves the
  decomposition $\CI(\Sphere;\C^N)=\Cp\oplus\CI_-$, so
  $\ind P=\ind(SPS|_{\Cp})+\ind(S'PS'|_{\CI_-})$. The first is an elliptic Toeplitz
  operator with symbol $\sigma_+$, of index $-\wind\sigma_+$ by
  \cref{bott:cor:index-circle}. For the second, let $Ru(\theta)=e^{-i\theta}u(-\theta)$.
  Then $R(e^{ik\theta})=e^{-i(k+1)\theta}$, so $R$ is an isomorphism of $\CI(\Sphere;\C^N)$
  exchanging $\CI_-$ and $\Cp$, and $RS'R^{-1}=S$. Thus $R(S'PS')R^{-1}=STS$ with
  $T=RPR^{-1}\in\Psido{0}(\Sphere;\C^N)$; as $R$ is the composite of the pull-back by
  $\theta\mapsto-\theta$ and a multiplication, $\sigma_0(T)(\theta,\xi)=\sigma_0(P)(-\theta,-\xi)$.
  So $S'PS'|_{\CI_-}$ is conjugate to the elliptic Toeplitz operator with symbol
  $\sigma_-(-\theta)$, whose index is $-\wind(\sigma_-(-\cdot))=\wind\sigma_-$.
\end{proof}

\begin{remark}
  The index of an elliptic \emph{differential} operator $P\in\Diff{m}(\Sphere;\C^N)$
  vanishes. Indeed $\sigma_m(P)=p(\theta)\xi^m$ with $p$ invertible, and
  $P(\Id+D_\theta^2)^{-m/2}\in\Psido{0}$ has the same index as $P$ (the second factor
  is invertible) and symbol $\sigma_\pm=(\pm1)^mp$. Since $\wind((-1)^mp)=\wind p$,
  the two terms in \cref{bott:prop:index-psido-circle} cancel.
\end{remark}

\section{The families index}\label{bott:sec:families}

We first record that inverses of smooth families are smooth. The same argument applies
to the operators with values in smoothing operators of \cref{bott:sec:classifying}.

\begin{lemma}\label{bott:lem:inverse}
  Let $Z$ be a manifold and $A\in\CI(Z;\PsiT{0}(\Sphere;\C^N))$ a smooth family of
  elliptic operators, each invertible on $\Cp(\Sphere;\C^N)$ (equivalently on
  $L^2_+(\Sphere;\C^N)$). Then $A^{-1}\in\CI(Z;\PsiT{0}(\Sphere;\C^N))$.
\end{lemma}
\begin{proof}
  The symbolic parametrix construction (\cref{ch:manifolds}) works with smooth
  dependence on parameters: it uses only the multiplicativity of the symbol, the
  inversion of the symbol and asymptotic summation. This gives a smooth family
  $B$ with $BA=\Id-R$, $AB=\Id-R'$, where $R,R'\in\CI(Z;\PsiT{-\infty}(\Sphere;\C^N))$.
  The inclusion of $\PsiT{0}$ into the bounded operators on $L^2_+$ is continuous, so
  $z\mapsto A(z)$ is smooth in operator norm, and hence so is $V=A^{-1}$. From
  $V=B+RV$ and $V=B+VR'$ we get $V=B+RB+RVR'$. The bilinear map $T\mapsto RTR'$ is
  continuous from bounded operators to smoothing operators (its kernel is
  $(w,w')\mapsto\inner{TR'(\cdot,w')}{\overline{R(w,\cdot)}}$ and all derivatives fall on
  $R$ and $R'$), so $RVR'$, and hence $V$, is a smooth family in $\PsiT{0}$.
\end{proof}

Now let $Z$ be a compact manifold (possibly with boundary) and
$A\in\CI(Z;\PsiT{0}(\Sphere;\C^N))$ a smooth family of elliptic Toeplitz operators.
Since $A(z)$ is Fredholm, the range of $A(z)(\Id-P_k)=A(z)-A(z)P_k$ is closed and of
finite codimension, equal to $(k+1)N-\ind A(z)$ if its null space is $\Ran P_k$; the
same holds on $\Cp$ and on $L^2_+$.

\begin{definition}\label{bott:def:index-data}
  An integer $k\ge0$ is \emph{admissible} for $A$ if $\Nul(A(z)(\Id-P_k))=\Ran P_k$ for
  all $z\in Z$. A \emph{complement} for $(A,k)$ is a smooth vector bundle $W$ over $Z$
  whose fibres $W(z)\subset\Cp(\Sphere;\C^N)$ satisfy
  $\Cp(\Sphere;\C^N)=A(z)(\Id-P_k)\Cp\oplus W(z)$; smoothness means that $W$ has local
  frames given by smooth maps into $\Cp(\Sphere;\C^N)$.
\end{definition}

If $k$ is admissible then so is every $k'\ge k$, since $\Ran(\Id-P_{k'})\subset\Ran(\Id-P_k)$.

\begin{lemma}\label{bott:lem:index-bundle}
  Let $A$ be a smooth elliptic family as above.
  \begin{enumerate}[label=(\roman*)]
  \item There is $k_0$ such that every $k\ge k_0$ is admissible.
  \item If $k$ is admissible, the orthocomplements $F_k(z)$ of the ranges of
    $A'(z)=A(z)(\Id-P_k)$ in $L^2_+$ form a complement for $(A,k)$, and the orthogonal
    projections onto $F_k(z)$ form a smooth family. Every complement $W$ for $(A,k)$ is
    isomorphic to $F_k$.
  \item If $k$ is admissible with complement $W$, then
    $W\oplus A(e^{i(k+1)\theta}\C^N)\cong W\oplus\C^N$ is a complement for $(A,k+1)$.
  \end{enumerate}
  Consequently the class
  \begin{equation*}
    \ind(A)=[\Ran P_k]-[W]=[\C^{(k+1)N}]-[W]\in K^0(Z)
  \end{equation*}
  is independent of the admissible $k$ and of the complement $W$ for $(A,k)$.
\end{lemma}
\begin{proof}
  (i) Choose a smooth family of parametrices $B$, so $BA=\Id-R$ with
  $R\in\CI(Z;\PsiT{-\infty})$. Then $BA(\Id-P_k)=(\Id-R(\Id-P_k))(\Id-P_k)$. By
  \ref{bott:B:shifts}, $R(\Id-P_k)\to0$ in $\PsiT{-\infty}$, hence in $L^2$ operator norm,
  uniformly on $Z$, so for $k\ge k_0$ the operators $\Id-R(\Id-P_k)$ are invertible and
  $(\Id-P_k)(\Id-R(\Id-P_k))^{-1}BA(\Id-P_k)=\Id-P_k$. Hence $\Nul A(\Id-P_k)=\Ran P_k$.

  (ii) The family $C=A'^*A'+P_k$ is elliptic (its symbol is $\sigma_0(A)^*\sigma_0(A)$) and
  injective: $\inner{Cu}{u}=\norm{A'u}^2+\norm{P_ku}^2=0$ forces $P_ku=0$ and $A'u=0$, so
  $u\in\Ran(\Id-P_k)\cap\Ran P_k=\{0\}$. Being self-adjoint and Fredholm it is invertible,
  and $C^{-1}$ is a smooth family by \cref{bott:lem:inverse}. Since $CP_k=P_k$ and
  $C(\Id-P_k)=A'^*A'$, the family $G=C^{-1}A'^*$ satisfies $GA'=\Id-P_k$. Then $A'G$ is the
  orthogonal projection onto $\Ran A'$ (it is idempotent, self-adjoint as
  $A'G=A'C^{-1}A'^*$, and has range $\Ran A'$), so $\Pi=\Id-A'G$ is a smooth family of
  orthogonal projections onto $F_k$. Its range consists of smooth functions (it is the null
  space of the elliptic operator $A'^*$) and its rank $(k+1)N-\ind A(z)$ is locally
  constant. If $v_1,\dots,v_r$ is a basis of $F_k(z_0)$, then $\Pi(z)v_1,\dots,\Pi(z)v_r$
  are smooth maps into $\Cp$, linearly independent near $z_0$, so form a local frame:
  $F_k$ is a smooth bundle. As $\Ran A'$ is closed and $A'(\Cp)=\Ran A'\cap\Cp$ (if
  $A'v\in\Cp$ then $GA'v=(\Id-P_k)v\in\Cp$), $F_k$ is a complement. Finally, if $W$ is a
  complement then $\Pi|_W\colon W\to F_k$ is a smooth bundle map; it is injective because
  $\Nul\Pi\cap\Cp=A'(\Cp)$ meets $W$ only in $0$, and $\dim W=\dim F_k$, both being the
  codimension of $A'(\Cp)$. So $W\cong F_k$.

  (iii) Since $A$ is injective on $\Ran(\Id-P_k)$, we have
  $A(\Id-P_k)\Cp=A(\Id-P_{k+1})\Cp\oplus A(e^{i(k+1)\theta}\C^N)$. Hence
  $W\oplus A(e^{i(k+1)\theta}\C^N)$ is a complement for $(A,k+1)$, and
  $v\mapsto A(e^{i(k+1)\theta}v)$ trivializes the second summand.

  For the consequence: given admissible $k\le k'$ with complements $W,W'$, (iii)
  applied $k'-k$ times gives a complement $W\oplus\C^{(k'-k)N}$ for $k'$, isomorphic to
  $W'$ by (ii); since $\Ran P_{k'}=\Ran P_k\oplus\C^{(k'-k)N}$ the classes agree.
\end{proof}

\begin{definition}\label{bott:def:families-index}
  The class $\ind(A)\in K^0(Z)$ of \cref{bott:lem:index-bundle} is the \emph{families
  index} of the elliptic family $A$.
\end{definition}

\begin{question}
  Why is the families index a class in $K^0(Z)$ and not just an integer?
\end{question}
\begin{answer}
  The integer $\ind A(z)$ is the rank of $\ind(A)$ at $z$, and it is constant on
  each component of $Z$. But the null spaces and the cokernels can twist as $z$
  moves. For the normal form operator $Q$ of \eqref{bott:eq:normal-op}, the null
  space is the bundle $\Ran\pi_-$ and the cokernel is $\Ran\pi_+$
  (\cref{bott:lem:normal-op}). If $\Ran\pi_-$ is a non-trivial line bundle and
  $\Ran\pi_+=0$, the integer index is $1$ at every point, yet
  $\ind(Q)=[\Ran\pi_-]$ differs from $[\C]$ whenever this line bundle is not
  stably trivial, as for the tautological line bundle over $\Sphere^2$. The class
  in $K^0(Z)$ records this twisting.
\end{answer}

For $Z$ a point this is the numerical index, since $(k+1)N-\dim W=\ind A$. In the
notation of \cref{ch:index}, $F_k=\Nul((\Id-P_k)A^*)$, and \cref{bott:lem:index-bundle}
shows that the bundle considered there represents, together with the trivial bundle
$\Ran P_k$, a well-defined element of $K^0(Z)$. The families index has the expected
properties.

\begin{proposition}\label{bott:prop:index-properties}
  Let $A,A'$ be smooth elliptic families of Toeplitz operators over $Z$.
  \begin{enumerate}[label=(\roman*)]
  \item (Homotopy) If $A_t$ is an elliptic family over $[0,1]\times Z$, then
    $\ind(A_0)=\ind(A_1)$. In particular $\ind(A)$ depends only on the homotopy class
    of $\sigma_0(A)$ in $\CI(Z\times\Sphere;\GL(N,\C))$.
  \item (Additivity) $\ind(A\oplus A')=\ind(A)+\ind(A')$; $\ind(A\oplus\Id)=\ind(A)$.
  \item (Naturality) $\ind(f^*A)=f^*\ind(A)$ for smooth $f\colon Z'\to Z$.
  \item If each $A(z)$ is invertible then $\ind(A)=0$. If $b,c\colon Z\to\GL(N,\C)$
    then $\ind(bAc)=\ind(A)$.
  \item (Normal form) For $Q$ as in \eqref{bott:eq:normal-op},
    $\ind(Q)=[\Ran\pi_-]-[\Ran\pi_+]$.
  \end{enumerate}
\end{proposition}
\begin{proof}
  (i) Choose $k$ admissible for the family over $[0,1]\times Z$; the restrictions of its
  complement $F_k$ to $t=0,1$ are complements for $A_0$, $A_1$ and are isomorphic by
  \cref{bott:lem:bundle-homotopy}. Two families with homotopic symbols are joined by an
  elliptic family: first the linear homotopy between $A$ and $T_{\sigma_0(A)}$, which has
  constant symbol, then $T_{a_t}$ for a homotopy $a_t$ of symbols. (ii) and (iii) are clear
  from the definition, using $P_k$ diagonally; a complement for $(\Id,k)$ is $\Ran P_k$.
  (iv) If $A$ is invertible every $k$ is admissible and $A(\Ran P_k)\cong\C^{(k+1)N}$ is a
  complement. Multiplication by $b(z)$ and $c(z)$ commutes with $P_k$, so
  $(bAc)(\Id-P_k)$ has null space $c^{-1}\Ran P_k=\Ran P_k$, range $bA(\Id-P_k)\Cp$ and
  complement $bW\cong W$. (v) By \cref{bott:lem:normal-op} every $k$ is admissible, with
  complement $W_k\cong\C^{kM}\oplus\Ran(\pi_0+\pi_+)\oplus\Ran\pi_+$ (a smooth bundle by
  \cref{bott:lem:bundles}(ii)); since $\C^M=\Ran\pi_-\oplus\Ran(\pi_0+\pi_+)$,
  $[\C^{(k+1)M}]-[W_k]=[\Ran\pi_-]-[\Ran\pi_+]$.
\end{proof}

The families index is the only obstruction to invertibility after perturbation; this
is the criterion of \cref{ch:index} (there phrased as triviality of the bundle $F_k$),
up to stabilization.

\begin{proposition}\label{bott:prop:perturb}
  Let $A$ be a smooth elliptic family of Toeplitz operators on $\C^N$ over compact $Z$.
  Then $\ind(A)=0$ in $K^0(Z)$ if and only if for some $L$ there is a smooth family
  $B\colon Z\to\PsiT{-\infty}(\Sphere;\C^{N+L})$ such that $A\oplus\Id_L+B$ is
  invertible for all $z\in Z$.
\end{proposition}
\begin{proof}
  If $A\oplus\Id_L+B$ is invertible, then $\ind(A)=\ind(A\oplus\Id_L+B)=0$ by
  \cref{bott:prop:index-properties}. Conversely, let $k$ be admissible with complement
  $W$. If $\ind(A)=0$ then $W\oplus\C^J\cong\C^{(k+1)N+J}$ for some $J$. Choose $L$ with
  $(k+1)L\ge J$; then $\tilde A=A\oplus\Id_L$ has complement
  $\tilde W=W\oplus\C^{(k+1)L}\cong\C^{(k+1)(N+L)}=\Ran P_k$, so there is a smooth bundle
  isomorphism $\Theta\colon\Ran P_k\to\tilde W$. Put $B=\Theta P_k-\tilde AP_k$; this is a
  smooth family of finite rank operators with smooth kernels, so $B\in\PsiT{-\infty}$,
  and $\tilde A+B=\tilde A(\Id-P_k)+\Theta P_k$ maps $\Ran(\Id-P_k)$ injectively onto
  $\tilde A(\Id-P_k)\Cp$ and $\Ran P_k$ onto the complement $\tilde W$. So it is a
  bijection of $\Cp$, hence invertible, and the inverses form a smooth family by
  \cref{bott:lem:inverse}.
\end{proof}

\section[The index isomorphism]{The index isomorphism and Bott periodicity, first form}\label{bott:sec:index-iso}

We can now assemble the pieces. Quantizing a loop gives a family of Toeplitz
operators, and its families index is a class in $K^0(X)$. The main theorem says
that this is a bijection between loops (up to stable pointed homotopy) and
$K^0(X)$. Injectivity comes from the linearization and
\cref{bott:lem:normal-trivial}; surjectivity from an explicit inverse, the Bott
map, which turns a pair of bundles into a normal form loop.

For a pointed loop $a\colon X\times\Sphere\to\GL(N,\C)$, $T_a(x)=Sa(x,\cdot)S$ is a
smooth elliptic family of Toeplitz operators with symbol $a$; set
\begin{equation*}
  \ind(a)=\ind(T_a)\in K^0(X).
\end{equation*}
By \cref{bott:prop:index-properties}(i) any elliptic family with symbol $a$ has the same
index. By \cref{bott:prop:index-properties}(i),(ii), $\ind(a)$ depends only on the stable
pointed homotopy class of $a$, so by \cref{bott:lem:stable} it defines a map
\begin{equation*}
  \ind\colon K^{-2}(X)\longrightarrow K^0(X),
\end{equation*}
which is a homomorphism since $[a]+[b]=[\diag(a,b)]$ and $T_{\diag(a,b)}=T_a\oplus T_b$.
For $X$ a point, $\ind(a)=-\wind(a)$ by \cref{bott:cor:index-circle}; in general the
rank of $\ind(a)$ at $x$ is $-\wind a(x,\cdot)$.

\begin{proposition}\label{bott:prop:injective}
  If $a$ is a pointed loop over $X$ with $\ind(a)=0$ in $K^0(X)$, then $a$ is stably
  pointed homotopic to $\Id$.
\end{proposition}
\begin{proof}
  By \cref{bott:thm:linearization} and homotopy invariance we may assume
  $a=\pi_-z^{-1}+\pi_0+\pi_+z$, and then $[\Ran\pi_-]-[\Ran\pi_+]=\ind(a)=0$ by
  \cref{bott:prop:index-properties}(v), i.e.\ $\Ran\pi_-\oplus H\cong\Ran\pi_+\oplus H$ for
  some bundle $H$. By \cref{bott:lem:bundles} we may take $H=\Ran q$ for a smooth family of
  idempotents $q$ on some $\C^K$. The loop
  \begin{equation*}
    \beta_q=\diag\big(qz^{-1}+(\Id-q),\,(\Id-q)+qz\big)
  \end{equation*}
  on $\C^{2K}$ is in normal form with $\Ran\pi_-^{\beta}=\Ran q\oplus0\cong0\oplus\Ran q=\Ran\pi_+^\beta$,
  so $\beta_q$ is pointed homotopic to $\Id$ by \cref{bott:lem:normal-trivial}. Hence
  $a\sim\diag(a,\beta_q)$, and $\diag(a,\beta_q)$ is in normal form with idempotents
  $\pi_-\oplus q\oplus0$, $\pi_0\oplus(\Id-q)\oplus(\Id-q)$, $\pi_+\oplus0\oplus q$, whose
  ranges satisfy $\Ran\pi_-\oplus H\cong\Ran\pi_+\oplus H$. By
  \cref{bott:lem:normal-trivial} again, $\diag(a,\beta_q)$ is pointed homotopic to $\Id$.
\end{proof}

\begin{theorem}[Bott periodicity, first form]\label{bott:thm:index-iso}
  For every compact manifold $X$ the index map $\ind\colon K^{-2}(X)\to K^0(X)$ is an
  isomorphism of abelian groups, natural with respect to pull-back. Its inverse, the
  \emph{Bott map} $\beta$, is given as follows: if $E_\pm\cong\Ran p_\pm$ with
  $p_\pm\colon X\to M(N_\pm,\C)$ smooth families of projections, then
  \begin{equation}\label{bott:eq:bott-map}
    \beta([E_+]-[E_-])=\big[\big(p_+z^{-1}+(\Id-p_+)\big)\oplus\big((\Id-p_-)+p_-z\big)\big].
  \end{equation}
\end{theorem}
\begin{proof}
  Injectivity is \cref{bott:prop:injective}, and naturality follows from
  \cref{bott:prop:index-properties}(iii). The loop in \eqref{bott:eq:bott-map} is in normal
  form with $\pi_-=p_+\oplus0$ and $\pi_+=0\oplus p_-$, so by
  \cref{bott:prop:index-properties}(v) its index is $[\Ran p_+]-[\Ran p_-]=[E_+]-[E_-]$. As
  every element of $K^0(X)$ is of the form $[E_+]-[E_-]$ and every bundle is the range
  of a family of projections (\cref{bott:lem:bundles}), $\ind$ is surjective, and
  $\beta$ is its inverse.
\end{proof}

Concretely, for the operator with symbol \eqref{bott:eq:bott-map}, namely
$(p_+L+\Id-p_+)\oplus(\Id-p_-+p_-U)$, the null space is the bundle $E_+$ of constant
sections and the null space of the adjoint is $E_-$. The Toeplitz operator $L$, of index
$1=[\C]\in K^0(\mathrm{pt})$, is the \emph{Bott element}: $\beta[\C]=[z^{-1}]$.
\Cref{bott:fig:bott-map} summarizes the construction.

\begin{figure}[ht]
\centering
\begin{tikzcd}[column sep=3.2em,row sep=2.6em,ampersand replacement=\&,
  every label/.append style={font=\footnotesize}]
  {[E_+]-[E_-]}\in K^0(X)
    \arrow[r,"\beta"]
  \& {\big[(p_+z^{-1}+\Id-p_+)\oplus(\Id-p_-+p_-z)\big]}\in K^{-2}(X)
    \arrow[d,"{\text{quantize: }z^{-1}\mapsto L,\ z\mapsto U}"] \\
  \& (p_+L+\Id-p_+)\oplus(\Id-p_-+p_-U)
    \arrow[ul,bend left=12,"{\ind=[\Nul]-[\text{coker}]}"]
\end{tikzcd}
\caption{The Bott map $\beta$ and the index. A pair of bundles $E_\pm=\Ran p_\pm$
is turned into a loop in normal form, with $z^{-1}$ on $E_+$ and $z$ on $E_-$.
The corresponding Toeplitz family has null bundle $E_+$ (the constant sections
killed by $L$) and cokernel $E_-$ (the constants missed by $U$), so its families
index is $[E_+]-[E_-]$.}
\label{bott:fig:bott-map}
\end{figure}

\begin{discussion}
  Why ``periodicity''? Each step from $K^0$ to $K^{-1}$ to $K^{-2}$ adds one
  real variable: by clutching, $K^{-1}(X)\cong\Kc^0(\R\times X)$
  (\cref{bott:lem:clutching}), and $K^{-2}(X)\cong\Kc^{-1}(\R\times X)$
  (\cref{bott:lem:suspension}). \Cref{bott:thm:index-iso} says that after two
  such steps we are back at $K^0(X)$. So the K-groups repeat with period two;
  in terms of homotopy groups, $\pi_{k+2}(\Ginf)\cong\pi_k(\Ginf)$, which is
  proved in \cref{bott:cor:homotopy-groups}.
\end{discussion}

\begin{remark}
  The signs matter here: the loop $\pi_-z^{-1}+\pi_0+\pi_+z$ with $\Ran\pi_\pm\cong E_\pm$
  represents $-\beta([E_+]-[E_-])$, because $\ind U=-1$.
\end{remark}

\begin{corollary}\label{bott:cor:reduced}
  If $x_0\in X$ is a base point and $\tilde K^{-2}(X)$ denotes the kernel of restriction
  $K^{-2}(X)\to K^{-2}(\{x_0\})$, then $\ind$ restricts to an isomorphism
  $\tilde K^{-2}(X)\cong\tilde K^0(X)$. Moreover $K^{-2}(\mathrm{pt})\cong\Z$, i.e.\
  pointed loops in $\Ginf(Y;E)$ (equivalently in $\GL(N,\C)$, stably) are classified up to
  pointed homotopy by their winding number, and every winding number occurs.
\end{corollary}
\begin{proof}
  Restriction and pull-back along $X\to\{x_0\}$ commute with $\ind$, and $\ind$ is an
  isomorphism for $X$ and for $\{x_0\}$, so it maps the kernel of one restriction
  bijectively onto the kernel of the other. The second statement is the case $X=\mathrm{pt}$,
  where $\ind(a)=-\wind(a)$.
\end{proof}

\section{The isotropic index map}\label{bott:sec:isotropic}

Toeplitz operators are one way to `quantize' loops of matrices. A second way uses the
isotropic calculus on $\R$ of \cref{ch:isotropic}: the principal symbol of an isotropic
operator of order $0$ on $\R$ is a function on the circle at infinity of the phase
plane $\R^2$, so loops are again symbols. The role of the projections $P_k$ onto the
first Fourier modes is taken by the eigenprojections of the harmonic oscillator, and
the role of the shift $L$ by the annihilation operator. In this section we carry this
out and obtain a second proof of \cref{bott:thm:index-iso}; the two index maps are
compared in \cref{bott:prop:iso-toeplitz}. Unlike the Toeplitz construction, the
isotropic one extends directly to $\R^n$, which is exploited in
\cref{bott:sec:quantization}.

\subsection{Conventions}
The phase plane is $\R^2$ with coordinates $w=(x,\xi)$, oriented by $dx\wedge d\xi$.
We identify $\Sphere$ with the unit circle by $\theta\mapsto(\cos\theta,\sin\theta)$, so
that $\Sphere$ is traversed counterclockwise, and on it
\begin{equation*}
  z=e^{i\theta}=\frac{x+i\xi}{\abs w};
\end{equation*}
the base point $\theta=0$ is $(x,\xi)=(1,0)$. For $A\in\Psiisoph{0}(\R;\C^N)$, a
polyhomogeneous isotropic operator of order $0$, the principal symbol is
$\sigma_0(A)\in\CI(\Sphere;M(N,\C))$, and $A$ is elliptic if $\sigma_0(A)$ is invertible.
Conversely, for $a\in\CI(Z\times\Sphere;M(N,\C))$ the family
\begin{equation}\label{bott:eq:iso-quantize}
  \Op_\chi(a)(y)=\Op\bigl(\chi(\abs w)\,a(y,w/\abs w)\bigr),\qquad y\in Z,
\end{equation}
where $\chi\in\CI(\R)$, $\chi=0$ near $0$ and $\chi=1$ on $[1,\infty)$,
is a smooth family in $\Psiisoph{0}(\R;\C^N)$ with principal symbol $a(y,\cdot)$, the
symbol being homogeneous of degree $0$ outside the unit ball. We write
$u_k$, $k\in\N_0$, for the Hermite functions on $\R$ (the normalized eigenfunctions of
$H=D_x^2+x^2$, $Hu_k=(2k+1)u_k$) and $\Pi_k$ for the orthogonal projection onto
$\operatorname{span}\{u_0,\dots,u_k\}\otimes\C^N$.

\subsection{The dictionary}
The constructions of \cref{bott:sec:families} used the following properties of the
Toeplitz algebra: the existence of parametrices depending smoothly on parameters,
the continuity of the inclusion into the bounded operators on $L^2_+$, elliptic
regularity, the fact that $R\,T\,R'$ is a smoothing operator, continuously in $T$, when
$R,R'$ are smoothing and $T$ is bounded, and \ref{bott:B:shifts}. Each of these holds
for the isotropic calculus by the results of \cref{ch:isotropic}:
\begin{enumerate}[label=(I\arabic*),leftmargin=*]
\item\label{bott:I:parametrix} An elliptic $A\in\Psiisoph{0}(\R;\C^N)$ has a parametrix
  $B\in\Psiisoph{0}(\R;\C^N)$, $AB-\Id,\ BA-\Id\in\Psiiso{-\infty}(\R;\C^N)$; the
  construction is symbolic (inversion of the symbol and asymptotic summation), so for
  a smooth family $A$ the parametrices may be chosen to form a smooth family.
\item\label{bott:I:fredholm} An elliptic $A$ is Fredholm on $L^2(\R;\C^N)$ and on
  $\cS(\R;\C^N)$; the null spaces of $A$ and $A^*$ lie in $\cS$, they are the same on
  $\cS$, $L^2$ and $\cSp$, and the range on each of these spaces is the annihilator of
  $\Nul A^*$. The inclusion $\Psiisoph{0}(\R;\C^N)\to\Lin(L^2)$ is continuous.
\item\label{bott:I:corner} If $R,R'\in\Psiiso{-\infty}(\R;\C^N)$ and $T$ is bounded on
  $L^2$ then $RTR'\in\Psiiso{-\infty}(\R;\C^N)$, with seminorms bounded by seminorms of
  $R$, $R'$ times $\norm T$.
\item\label{bott:I:hermite} For a smooth family $K$ in $\Psiiso{-\infty}(\R;\C^N)$
  parametrized by a compact manifold, $K\Pi_k\to K$ and $\Pi_kK\to K$ in
  $\Psiiso{-\infty}$, uniformly, as $k\to\infty$ (the Hermite expansion of a Schwartz
  kernel converges in $\cS(\R^2)$). Moreover $\Ran\Pi_{k+1}=\Ran\Pi_k\oplus u_{k+1}\C^N$.
\end{enumerate}
The dictionary is therefore
\begin{gather*}
  \Cp(\Sphere;\C^N)\leadsto\cS(\R;\C^N),\quad L^2_+(\Sphere;\C^N)\leadsto L^2(\R;\C^N),\quad
  \PsiT{0}\leadsto\Psiisoph{0},\\
  \PsiT{-\infty}\leadsto\Psiiso{-\infty},\quad
  P_k\leadsto\Pi_k,\quad e^{ij\theta}\leadsto u_j,
\end{gather*}
with $T_a$ replaced by $\Op_\chi(a)$. Making these replacements in
\cref{bott:lem:inverse,bott:lem:index-bundle,bott:prop:perturb} and in
\cref{bott:prop:index-properties}(i)--(iv) we obtain, with the same proofs:

\begin{proposition}\label{bott:prop:iso-families}
  Let $Z$ be a compact manifold and $A\in\CI(Z;\Psiisoph{0}(\R;\C^N))$ a smooth family
  of elliptic operators.
  \begin{enumerate}[label=(\roman*)]
  \item If each $A(z)$ is invertible then $A^{-1}\in\CI(Z;\Psiisoph{0}(\R;\C^N))$.
  \item For $k$ large, $\Nul(A(z)(\Id-\Pi_k))=\Ran\Pi_k$ for all $z$; the orthocomplements
    $F_k(z)\subset\cS(\R;\C^N)$ of the ranges of $A(z)(\Id-\Pi_k)$ form a smooth vector
    bundle, and
    \begin{equation*}
      \indiso(A)=[\Ran\Pi_k]-[F_k]=[\C^{(k+1)N}]-[F_k]\in K^0(Z)
    \end{equation*}
    is independent of $k$; $F_k$ may be replaced by any smooth complement of the range of
    $A(\Id-\Pi_k)$ in $\cS(\R;\C^N)$. For $Z$ a point, $\indiso(A)=\ind A$.
  \item $\indiso$ is invariant under homotopies of elliptic families, depends only on
    the homotopy class of $\sigma_0(A)$ in $\CI(Z\times\Sphere;\GL(N,\C))$, is additive
    under direct sums, unchanged by stabilization $A\mapsto A\oplus\Id$, natural under
    pull-back, vanishes for invertible families, and $\indiso(bAc)=\indiso(A)$ for
    $b,c\in\CI(Z;\GL(N,\C))$.
  \item $\indiso(A)=0$ if and only if $A\oplus\Id_L+B$ is invertible for some $L$ and some
    smooth family $B$ in $\Psiiso{-\infty}(\R;\C^{N+L})$.
  \end{enumerate}
\end{proposition}

In (iii), two families with homotopic symbols are joined by the linear homotopy from
$A$ to $\Op_\chi(\sigma_0(A))$, which has constant symbol, followed by
$\Op_\chi(a_t)$ for a homotopy $a_t$ of symbols. The following lemma, valid in both the
Toeplitz and the isotropic setting, is the most convenient way to compute families
indices.

\begin{lemma}\label{bott:lem:null-complement}
  Let $A$ be a smooth elliptic family over a compact manifold $Z$, either of Toeplitz
  operators on $V=\Cp(\Sphere;\C^N)$ or of isotropic operators on $V=\cS(\R;\C^N)$.
  Suppose that the null spaces $\Nul A(z)$ form a smooth vector bundle $\mathcal N$
  (local frames given by smooth maps into $V$) and that $W$ is a smooth vector bundle
  with $V=A(z)V\oplus W(z)$ for all $z$. Then $\ind(A)=[\mathcal N]-[W]$.
\end{lemma}
\begin{proof}
  Let $P=P_k$, resp.\ $P=\Pi_k$, with $k$ admissible, and $A'=A(\Id-P)$. If
  $n\in\mathcal N(z)$ and $Pn=0$, then $A'n=An=0$, so $n\in\Nul A'=\Ran P$ and $n=0$.
  Hence $P$ maps $\mathcal N$ isomorphically onto a smooth subbundle $P\mathcal N$ of
  the trivial bundle $\Ran P$; let $Y$ be its orthocomplement in $\Ran P$, again a
  smooth bundle. We claim that $A(V)=A(Y)\oplus A'(V)$ with $A$ injective on $Y$. Indeed,
  if $y\in Y$ and $Ay=A'v$, then $y-(\Id-P)v\in\mathcal N$, and applying $P$,
  $y\in P\mathcal N\cap Y=\{0\}$. Conversely, since $V=\Ran P+\Ran(\Id-P)$, every element
  of $A(V)$ is $Ar+A'v$ with $r\in\Ran P$; writing $r=y+Pn$ with $y\in Y$,
  $n\in\mathcal N$, we get $Ar=Ay-A'n$. Thus $W\oplus A(Y)$ is a complement for $(A,k)$;
  it is smooth since $A$ maps local frames of $Y$ to local frames of $A(Y)$. By
  \cref{bott:lem:index-bundle} (or \cref{bott:prop:iso-families})
  $\ind(A)=[\Ran P]-[W]-[Y]=[P\mathcal N]-[W]=[\mathcal N]-[W]$.
\end{proof}

\subsection{The annihilation operator}
The operator $x+\partial_x=i(D_x-ix)$ is $i$ times the annihilation operator of the
harmonic oscillator. Its symbol is $x+i\xi$, of order one; we normalize it to order
$0$.

\begin{lemma}\label{bott:lem:annihilation}
  Let $B=\frac12\bigl(\Op(\ang w^{-1})+\Op(\ang w^{-1})^*\bigr)$ and
  $\ann=(x+\partial_x)B$. Then $\ann\in\Psiisoph{0}(\R)$ is elliptic, $\sigma_0(\ann)=z$,
  and $\ind\ann=1$. Moreover there is a finite rank $K\in\Psiiso{-\infty}(\R)$ such that
  $\ann'=\ann+K$ is surjective on $\cS(\R)$ and on $L^2(\R)$, with
  $\Nul\ann'=\C\phi_0$ for some $\phi_0\in\cS(\R)$ with $\norm{\phi_0}=1$; then
  $\ann'^*$ is injective, with range $\{f;\ (f,\phi_0)=0\}$.
\end{lemma}
\begin{proof}
  $B\in\Psiisoph{-1}(\R)$ is formally self-adjoint with principal symbol $\abs w^{-1}$, so it is
  elliptic of order $-1$, and $\ann$ is polyhomogeneous of order $0$ with principal symbol
  $(x+i\xi)/\abs w=z$. By \ref{bott:I:fredholm} (which holds for elliptic operators of any
  order on $\cS$), $B$ is Fredholm on $\cS(\R)$ with index $\dim\Nul B-\dim\Nul B^*=0$.
  The operator $x+\partial_x$ is elliptic of order one; $(x+\partial_x)u=0$ means
  $\partial_x(e^{x^2/2}u)=0$, so $\Nul(x+\partial_x)=\C e^{-x^2/2}$, while the adjoint
  $x-\partial_x$ has no tempered null space ($e^{x^2/2}$ is not tempered), so
  $x+\partial_x$ is surjective on $\cS$ and its index is $1$. For linear maps with finite
  dimensional null spaces and cokernels the index of a composite is the sum of the
  indices, so $\ind\ann=1$ on $\cS$, and also on $L^2$ by \ref{bott:I:fredholm}.

  Let $N_0=\Nul\ann$, $N_1=\Nul\ann^*$, both in $\cS$, with $\dim N_0=\dim N_1+1$. Choose
  a unit vector $\phi_0\in N_0$, let $N_0'=N_0\ominus\C\phi_0$, and let $K$ map $N_0'$
  isomorphically onto $N_1$ and vanish on $(N_0')^\perp$; $K$ has finite rank with range
  and co-range in $\cS$, so $K\in\Psiiso{-\infty}(\R)$. If $(\ann+K)u=0$, write $u=u'+u''$
  with $u'\in N_0'$, $u''\perp N_0'$; then $\ann u''\in\Ran\ann\perp N_1\ni Ku'$, so both
  vanish, $u'=0$ and $u''\in N_0\ominus N_0'=\C\phi_0$. The range of $\ann+K$ contains
  $\Ran\ann$ and $N_1$, hence everything. The statements about $\ann'^*$ follow from
  \ref{bott:I:fredholm}.
\end{proof}

\subsection{Normal forms and Bott periodicity}

\begin{proposition}\label{bott:prop:iso-normal}
  Let $\pi_-,\pi_0,\pi_+\colon Z\to M(M,\C)$ be smooth commuting idempotents with
  $\pi_-+\pi_0+\pi_+=\Id$, $V_\bullet=\Ran\pi_\bullet$. The family
  \begin{equation*}
    Q=\pi_-\ann'^*+\pi_0+\pi_+\ann',\qquad\sigma_0(Q)=\pi_-z^{-1}+\pi_0+\pi_+z,
  \end{equation*}
  has $\Nul Q=\phi_0\otimes V_+$, the bundle $\phi_0\otimes V_-$ is a complement of its
  range, and
  \begin{equation*}
    \indiso(Q)=[V_+]-[V_-]\in K^0(Z).
  \end{equation*}
\end{proposition}
\begin{proof}
  The matrices act on the factor $\C^M$ of $\cS(\R;\C^M)=\cS(\R)\otimes\C^M$ and commute
  with $\ann'$, $\ann'^*$, so with respect to
  $\cS(\R;\C^M)=\cS(\R)\otimes V_-\oplus\cS(\R)\otimes V_0\oplus\cS(\R)\otimes V_+$ the
  operator $Q$ is $\ann'^*\otimes\Id\oplus\Id\oplus\ann'\otimes\Id$. By
  \cref{bott:lem:annihilation} the null space is $\phi_0\otimes V_+$ and the range is
  $\{u\in\cS(\R)\otimes V_-;\ (u,\phi_0)=0\}\oplus\cS(\R)\otimes V_0\oplus\cS(\R)\otimes V_+$,
  where $(u,\phi_0)\in V_-$ denotes the pairing in the $\R$-variable; so
  $\phi_0\otimes V_-$ is a complement. Both are smooth bundles, with local frames
  $\phi_0\otimes\pi_\pm e_j$ (\cref{bott:lem:bundles}). Apply
  \cref{bott:lem:null-complement}.
\end{proof}

\begin{corollary}[Index formula for isotropic operators on $\R$]\label{bott:cor:iso-index-formula}
  If $A\in\Psiisoph{0}(\R;\C^N)$ is elliptic, then
  \begin{equation*}
    \ind A=\wind\sigma_0(A)=\frac1{2\pi i}\int_0^{2\pi}\Tr\Big(\sigma_0(A)^{-1}\frac{d\sigma_0(A)}{d\theta}\Big)d\theta,
  \end{equation*}
  the circle at infinity being traversed counterclockwise in the $(x,\xi)$-plane.
\end{corollary}
\begin{proof}
  As in the proof of \cref{bott:cor:index-circle} we may assume $\sigma_0(A)$ pointed and,
  by \cref{bott:thm:linearization} and \cref{bott:prop:iso-families}(iii), in normal form
  $\pi_-z^{-1}+\pi_0+\pi_+z$ with constant idempotents. Then
  $\ind A=\rank\pi_+-\rank\pi_-$ by \cref{bott:prop:iso-normal}, and this is the winding
  number of $\det(\pi_-z^{-1}+\pi_0+\pi_+z)=z^{\rank\pi_+-\rank\pi_-}$.
\end{proof}

For example $x+\partial_x$, the generator of the index, has symbol of winding number
one, while the creation operator $x-\partial_x$ has symbol $x-i\xi$ of winding number
$-1$, in accordance with $\Nul(x-\partial_x)=0$ and $\Nul(x+\partial_x)=\C e^{-x^2/2}$.

For a pointed loop $a\colon X\times\Sphere\to\GL(N,\C)$, $X$ compact, set
$\indiso(a)=\indiso(\Op_\chi(a))\in K^0(X)$. By \cref{bott:prop:iso-families}(iii) and
\cref{bott:lem:stable} this defines a homomorphism
\begin{equation*}
  \indiso\colon K^{-2}(X)\longrightarrow K^0(X),
\end{equation*}
natural under pull-back.

\begin{theorem}[Bott periodicity, isotropic form]\label{bott:thm:iso-bott}
  For every compact manifold $X$ the isotropic index map
  $\indiso\colon K^{-2}(X)\to K^0(X)$ is an isomorphism. Its inverse is
  \begin{equation}\label{bott:eq:iso-bott-map}
    \beta_{\mathrm{iso}}([E_+]-[E_-])=\big[\big(p_+z+(\Id-p_+)\big)\oplus\big((\Id-p_-)+p_-z^{-1}\big)\big],
  \end{equation}
  where $E_\pm\cong\Ran p_\pm$ for smooth families of projections $p_\pm$. The element
  $[z]\in K^{-2}(\mathrm{pt})$, the class of the symbol of $\ann$, corresponds to
  $[\C]\in K^0(\mathrm{pt})$.
\end{theorem}
\begin{proof}
  The loop in \eqref{bott:eq:iso-bott-map} is in normal form with $\pi_+=p_+\oplus0$ and
  $\pi_-=0\oplus p_-$, so by \cref{bott:prop:iso-normal} its isotropic index is
  $[E_+]-[E_-]$; thus $\indiso\circ\beta_{\mathrm{iso}}=\Id$ and $\indiso$ is surjective.
  For injectivity let $\indiso(a)=0$. By \cref{bott:thm:linearization} we may assume
  $a=\pi_-z^{-1}+\pi_0+\pi_+z$, and then $[V_+]=[V_-]$ by \cref{bott:prop:iso-normal}.
  The argument of \cref{bott:prop:injective}, which only involves loops (and
  \cref{bott:lem:normal-trivial}, which is symmetric in $\pi_\pm$), shows that $a$ is
  stably pointed homotopic to $\Id$.
\end{proof}

The proof shows that the only analytic input needed for Bott periodicity, besides the
linearization of loops, is the computation of the index of one operator of index one,
here $\ann$, and of its twisted versions in \cref{bott:prop:iso-normal}.

\begin{proposition}\label{bott:prop:iso-toeplitz}
  Let $\rho\colon\Sphere\to\Sphere$, $\theta\mapsto-\theta$. For every pointed loop $a$ over
  a compact manifold $X$,
  \begin{equation*}
    \indiso(a)=\ind(a\circ\rho)=-\ind(a),
  \end{equation*}
  where $\ind$ is the Toeplitz index map of \cref{bott:sec:index-iso}. In other words,
  after the orientation reversing identification of the circle at infinity of the phase
  plane with the circle $\Sphere_+$ of the Toeplitz calculus, the isotropic and the Toeplitz
  index maps coincide.
\end{proposition}
\begin{proof}
  All three expressions define homomorphisms $K^{-2}(X)\to K^0(X)$, so by
  \cref{bott:thm:linearization} it suffices to evaluate them on
  $a=\pi_-z^{-1}+\pi_0+\pi_+z$. Then $a\circ\rho=\pi_-z+\pi_0+\pi_+z^{-1}$, and by
  \cref{bott:prop:index-properties}(v), \cref{bott:prop:iso-normal},
  $\ind(a\circ\rho)=[V_+]-[V_-]=\indiso(a)=-\ind(a)$.
\end{proof}

\begin{remark}\label{bott:rem:iso-toeplitz-unitary}
  The comparison can be made at the level of operators. The unitary map
  $L^2(\R)\to L^2_+(\Sphere)$, $u_k\mapsto e^{ik\theta}$, carries $H$ to $2D_\theta+1$ and
  $\Psiiso{-\infty}(\R)$ onto $\PsiT{-\infty}(\Sphere)$ (both correspond to the matrices
  $(c_{jk})_{j,k\ge0}$ with rapidly decreasing entries), and $\ann'$ to an operator
  which, like $L$, is surjective with one-dimensional null space. The reversal of
  orientation reflects the fact that the annihilation operator lowers the eigenvalue of
  $H$, while $L$ lowers the Fourier mode, although their symbols are $z$ and $z^{-1}$.
\end{remark}

\begin{remark}[The rotation argument]\label{bott:rem:rotation-trick}
  Injectivity of the index map can also be proved without the linearization of loops,
  by an argument due to Atiyah. In the present setting it runs as follows. Let $a$ be a
  loop with $\indiso(a)=0$; by \cref{bott:prop:iso-families}(iv) there is (after
  stabilization) an \emph{invertible} family $A(x)$ with symbol $a(x,\cdot)$. Clutch $a$
  to a family of projections $\Pi_a$ on $\R^2\times X$, given by a matrix whose entries
  are affine in $a$ and $a^{-1}$ (as in \cref{bott:lem:sl-clutching} below). Using a
  loop $\beta$ of winding number one in a third variable $r$, the family
  $e=\beta(r)\Pi_a+\Id-\Pi_a$ of loops in $r$ has isotropic index the class of
  $\Pi_a$, i.e.\ the clutching class of $a$. Rotating the plane of $r$ and the variable
  of $a$ by a right angle does not change the index, but after the rotation the loop
  variable is the variable of $a$, and replacing $a$, $a^{-1}$ by $A$, $A^{-1}$ in the
  entries of $\Pi_a$ gives an invertible quantization; so the index, hence the clutching
  class of $a$, vanishes. Making this precise requires the families index over the
  non-compact parameter space $\R^2\times X$ with compact supports, and careful
  normalization at infinity; we do not pursue it.
\end{remark}

\section[The classifying sequence]{The stabilized Toeplitz group and the classifying sequence}\label{bott:sec:classifying}

We now pass from matrices to operators with values in the smoothing operators on
$Y$, to obtain groups which contain all the $\GL(N,\C)$ at once.

The aim is a short exact sequence of groups
$\GinfT\to\GT\to\GT/\GinfT$ in which the middle group is weakly contractible.
In topology, a contractible space $EG$ with a free action of a group $G$ has
quotient $BG$ with $\pi_k(BG)=\pi_{k-1}(G)$, a ``delooping'' of $G$. Here the
subgroup classifies odd K-theory, so the quotient, which is a space of symbols,
classifies even K-theory (\cref{bott:fig:classifying}). This is the analytic
form of the passage from $K^{-1}$ to $K^0$.

\subsection{Toeplitz operators with values in smoothing operators}
Let $Y$, $E$ be as in \cref{bott:sec:background}. On $\Psido{0}(\Sphere)$ we use the Fr\'echet
topology of \cref{ch:conormal,ch:manifolds}: an element is a conormal distribution at the
diagonal, i.e.\ the sum of a smooth kernel and, near the diagonal, the inverse Fourier
transform of a classical symbol, and the seminorms are those of uniform convergence of
all derivatives of these two pieces. Composition is continuous, and so is the inclusion
into the bounded operators on $L^2$.

\begin{definition}\label{bott:def:PsiTY}
  $\PsiT{0,-\infty}(\Sphere,Y;E)$ is the space of smooth sections
  $A\in\CI(Y^2;\PsiT{0}(\Sphere)\otimes\Hom(E)\otimes\Omega_R)$, i.e.\ smooth maps
  $(y,y')\mapsto A(y,y')$ with values in Toeplitz operators tensored with
  $\Hom(E_{y'},E_y)$ and a density in $y'$. It acts on $\CI(\Sphere\times Y;E)$, and on
  $\Cp(\Sphere\times Y;E)=S\CI(\Sphere\times Y;E)$, by
  $(Au)(\cdot,y)=\int_YA(y,y')u(\cdot,y')$, and is an algebra with
  $(AB)(y,y')=\int_YA(y,y'')\circ B(y'',y')$. The symbol
  \begin{equation*}
    \sigma_0\colon\PsiT{0,-\infty}(\Sphere,Y;E)\longrightarrow\CI(\Sphere;\Psido{-\infty}(Y;E)),
    \quad\sigma_0(A)(\theta)(y,y')=\sigma_0(A(y,y'))(\theta),
  \end{equation*}
  is multiplicative. The elements with $A(y,y')\in\PsiT{-\infty}(\Sphere)$ form the ideal
  $\PsiT{-\infty}(\Sphere\times Y;E)=S\Psido{-\infty}(\Sphere\times Y;E)S$, the Toeplitz
  smoothing operators on $\Sphere\times Y$.
\end{definition}

If $V\subset\CI(Y;E)$ has orthonormal basis $\phi_1,\dots,\phi_N$, then
$(A_{lm})\mapsto\sum_{l,m}A_{lm}\,\phi_l(y)\phi_m(y')^*$ embeds $\PsiT{0}(\Sphere;\C^N)$ as a
subalgebra (without unit); correspondingly an operator $A$ on $\Cp(\Sphere;\C^N)$ is
identified with $(\Id-\Pi_V)+A$ on $\Cp(\Sphere\times Y;E)$. Different choices of $V$ give
embeddings which are conjugate by unitaries joined to the identity, and enlarging $V$
corresponds to stabilization $A\mapsto A\oplus\Id$.

\begin{definition}\label{bott:def:GT}
  The \emph{stabilized normalized Toeplitz group} is
  \begin{multline*}
    \GT(\Sphere,Y;E)=\{\Id+A;\ A\in\PsiT{0,-\infty}(\Sphere,Y;E),\ \sigma_0(A)(1)=0,\\
    (\Id+A)^{-1}=\Id+B\text{ with }B\in\PsiT{0,-\infty}(\Sphere,Y;E)\},
  \end{multline*}
  and $\GinfT(\Sphere\times Y;E)$ is its subgroup of elements with
  $A\in\PsiT{-\infty}(\Sphere\times Y;E)$. Similarly $\GT(\Sphere;\C^N)$ denotes the group
  of invertible $A\in\PsiT{0}(\Sphere;\C^N)$ with $\sigma_0(A)(1)=\Id$, and
  $\GinfT(\Sphere;\C^N)$ the subgroup with $A-\Id\in\PsiT{-\infty}$.
\end{definition}

Elements of $\GT$ are elliptic: $\Id+\sigma_0(A)(\theta)$ has inverse $\Id+\sigma_0(B)(\theta)$, so
$\Id+\sigma_0(A)\in\Omega\Ginf(Y;E)$. The normalization $\sigma_0(A)(1)=0$ removes the
evaluation map $\GT\to\Ginf(Y;E)$ at $1\in\Sphere$, which would otherwise carry odd
K-theory. A map $f\colon X\to\GT(\Sphere,Y;E)$ is smooth if $f-\Id$ is smooth into the
Fr\'echet space $\PsiT{0,-\infty}(\Sphere,Y;E)$; by \cref{bott:lem:inverse-Y} below, $f^{-1}$ is then
smooth too, and products of smooth maps are smooth.

\begin{lemma}\label{bott:lem:inverse-Y}
  Let $A\in\CI(Z;\PsiT{0,-\infty}(\Sphere,Y;E))$ be a smooth family such that, for
  all $z\in Z$ and $\theta\in\Sphere$, $\Id+\sigma_0(A(z))(\theta)\in\Ginf(Y;E)$ and
  $\Id+A(z)$ is invertible on $L^2_+(\Sphere\times Y;E)$. Then $(\Id+A)^{-1}=\Id+B$ with
  $B\in\CI(Z;\PsiT{0,-\infty}(\Sphere,Y;E))$.
\end{lemma}
\begin{proof}
  The proof of \cref{bott:lem:inverse} applies verbatim. The parametrix construction
  goes through for operators with values in $\Psido{-\infty}(Y;E)$; the symbol of a parametrix
  is $(\Id+\sigma_0(A))^{-1}-\Id\in\CI(\Sphere;\Psido{-\infty}(Y;E))$, inversion in
  $\Ginf(Y;E)$ being smooth, and the lower order corrections are obtained by the same
  asymptotic summation. This gives $\Id+B_0$ with errors $R,R'$ in
  $\PsiT{-\infty}(\Sphere\times Y;E)$, and then
  \begin{equation*}
    V=(\Id+A)^{-1}=(\Id+B_0)+R(\Id+B_0)+RVR'.
  \end{equation*}
\end{proof}

\begin{lemma}\label{bott:lem:open}
  Let $A\in\CI(X;\PsiT{0,-\infty}(\Sphere,Y;E))$ with $\Id+A(x)\in\GT$ for all $x$, $X$
  compact. If $A'\in\CI(X;\PsiT{0,-\infty})$ satisfies $\sigma_0(A')(1)=0$ and is
  sufficiently close to $A$ uniformly on $X$ (in a suitable seminorm), then
  $\Id+A'(x)\in\GT$ for all $x$. The same holds for $\GT(\Sphere;\C^N)$ and for the groups
  $\GinfT$.
\end{lemma}
\begin{proof}
  The symbols $\Id+\sigma_0(A')(\theta)$ are uniformly close to the compact set of values of
  $\Id+\sigma_0(A)$ in the open set $\Ginf(Y;E)$, so $\Id+A'$ is elliptic; the operators
  are close in $L^2$ operator norm to invertible ones, so are invertible on
  $L^2_+(\Sphere\times Y;E)$. By \cref{bott:lem:inverse-Y} (or \cref{bott:lem:inverse}) the inverses lie in the algebra.
\end{proof}

\subsection{Weak contractibility}

\begin{proposition}\label{bott:prop:deform-smoothing}
  Let $g\colon X\to\GT(\Sphere;\C^N)$ be smooth. Then for some $M\ge N$, $g\oplus\Id_{M-N}$ is
  smoothly homotopic in $\GT(\Sphere;\C^M)$ to a map into $\GinfT(\Sphere;\C^M)$.
\end{proposition}
\begin{proof}
  The symbol $a=\sigma_0(g)$ is a pointed loop, and $\ind(a)=\ind(g)=0$ since $g$ is
  invertible (\cref{bott:prop:index-properties}(iv)). By \cref{bott:prop:injective} there is a
  pointed homotopy $a_t$, $t\in[0,1]$, in some $\GL(M,\C)$ from $a\oplus\Id$ to $\Id$. The
  family $A_t=g\oplus\Id+T_{a_t-a_0}$ is elliptic with symbol $a_t$, and $A_0=g\oplus\Id$
  is invertible. By \ref{bott:B:lift} there is a smooth family $B_t$ of Toeplitz smoothing
  operators with $B_0=0$ such that $A_t+B_t$ is invertible; its inverse is smooth by
  \cref{bott:lem:inverse}, and its symbol $a_t$ is pointed. So $A_t+B_t$ is a homotopy
  in $\GT(\Sphere;\C^M)$, and $A_1+B_1$ has symbol $\Id$, i.e.\ lies in $\GinfT$.
\end{proof}

\begin{theorem}\label{bott:thm:contractible}
  The group $\GT(\Sphere,Y;E)$ is weakly contractible: every smooth map
  $f\colon X\to\GT(\Sphere,Y;E)$ from a compact manifold is smoothly homotopic in
  $\GT(\Sphere,Y;E)$ to the constant map $\Id$.
\end{theorem}
\begin{proof}
  \emph{Step 1: reduction to matrices.} Write $f=\Id+A$. Let $\Pi_j$ be the projections of
  \ref{bott:B:finite}. The expansion of the smooth map $(x,y,y')\mapsto A(x)(y,y')$ with values
  in the Fr\'echet space $\PsiT{0}(\Sphere)$ (tensored with $\Hom(E)$) shows, exactly as
  for smoothing operators, that $\Pi_jA\Pi_j\to A$ uniformly on $X$ in
  $\PsiT{0,-\infty}$. Since $\sigma_0(\Pi_jA\Pi_j)(1)=\Pi_j\sigma_0(A)(1)\Pi_j=0$,
  \cref{bott:lem:open} shows that $\Id+(1-t)A+t\Pi_jA\Pi_j$, $t\in[0,1]$, is a homotopy in
  $\GT$ for $j$ large. Its end point is $(\Id-\Pi_j)+g$, with
  $g\colon X\to\GT(\Sphere;\C^N)$, $N=\dim V_j$.

  \emph{Step 2: reduction to smoothing operators.} By
  \cref{bott:prop:deform-smoothing}, after enlarging $V_j$ (which does not change the
  element of $\GT(\Sphere,Y;E)$), $g$ is homotopic in $\GT(\Sphere;\C^M)$ to a map
  $g'=\Id+K$, $K\in\CI(X;\PsiT{-\infty}(\Sphere;\C^M))$.

  \emph{Step 3: reduction to the zero Fourier mode.} By \ref{bott:B:shifts},
  $P_kKP_k\to K$ uniformly, so by \cref{bott:lem:open} $\Id+(1-t)K+tP_kKP_k$ is a homotopy in
  $\GinfT(\Sphere;\C^M)$ for $k$ large. Its end point acts on $\Ran P_k$, spanned by
  $e^{ij\theta}e_l$, $0\le j\le k$, $1\le l\le M$, and is the identity on $\Ran(\Id-P_k)$.
  Stabilize once more, to $\C^{(k+1)M}$ with orthonormal basis $e_{l,j}$,
  $1\le l\le M$, $0\le j\le k$, where $e_{l,0}=e_l$. For $1\le j\le k$ the pairs
  $(e^{ij\theta}e_{l,0},e_{l,j})$, the second being a constant function, are orthonormal
  and span mutually orthogonal planes. Let $\mathcal R_\tau$ be the unitary rotating each
  of these planes by the angle $\tau$ and fixing their orthocomplement. Then
  $\mathcal R_\tau-\Id$ has finite rank and smooth kernel, so
  $\mathcal R_\tau\in\GinfT(\Sphere;\C^{(k+1)M})$, and conjugation by $\mathcal R_\tau$,
  $\tau\in[0,\pi/2]$, is a homotopy in $\GinfT$ ending at an operator of the form
  \begin{equation*}
    g''=\Id-P_0+P_0hP_0,\qquad h\colon X\to\GL((k+1)M,\C),
  \end{equation*}
  since $\mathcal R_{\pi/2}$ maps $\Ran P_k\otimes e_{\cdot,0}$ onto the constants $\Ran P_0$.

  \emph{Step 4: rotating away the matrix.} For $h\in\GL(N',\C)$ consider on
  $\Cp(\Sphere;\C^{N'})^{2}$ the two curves of Toeplitz operators, $c=\cos\tau$, $s=\sin\tau$,
  \begin{equation*}
    M_\tau=\begin{pmatrix}c(\Id-P_0)+P_0hP_0&s\,hU\\-s\,h^{-1}L&c\end{pmatrix},\quad
    M'_\tau=\begin{pmatrix}c\,h(\Id-P_0)+P_0hP_0&s\,hU\\-s\,h^{-1}L&c\,h^{-1}\end{pmatrix},
  \end{equation*}
  $\tau\in[0,\pi/2]$; they agree at $\tau=\pi/2$, $M_0=\diag(g'',\Id)$ and $M'_0=\diag(h,h^{-1})$.
  Both are invertible: on $(u_ke^{ik\theta},v_ke^{i(k-1)\theta})$, $k\ge1$, they act by the
  matrices $\begin{pmatrix}c&sh\\-sh^{-1}&c\end{pmatrix}=\diag(h,1)R'_\tau\diag(h^{-1},1)$ and
  $\diag(h,h^{-1})R'_\tau$, $R'_\tau=\begin{pmatrix}c&s\\-s&c\end{pmatrix}$, and they map
  $(u_0,0)$ to $(hu_0,0)$. Let $C_\tau=\sigma_0(M_\tau)(1)$, a constant invertible matrix (it
  commutes with $S$). Then $C_\tau^{-1}M_\tau$ followed by $(C'_\tau)^{-1}M'_\tau$ (run from
  $\tau=\pi/2$ to $0$, $C'_\tau=\sigma_0(M'_\tau)(1)$) is a homotopy in $\GT(\Sphere;\C^{2N'})$,
  smooth in $h=h(x)$, from $\diag(g'',\Id)$ to $\diag(h,h^{-1})^{-1}\diag(h,h^{-1})=\Id$.
  (This is the deformation of \cref{ch:index}, normalized at $1\in\Sphere$.)

  Composing the homotopies of Steps 1--4, all of which take place in $\GT(\Sphere,Y;E)$ via
  the embeddings above, gives a homotopy from $f$ to $\Id$.
\end{proof}

\begin{remark}
  Continuous maps from compact manifolds into $\GT$ can be approximated by smooth ones,
  and nearby maps are homotopic (linear homotopy and \cref{bott:lem:open}), so all homotopy
  groups of $\GT(\Sphere,Y;E)$ vanish. Whether $\GT$ is actually contractible is discussed in
  the Notes.
\end{remark}

\subsection{The classifying sequence}

To describe the quotient $\GT/\GinfT$ we use full symbols. On the circle, with the
normal fibration of the diagonal given by $\theta-\theta'$, an operator
$A\in\Psido{0}(\Sphere)$ has, modulo $\Psido{-\infty}$, a kernel given near the diagonal by
the inverse Fourier transform of a classical symbol $a(\theta,\xi)$; its Taylor series at
$\xi=+\infty$ in $\rho=1/\xi$ is the \emph{full symbol} at $\Sphere_+$,
$\Phi(A)=\sum_{j\ge0}\rho^ja_j(\theta)$, with $a_0=\sigma_0(A)|_{\Sphere_+}$. For Toeplitz
operators the full symbol at $\Sphere_-$ vanishes, so by Borel's lemma (asymptotic
completeness, \cref{ch:conormal}) there is a short exact sequence
\begin{equation*}
  \PsiT{-\infty}(\Sphere)\longrightarrow\PsiT{0}(\Sphere)\xrightarrow{\ \Phi\ }\CI(\Sphere)[[\rho]],
\end{equation*}
and by the composition formula $\Phi(AB)=\Phi(A)*\Phi(B)$ for a bilinear product $*$ on
formal series whose leading term is the pointwise product. Applying $\Phi$ to the kernels
$A(y,y')$ gives, for $\PsiT{0,-\infty}(\Sphere,Y;E)$, the exact sequence
\begin{equation}\label{bott:eq:full-symbol}
  \PsiT{-\infty}(\Sphere\times Y;E)\longrightarrow\PsiT{0,-\infty}(\Sphere,Y;E)
  \xrightarrow{\ \Phi\ }\CI(\Sphere;\Psido{-\infty}(Y;E))[[\rho]].
\end{equation}

\begin{lemma}\label{bott:lem:index-Y}
  If $A\in\PsiT{0,-\infty}(\Sphere,Y;E)$ and $\Id+\sigma_0(A)$ is a loop in $\Ginf(Y;E)$, then
  $\Id+A$ is Fredholm on $L^2_+(\Sphere\times Y;E)$ and $\ind(\Id+A)=-\wind(\Id+\sigma_0(A))$.
  If this index vanishes there is $K\in\PsiT{-\infty}(\Sphere\times Y;E)$ of finite rank with
  $\Id+A+K$ invertible; if moreover $\sigma_0(A)(1)=0$ then $\Id+A+K\in\GT(\Sphere,Y;E)$.
\end{lemma}
\begin{proof}
  A parametrix (\cref{bott:lem:inverse-Y}) inverts $\Id+A$ modulo compact operators, so it is
  Fredholm, with null space and orthocomplement of the range (the null space of the
  adjoint, which is of the same type) consisting of smooth sections. For the index, the
  family $\Id+(1-t)A+t\Pi_jA\Pi_j$ consists of elliptic, hence Fredholm, operators for $j$
  large and is continuous in norm, so the index is constant; the winding number of the
  symbol is constant as well. At $t=1$ the operator is $\Id-\Pi_j$ plus an elliptic Toeplitz
  operator on $\C^{\dim V_j}$, and \cref{bott:cor:index-circle} applies. If the index
  vanishes, let $K$ map an orthonormal basis of the null space to one of the orthocomplement
  of the range and vanish on the orthocomplement of the null space; it has finite rank and
  smooth kernel, $\Id+A+K$ is bijective on $L^2_+$, and \cref{bott:lem:inverse-Y} applies.
\end{proof}

\begin{theorem}[Classifying sequence]\label{bott:thm:classifying}
  For $Y$, $E$ as above there is a short exact sequence of groups
  \begin{equation}\label{bott:eq:classifying}
    \GinfT(\Sphere\times Y;E)\longrightarrow\GT(\Sphere,Y;E)\xrightarrow{\ \Phi\ }\mathcal Q,
  \end{equation}
  where $\mathcal Q$ is the set of formal series $\sum_{j\ge0}\rho^jb_j$ with
  $b_j\in\CI(\Sphere;\Psido{-\infty}(Y;E))$ and $\Id+b_0\in\Omega_0\Ginf(Y;E)$ (pointed loops of
  winding number zero), with the group structure induced by $*$. Moreover, for every compact
  manifold $X$:
  \begin{enumerate}[label=(\roman*)]
  \item $[X;\GinfT(\Sphere\times Y;E)]\cong K^{-1}(X)$;
  \item $\GT(\Sphere,Y;E)$ is weakly contractible;
  \item if $X$ is connected, $[X;\mathcal Q]\cong[X;\Omega_0\Ginf(Y;E)]\cong\tilde K^0(X)$,
    where $\mathcal Q$ has the product topology of the coefficients.
  \end{enumerate}
  Thus $\GinfT$ classifies odd K-theory, $\GT$ is weakly contractible, and $\mathcal Q$
  classifies reduced even K-theory.
\end{theorem}

\begin{figure}[ht]
\centering
\begin{tikzpicture}[
  box/.style={draw,rounded corners,align=center,inner sep=5pt,minimum height=1.1cm},
  note/.style={font=\footnotesize,align=center,text width=3.6cm},
  >=Stealth]
\node[box] (a) at (0,0) {$\GinfT(\Sphere\times Y;E)$\\ \footnotesize$\Id+{}$smoothing};
\node[box] (b) at (4.7,0) {$\GT(\Sphere,Y;E)$\\ \footnotesize invertible Toeplitz};
\node[box] (c) at (9.4,0) {$\mathcal Q$\\ \footnotesize full symbols};
\draw[{Hooks[right]}->] (a)--(b);
\draw[->>] (b)--node[above]{$\Phi$}(c);
\node[note] at (0,-1.4) {$[X;\GinfT]=K^{-1}(X)$};
\node[note] at (4.7,-1.4) {weakly contractible};
\node[note] at (9.4,-1.4) {$[X;\mathcal Q]=\tilde K^0(X)$\\ ($X$ connected)};
\node[note] at (0,1.2) {kernel of the symbol};
\node[note] at (9.4,1.2) {leading term: loop in $\Ginf(Y;E)$ of winding $0$};
\end{tikzpicture}
\caption{The classifying sequence of \cref{bott:thm:classifying}, a group-level
version of the symbol sequence $\Psi^{-\infty}\to\Psi^0\to$ symbols. Ellipticity
and invertibility are built into the middle group, so the index obstruction of
\eqref{index:eq:group-sequence} has been removed by restricting the quotient to
loops of winding number zero.}
\label{bott:fig:classifying}
\end{figure}

\begin{question}
  Why is a weakly contractible middle group useful?
\end{question}
\begin{answer}
  It converts operators into topology. A map $X\to\mathcal Q$ is a family of
  symbols; it lifts to a family of invertible operators in $\GT$, and any two
  lifts, or any lift of a homotopy, are related through the contractible group
  $\GT$. The ambiguity in the lift is a map into the subgroup $\GinfT$, which
  carries odd K-theory. If the sequence were a fibration this would give
  $\pi_k(\mathcal Q)\cong\pi_{k-1}(\GinfT)$, i.e.\ the delooping relation
  between even and odd K-theory, as in the remark below.
\end{answer}
\begin{proof}
  By \eqref{bott:eq:full-symbol} the kernel of $\Phi$ on $\GT$ is $\GinfT(\Sphere\times Y;E)$
  (an element of $\GT$ with $A$ smoothing has inverse $\Id+B$ with
  $B=-A-AB\in\PsiT{-\infty}$). If $\Id+A\in\GT$, then $\Id+\sigma_0(A)$ is a pointed loop in
  $\Ginf(Y;E)$, and its winding number is $-\ind(\Id+A)=0$ by \cref{bott:lem:index-Y}; so
  $\Phi(\GT)\subset\mathcal Q$. Conversely, given $\sum\rho^jb_j\in\mathcal Q$, Borel's lemma
  gives $A$ with $\Phi(A)=\sum\rho^jb_j$, and by \cref{bott:lem:index-Y} there is
  $K\in\PsiT{-\infty}(\Sphere\times Y;E)$ with $\Id+A+K\in\GT$; since $\Phi(K)=0$, $\Phi$ is
  surjective. This proves exactness, and $\Phi$ is a homomorphism for the product $*$.

  (i) $\GinfT(\Sphere\times Y;E)$ is the group of invertible perturbations of the identity
  by smoothing operators on the Hilbert space $L^2_+(\Sphere\times Y;E)$. The orthonormal
  system $e^{il\theta}\phi_m$, $l\ge0$, with $\phi_m$ the eigensections of
  \ref{bott:B:finite}, gives finite rank approximations exactly as in \ref{bott:B:finite}, so
  the proof of \cref{bott:lem:stable} identifies $[X;\GinfT]$ with the stable homotopy
  classes of maps $X\to\GL(N,\C)$, which is $K^{-1}(X)$.

  (ii) is \cref{bott:thm:contractible}.

  (iii) As a space, $\mathcal Q=\Omega_0\Ginf(Y;E)\times\prod_{j\ge1}\CI(\Sphere;\Psido{-\infty}(Y;E))$,
  and the second factor is a vector space; scaling it by $t\in[0,1]$ shows that projection
  onto $b_0$ induces a bijection $[X;\mathcal Q]\to[X;\Omega_0\Ginf]$. A smooth map
  $X\to\Omega\Ginf$ is the same as a pointed map $X\times\Sphere\to\Ginf$; for connected $X$
  its winding number is constant, and it is zero exactly when the rank of the index class
  is zero. By \cref{bott:thm:index-iso} the maps into $\Omega_0\Ginf$ thus correspond to
  $\{\xi\in K^0(X);\ \rank\xi=0\}=\tilde K^0(X)$.
\end{proof}

\begin{remark}
  If \eqref{bott:eq:classifying} were known to be a fibration, the long exact sequence of
  homotopy groups and (ii) would give $\pi_k(\mathcal Q)\cong\pi_{k-1}(\GinfT)$. This is
  consistent with the results above: by (iii) and \cref{bott:cor:homotopy-groups} below,
  $\pi_k(\mathcal Q)=\pi_k(\Omega_0\Ginf)=\pi_{k+1}(\Ginf)$, which by Bott periodicity equals
  $\pi_{k-1}(\Ginf)=\pi_{k-1}(\GinfT)$.
\end{remark}

\section[Compactly supported K-theory]{Compactly supported K-theory and Bott periodicity}\label{bott:sec:Kc}

The symbol of an elliptic operator on $Z$ lives on $T^*Z$, which is not compact,
and it is invertible only outside a compact set (the zero section). The natural
K-group for such data is K-theory with compact supports: pairs of bundles
together with an isomorphism between them away from a compact set. This section
defines these groups, relates them to the loop models by clutching, and proves
Bott periodicity in its usual form, $\Kc^0(\R^2\times X)\cong K^0(X)$.

\subsection{Definitions}
Let $M$ be a manifold, not necessarily compact. We write $K\Subset M$ for a compact subset.

\begin{definition}\label{bott:def:Kc}
  \begin{enumerate}[label=(\roman*)]
  \item $\Kc^{-1}(M)$ is the set of smooth maps $f\colon M\to\Ginf(Y;E)$ with $f=\Id$ on
    $M\setminus K$ for some $K\Subset M$, modulo homotopies $f_t$ with $f_t=\Id$ on
    $M\setminus K$ for one $K\Subset M$ and all $t$.
  \item $\Kc^{-2}(M)$ is defined in the same way, using maps $f\colon M\times\Sphere\to\Ginf(Y;E)$
    with $f=\Id$ on $(M\setminus K)\times\Sphere\cup M\times\{1\}$.
  \item $\Kc^0(M)$ is the set of equivalence classes of triples $(E_+,E_-,a)$, where $E_\pm$
    are vector bundles over $M$ and $a\colon E_+\to E_-$ is a smooth bundle isomorphism over
    $M\setminus K$ for some $K\Subset M$. The equivalence relation is generated by:
    \begin{enumerate}[label=(\alph*)]
    \item $(E_+,E_-,a)\sim(E'_+,E'_-,\phi_-a\phi_+^{-1})$ for bundle isomorphisms
      $\phi_\pm\colon E_\pm\to E'_\pm$ over $M$;
    \item $(E_+,E_-,a)\sim(E_+,E_-,a')$ if $a=a'$ outside a compact set;
    \item $(E_+,E_-,a_0)\sim(E_+,E_-,a_1)$ if $a_t$, $t\in[0,1]$, is a smooth family of
      isomorphisms over $M\setminus K$ for one $K\Subset M$;
    \item $(E_+,E_-,a)\sim(E_+\oplus H,E_-\oplus H,a\oplus\Id_H)$ for any bundle $H$.
    \end{enumerate}
    The sum is $[(E_+,E_-,a)]+[(F_+,F_-,b)]=[(E_+\oplus F_+,E_-\oplus F_-,a\oplus b)]$.
  \end{enumerate}
\end{definition}

The groups $\Kc^{-1}$, $\Kc^{-2}$ are abelian groups by pointwise multiplication (as in
\cref{bott:lem:stable}). $\Kc^0(M)$ is an abelian monoid with zero $[(0,0,\cdot)]$; a
triple $(F,F,c)$ with $c$ an isomorphism over all of $M$ represents $0$ (use (a) with
$\phi_-=c^{-1}$ and then (d)). It is a group: $(E_+\oplus E_-,E_-\oplus E_+,a\oplus a^{-1})$ is
equivalent by (a), composing with the exchange of the summands of $E_-\oplus E_+$, to
$(F,F,b)$ with $F=E_+\oplus E_-$ and $b=\begin{pmatrix}0&a^{-1}\\a&0\end{pmatrix}$, and
$b_\tau=\begin{pmatrix}\cos\tau&\sin\tau\,a^{-1}\\\sin\tau\,a&-\cos\tau\end{pmatrix}$ satisfies
$b_\tau^2=\Id$ and joins $b$ to the globally defined isomorphism $\diag(\Id,-\Id)$.

For compact $M$ the conditions at infinity are void, and $\Kc^{-1}(M)=K^{-1}(M)$,
$\Kc^{-2}(M)=K^{-2}(M)$; also $\Kc^0(M)=K^0(M)$, since $(E_+,E_-,a)$ may be taken with $a$
defined on the empty set and then (a), (d) are exactly the relation of \cref{bott:def:K0}.

\subsection{Flattening and suspension}
\begin{lemma}[Flattening]\label{bott:lem:flatten}
  Let $W_1\subset M_1$ and $W_2\subset M_2$ be points (or empty) in compact manifolds and
  $W=W_1\times M_2\cup M_1\times W_2$. Every smooth $f\colon M_1\times M_2\to\Ginf$ with $f=\Id$
  on $W$ is homotopic, through such maps, to one which is $\Id$ in a neighbourhood of $W$;
  and two maps which are $\Id$ near $W$ and homotopic through maps equal to $\Id$ on $W$ are
  homotopic through maps equal to $\Id$ near $W$.
\end{lemma}
\begin{proof}
  In a coordinate ball centred at $W_1=\{w\}$, let $\psi_t(y)=((1-t)+t\chi(|y|))y$, where
  $\chi\in\CI([0,\infty))$ vanishes on $[0,\delta]$ and equals $1$ on $[2\delta,\infty)$, and
  extend $\psi_t$ by the identity. Then $\psi_t$ fixes $w$, $\psi_0=\Id$ and $\psi_1$ maps a
  neighbourhood of $w$ to $w$. With the analogous map on $M_2$, $\Psi_t=\psi_t\times\psi_t$ maps
  $W$ into $W$ for all $t$, and $f\circ\Psi_t$ is the required homotopy. For the second
  statement apply $\Psi_1$ to the homotopy; $h\circ\Psi_t$ joins a map $h$ which is $\Id$
  near $W$ to $h\circ\Psi_1$ through maps equal to $\Id$ near $W$, provided $\delta$ is
  small (by compactness $\Psi_t$ maps a small neighbourhood of $W$ into the set where $h=\Id$).
\end{proof}

Applied with $M_1=\mathrm{pt}$ or with the stereographic identification
$\Sphere^k\setminus\{e\}\cong\R^k$ this shows that the homotopy groups
$\pi_k(\Ginf)=[\Sphere^k,e;\Ginf,\Id]$ (pointed smooth homotopy classes) satisfy
$\pi_k(\Ginf)\cong\Kc^{-1}(\R^k)$. Similarly, identifying $\Sphere\setminus\{1\}$ with $\R$:

\begin{lemma}\label{bott:lem:suspension}
  For every manifold $M$, $\Kc^{-2}(M)\cong\Kc^{-1}(\R\times M)$.
\end{lemma}
\begin{proof}
  By (the proof of) \cref{bott:lem:flatten} in the circle factor, representatives and
  homotopies in $\Kc^{-2}(M)$ may be taken equal to $\Id$ near $M\times\{1\}$; restricted to
  $M\times(\Sphere\setminus\{1\})\cong M\times\R$ they are exactly the representatives and
  homotopies of $\Kc^{-1}(\R\times M)$, and conversely such maps extend by $\Id$.
\end{proof}

\subsection{Clutching}
For $n\ge1$ let $\Sphere^{n-1}\subset\R^n$ be the unit sphere with base point $e$, and write
$K^{-1}(\Sphere^{n-1}\times X,X)$ for the set of homotopy classes of maps
$g\colon\Sphere^{n-1}\times X\to\Ginf(Y;E)$ with $g(e,x)=\Id$, through such maps. Thus
$K^{-1}(\Sphere^0\times X,X)=K^{-1}(X)$ and $K^{-1}(\Sphere\times X,X)=K^{-2}(X)$.

\begin{lemma}[Clutching]\label{bott:lem:clutching}
  For compact $X$ and $n\ge1$ there is a natural isomorphism
  \begin{equation*}
    \Kc^0(\R^n\times X)\cong K^{-1}(\Sphere^{n-1}\times X,X),
  \end{equation*}
  whose inverse sends the class of a pointed $g$ with values in $\GL(N,\C)$ to the class of
  $(\C^N,\C^N,\tilde g)$, $\tilde g(r\omega,x)=g(\omega,x)$ for $r\ge1$, $\omega\in\Sphere^{n-1}$.
\end{lemma}
\begin{proof}
  Let $\varpi\colon\R^n\times X\to X$. By \cref{bott:lem:bundle-homotopy}(ii) every bundle on
  $\R^n\times X$ is isomorphic to a pull-back, so by (a) every class is represented by
  $(\varpi^*F_+,\varpi^*F_-,a)$ with $a$ an isomorphism on $\{|v|\ge R\}$. Choose a bundle $G$
  and an isomorphism $\tau\colon F_-\oplus G\to\C^N$, and set
  \begin{equation*}
    g(\omega,x)=\tau\,(a_R\oplus\Id_G)(\omega,x)\,(a_R\oplus\Id_G)(e,x)^{-1}\tau^{-1},\qquad
    a_R(\omega,x)=a(R\omega,x),
  \end{equation*}
  where $a_R(\omega,x)\colon F_{+,x}\to F_{-,x}$. This is a pointed map into $\GL(N,\C)$; let
  $\kappa[(E_+,E_-,a)]$ be its class.

  \emph{Well-definedness.} Changing $R$, or $a$ by a homotopy (c), changes $g$ by a pointed
  homotopy, and (b) is covered by changing $R$. Under (d), or a change of $(G,\tau)$, $g$ is
  changed by stabilization and conjugation by a map $X\to\GL$ independent of $\omega$, which
  does not change the class by \cref{bott:lem:conjugation}. Under (a) with isomorphisms
  $\phi_\pm$ of the pull-back bundles, $\phi_\pm(s R\omega,x)$, $s\in[0,1]$, joins
  $\phi_\pm|_{\{|v|=R\}}$ to $\phi_\pm(0,x)$, which is independent of $\omega$; along this
  homotopy $g$ is changed by a pointed homotopy, and at $s=0$ it is replaced by its conjugate
  by $\tau(\phi_-(0,\cdot)\oplus\Id)\tau^{-1}$ and again the class does not change. $\kappa$ is
  additive, since direct sums of triples give block sums of the $g$'s.

  \emph{Surjectivity.} For pointed $g$ with values in $\GL(N,\C)$ (\cref{bott:lem:stable})
  the triple $(\C^N,\C^N,\tilde g)$ of the statement is mapped to $[g]$ (with $G=0$, $\tau=\Id$,
  $R=1$).

  \emph{Injectivity.} Using (d) with $\varpi^*G$ and then (a) with $\tau$ we may assume
  $F_-=\C^N$; since $a_R(e,x)\colon F_+\to\C^N$ is an isomorphism, (a) allows us to assume
  $F_+=\C^N$ and $a_R(e,x)=\Id$. The homotopy
  $a_s(r\omega,x)=a\big(((1-s)r+sR)\omega,x\big)$, $r\ge R$, shows that the triple is equivalent
  to $(\C^N,\C^N,\tilde g)$ with $\tilde g(r\omega,x)=g(\omega,x)$ for $r\ge R$. If $[g]=0$, then
  by \cref{bott:lem:stable} $g\oplus\Id$ is joined to $\Id$ by a pointed homotopy $g_t$ in some
  $\GL(N',\C)$, and (d), (c) give
  $(\C^N,\C^N,\tilde g)\sim(\C^{N'},\C^{N'},\tilde g_1)=(\C^{N'},\C^{N'},\Id)$, which represents
  $0$.
\end{proof}

\subsection{Bott periodicity}
\begin{theorem}[Bott periodicity]\label{bott:thm:bott}
  For every compact manifold $X$ there is a natural isomorphism
  \begin{equation*}
    K^0(X)\longrightarrow\Kc^0(\R^2\times X).
  \end{equation*}
  Identifying $\R^2=\C$ with variable $w$, it maps $[E]$ to the class of
  $(\varpi^*E,\varpi^*E,\bar w/|w|)$; that is, it is multiplication by the Bott class
  $b=[(\C,\C,\bar w/|w|)]\in\Kc^0(\R^2)$.
\end{theorem}
\begin{proof}
  The composite of the inverse of the index isomorphism (\cref{bott:thm:index-iso}) and of
  the inverse of the clutching isomorphism (\cref{bott:lem:clutching}, $n=2$),
  $K^0(X)\to K^{-2}(X)=K^{-1}(\Sphere\times X,X)\to\Kc^0(\R^2\times X)$, is an isomorphism. By
  \eqref{bott:eq:bott-map}, if $E=\Ran p$, $p\colon X\to M(N,\C)$, then $\beta[E]$ is
  represented by $pz^{-1}+(\Id-p)$, and the clutching construction gives
  $(\C^N,\C^N,p\,\bar w/|w|+\Id-p)$ for $|w|\ge1$, since $z^{-1}=\bar w/|w|$ on the circle
  $\omega=w/|w|$. Writing $\C^N=E\oplus\Ran(\Id-p)$ and using (d) this is equivalent to
  $(\varpi^*E,\varpi^*E,\bar w/|w|)$.
\end{proof}

\begin{discussion}
  The Bott class $b$ is the simplest compactly supported class on $\R^2$: the
  trivial line bundle, compared with itself by multiplication by $\bar w/|w|$
  outside the unit disc. Restricted to the circle at infinity, this is a loop of
  winding number $-1$, the symbol of the Toeplitz operator $L$ of index $1$.
  The theorem says that every compactly supported class on $\R^2\times X$ is
  obtained from a class on $X$ by ``twisting'' with this one loop, so the plane
  $\R^2$ contributes nothing but a shift.
\end{discussion}

\begin{remark}
  Periodicity can also be deduced from the chain
  $\Kc^0(\R^2\times X)\cong\Kc^{-1}(\R\times X)\cong\Kc^{-2}(X)\cong\Kc^0(X)$, which uses the
  clutching isomorphism for the non-compact manifold $\R\times X$ and the index isomorphism
  for non-compact $X$ (\cref{bott:ex:noncompact}). For compact $X$ the clutching lemma with
  $n=1$ gives $\Kc^0(\R\times X)\cong K^{-1}(X)$, and \cref{bott:lem:suspension} gives
  $K^{-2}(X)\cong\Kc^{-1}(\R\times X)$, so the chain is
  $\Kc^0(\R^2\times X)\cong K^{-2}(X)\cong\Kc^{-1}(\R\times X)$ and
  $\Kc^0(\R\times X)\cong K^{-1}(X)$.
\end{remark}

\subsection{The homotopy groups of \texorpdfstring{$\Ginf$}{G}}
Spheres need one more form of the clutching construction.

\begin{lemma}\label{bott:lem:sphere-clutching}
  For $m\ge1$, $\tilde K^0(\Sphere^m)\cong\pi_{m-1}(\Ginf(Y;E))$.
\end{lemma}

\begin{figure}[ht]
\centering
\begin{tikzpicture}[>=Stealth]
\shade[ball color=gray!4] (0,0) circle (1.8);
\draw[thick] (0,0) circle (1.8);
\draw[very thick] (-1.8,0) arc (180:360:1.8 and 0.45);
\draw[very thick,densely dashed] (1.8,0) arc (0:180:1.8 and 0.45);
\fill (1.8,0) circle (1.6pt) node[right,font=\footnotesize]{$e$};
\node at (0,1.0) {$U_+$};
\node at (0,-1.1) {$U_-$};
\node[font=\footnotesize] at (-2.85,0.95) {$U_+\times\C^N$};
\node[font=\footnotesize] at (-2.85,-0.95) {$U_-\times\C^N$};
\draw[->] (-2.3,0.75)--(-1.15,0.6);
\draw[->] (-2.3,-0.75)--(-1.15,-0.65);
\draw[->] (2.0,-0.9) node[right,font=\footnotesize,align=left]
  {equator $\Sphere^{m-1}$:\\ $(y,v)_-\sim(y,g(y)v)_+$}--(1.05,-0.4);
\end{tikzpicture}
\caption{The clutching construction. Over the two hemispheres $U_\pm$ of
$\Sphere^m$ (both contractible) a bundle is trivial; it is determined by how the
two trivializations are glued along the equator, i.e.\ by a map
$g:\Sphere^{m-1}\to\GL(N,\C)$. For $m=2$ the equator is a circle and $g$ is a
loop; the loop $g(z)=z$ gives, up to conjugation, the tautological line bundle
over $\Sphere^2=\mathbb{CP}^1$.}
\label{bott:fig:clutching}
\end{figure}
\begin{proof}
  Let $U_\pm=\{\pm x_{m+1}>-\epsilon\}\subset\Sphere^m\subset\R^{m+1}$, so that $U_\pm$ are
  diffeomorphic to $\R^m$ and $U_+\cap U_-$ retracts, by a map $r$, onto the equator
  $\Sphere^{m-1}$, which carries the base point $e$. For pointed $g\colon\Sphere^{m-1}\to\GL(N,\C)$
  let $E_g$ be the bundle obtained from $U_+\times\C^N$ and $U_-\times\C^N$ by identifying
  $(y,v)_-$ with $(y,g(r(y))v)_+$ over $U_+\cap U_-$ (\cref{bott:fig:clutching}), and set $\Psi(g)=[E_g]-[\C^N]$, an
  element of $\tilde K^0(\Sphere^m)$. A pointed homotopy $g_t$ defines a bundle over
  $[0,1]\times\Sphere^m$, so $E_{g_0}\cong E_{g_1}$ by \cref{bott:lem:bundle-homotopy};
  $E_{g\oplus\Id}=E_g\oplus\C$ and $E_{\diag(g,h)}=E_g\oplus E_h$. By \cref{bott:lem:stable}
  and \cref{bott:lem:rotation}, $\Psi$ induces a homomorphism
  $\pi_{m-1}(\Ginf)\to\tilde K^0(\Sphere^m)$.

  It is surjective: an element of $\tilde K^0(\Sphere^m)$ has the form $[E]-[\C^N]$ with
  $\rank E=N$ ($\Sphere^m$ is connected). $E$ is trivial over $U_\pm$
  (\cref{bott:lem:bundle-homotopy}(ii)); for trivializations $\tau_\pm$ the transition map
  $\tilde g=\tau_+\tau_-^{-1}$ on $U_+\cap U_-$ is homotopic to $\tilde g\circ r$ (deform along
  the retraction), so $E\cong E_g$ with $g=\tilde g|_{\Sphere^{m-1}}$; replacing $\tau_-$ by
  $g(e)\tau_-$ makes $g$ pointed.

  It is injective: if $[E_g]=[\C^N]$, then $E_g\oplus\C^L\cong\C^{N+L}$, i.e.\ $E_{g'}$,
  $g'=g\oplus\Id_L$, has a global trivialization. In the local trivializations over $U_\pm$ it is given by maps
  $\tau_\pm\colon U_\pm\to\GL(N+L,\C)$ with $\tau_-=\tau_+g'$, i.e.\ $g'=\tau_+^{-1}\tau_-$, on the equator. As $U_\pm$ are
  star-shaped in suitable coordinates, $\tau_\pm$ are homotopic to constants, so $g'$ is
  homotopic to a constant, and then (as $\GL$ is connected) to $\Id$; normalizing by the value
  at $e$ (\cref{bott:lem:normalize}) makes the homotopy pointed. So $[g]=0$.
\end{proof}

\begin{lemma}\label{bott:lem:smash}
  For $m\ge0$, with base point $e\in\Sphere^m$, $\tilde K^{-2}(\Sphere^m)\cong\pi_{m+1}(\Ginf(Y;E))$.
\end{lemma}
\begin{proof}
  An element of $\tilde K^{-2}(\Sphere^m)$ is represented by $f\colon\Sphere^m\times\Sphere\to\Ginf$
  with $f=\Id$ on $\Sphere^m\times\{1\}$ and such that the loop $f(e,\cdot)$ is null-homotopic
  through pointed loops $h_s$, $h_0=f(e,\cdot)$, $h_1=\Id$. Choose $\chi\in\CI(\Sphere^m)$ with
  values in $[0,1]$ and $\chi(e)=1$. Then
  $F_s(x,\theta)=f(x,\theta)f(e,\theta)^{-1}h_{s\chi(x)}(\theta)$ is a homotopy in $K^{-2}$
  from $f$ to a map equal to $\Id$ on $W=\Sphere^m\times\{1\}\cup\{e\}\times\Sphere$. Similarly a
  homotopy $H_t$ between two maps equal to $\Id$ on $W$ is replaced by
  $H_t(x,\theta)H_t(e,\theta)^{-1}$, which is $\Id$ on $W$ and has the same ends. Hence
  $\tilde K^{-2}(\Sphere^m)$ is the set of homotopy classes, relative to $W$, of maps equal to
  $\Id$ on $W$. By \cref{bott:lem:flatten} these are the classes of maps which are $\Id$ near
  $W$, and since $(\Sphere^m\setminus\{e\})\times(\Sphere\setminus\{1\})\cong\R^{m+1}$ this is
  $\Kc^{-1}(\R^{m+1})\cong\pi_{m+1}(\Ginf)$. (For $m=0$, $\Sphere^0=\{e,-e\}$ and the statement
  is $K^{-2}(\{-e\})=\pi_1(\Ginf)$.)
\end{proof}

\begin{corollary}\label{bott:cor:homotopy-groups}
  For every compact $Y$ of positive dimension and every bundle $E$,
  \begin{equation*}
    \pi_k(\Ginf(Y;E))\cong\begin{cases}\Z&k\text{ odd},\\0&k\text{ even}.\end{cases}
  \end{equation*}
  The isomorphism $\pi_1(\Ginf)\cong\Z$ is given by the winding number. Consequently also
  $K^{-1}(\Sphere^k)=[\Sphere^k;\Ginf]=\pi_k(\Ginf)$, and $\tilde K^0(\Sphere^k)$ is $\Z$ for
  $k$ even and $0$ for $k$ odd, $k\ge1$.
\end{corollary}
\begin{proof}
  $\pi_0(\Ginf)=0$ since $\Ginf$ is connected. For $m\ge0$,
  $\pi_{m+1}(\Ginf)\cong\tilde K^{-2}(\Sphere^m)\cong\tilde K^0(\Sphere^m)$ by
  \cref{bott:lem:smash,bott:cor:reduced}. For $m=0$, $\tilde K^0(\Sphere^0)=K^0(\mathrm{pt})=\Z$,
  so $\pi_1(\Ginf)=\Z$, the isomorphism being minus the index, i.e.\ the winding number. For
  $m\ge1$, $\tilde K^0(\Sphere^m)\cong\pi_{m-1}(\Ginf)$ by \cref{bott:lem:sphere-clutching}, so
  $\pi_{m+1}(\Ginf)\cong\pi_{m-1}(\Ginf)$ and the result follows by induction. Free and pointed
  homotopy classes of maps $\Sphere^k\to\Ginf$ agree, since $f\mapsto f(e)^{-1}f$ is a
  homotopy inverse to the inclusion of pointed maps ($\Ginf$ is a connected group).
\end{proof}

This justifies calling $\Ginf$ a classifying space for odd K-theory: its homotopy groups are
those of the stable unitary group (Bott's original theorem).

\section{Isotropic--semiclassical quantization maps}\label{bott:sec:quantization}

The isotropic index map quantizes a loop of matrices, viewed as a function on the
circle at infinity of $\R^2$, to a Fredholm operator. There is a second, `odd', way to
pass from symbols to operators, which works on every $\R^{2n}$: a compactly supported
map $g\colon\R^{2n}\to\GL(N,\C)$ (equal to $\Id$ outside a compact set) is quantized to
an invertible perturbation of the identity by an isotropic smoothing operator on
$\R^n$, after rescaling the symbol by a small parameter $\eps$. The same procedure
applied to projections gives an even version. These maps are used in
\cref{ch:semiclassical-index}; here we define them, show that they commute with
clutching, and show that they realize Bott periodicity.

\subsection{Weyl operators and semiclassical scaling}
On $\R^{2n}$ we use the variables $w=(x,\xi)$ and dual variables $\zeta=(p,q)$, with
$\pair\zeta w=p\cdot x+q\cdot\xi$, and the Fourier transform
$\hat a(\zeta)=\int e^{-i\pair\zeta w}a(w)\,dw$. The \emph{Weyl operators}
\begin{equation*}
  \mathcal W(\zeta)u(x)=e^{ip\cdot x+ip\cdot q/2}\,u(x+q),\qquad\zeta=(p,q)\in\R^{2n},
\end{equation*}
are unitary on $L^2(\R^n)$, strongly continuous in $\zeta$, and by direct computation
\begin{equation}\label{bott:eq:weyl-ops}
  \mathcal W(\zeta)\mathcal W(\zeta')=e^{i\sigma(\zeta,\zeta')/2}\,\mathcal W(\zeta+\zeta'),\qquad
  \sigma\bigl((p,q),(p',q')\bigr)=p'\cdot q-p\cdot q'.
\end{equation}
From the kernel of the Weyl quantization $\OpW$ of \cref{ch:isotropic} one finds
$\OpW(e^{i\pair\zeta\cdot})=\mathcal W(\zeta)$: its kernel is
$e^{ip\cdot(x+y)/2}\delta(x-y+q)$. Since $u,v\mapsto\pair{\OpW(a)u}{v}$ is the pairing of
$a$ with a Schwartz function of $w$ (the kernel formula of the Weyl quantization), we
may insert the Fourier inversion formula $a=(2\pi)^{-2n}\int\hat a(\zeta)e^{i\pair\zeta\cdot}d\zeta$
and get, for $a\in\cS(\R^{2n})$,
\begin{equation}\label{bott:eq:weyl-superposition}
  \OpW(a)=(2\pi)^{-2n}\int\hat a(\zeta)\,\mathcal W(\zeta)\,d\zeta,\qquad
  \norm{\OpW(a)}_{\Lin(L^2)}\le(2\pi)^{-2n}\norm{\hat a}_{L^1},
\end{equation}
the integral converging strongly. The same holds for matrix-valued symbols, with
$\norm{\hat a}_{L^1}=\int\abs{\hat a(\zeta)}\,d\zeta$ and the operator norm on matrices.

\begin{definition}\label{bott:def:sl-quantization}
  For $0<\eps\le1$ and $a\in\cS(\R^{2n};M(N,\C))$ set
  \begin{equation*}
    q_\eps(a)=\OpW(a_\eps)\in\Psiiso{-\infty}(\R^n;\C^N),\qquad a_\eps(w)=a(\sqrt\eps\,w).
  \end{equation*}
\end{definition}

Conjugating by the unitary dilation $u\mapsto\eps^{-n/4}u(\eps^{-1/2}\cdot)$ transforms
$q_\eps(a)$ into the semiclassical Weyl quantization of $a$ with parameter $\eps$,
$(2\pi\eps)^{-n}\int e^{i(x-y)\cdot\xi/\eps}a(\frac{x+y}2,\xi)u(y)\,dy\,d\xi$, so $q_\eps$ is
the isotropic version of the semiclassical quantization of \cref{ch:semiclassical}.
The dilations form a smooth path of unitaries joining this conjugation to the identity
and preserving $\Psiiso{-\infty}$, so for the K-theoretic constructions below it makes
no difference which of the two is used.

\begin{lemma}\label{bott:lem:sl-product}
  For $a,b\in\cS(\R^{2n};M(N,\C))$ and $0<\eps\le1$:
  \begin{enumerate}[label=(\roman*)]
  \item $q_\eps(a)^*=q_\eps(a^*)$ and $\norm{q_\eps(a)}\le(2\pi)^{-2n}\norm{\hat a}_{L^1}$;
  \item $\norm{q_\eps(a)q_\eps(b)-q_\eps(ab)}\le\eps\,2^{-1}(2\pi)^{-4n}\,
    \norm{\abs\zeta\hat a}_{L^1}\norm{\abs\zeta\hat b}_{L^1}$;
  \item $q_\eps(a)\to a(0)$ strongly on $L^2(\R^n;\C^N)$ as $\eps\downarrow0$.
  \end{enumerate}
  If $a$ depends smoothly on parameters, with $a=0$ outside a compact subset of
  $\R^{2n}\times X$, the bounds are locally uniform in the parameters and $q_\eps(a)$ is a
  smooth function of $(\eps,x)\in(0,1]\times X$ with values in $\Psiiso{-\infty}(\R^n;\C^N)$.
\end{lemma}
\begin{proof}
  The adjoint formula is that of the Weyl quantization. Since
  $\widehat{a_\eps}(\zeta)=\eps^{-n}\hat a(\zeta/\sqrt\eps)$,
  \eqref{bott:eq:weyl-superposition} gives
  $q_\eps(a)=(2\pi)^{-2n}\int\hat a(\zeta)\mathcal W(\sqrt\eps\zeta)\,d\zeta$, which implies
  the norm bound and, by dominated convergence and $\mathcal W(\sqrt\eps\zeta)u\to u$, (iii).
  By \eqref{bott:eq:weyl-ops},
  \begin{equation*}
    q_\eps(a)q_\eps(b)=(2\pi)^{-4n}\iint\hat a(\zeta)\hat b(\zeta')\,
    e^{i\eps\sigma(\zeta,\zeta')/2}\,\mathcal W\bigl(\sqrt\eps(\zeta+\zeta')\bigr)\,d\zeta\,d\zeta',
  \end{equation*}
  while $\widehat{ab}=(2\pi)^{-2n}\hat a*\hat b$ gives the same formula for $q_\eps(ab)$
  without the factor $e^{i\eps\sigma/2}$. Since $\abs{e^{is}-1}\le\abs s$ and
  $\abs{\sigma(\zeta,\zeta')}\le\abs\zeta\abs{\zeta'}$, (ii) follows. The last statement is
  clear from the kernel of $\OpW(a_\eps)$, a Schwartz function depending smoothly on
  $(\eps,x)$.
\end{proof}

\begin{discussion}
  The parameter $\eps$ is a semiclassical parameter. The rescaled symbol
  $a(\sqrt\eps\,w)$ varies on the scale $\eps^{-1/2}$, which is large compared
  with the scale $1$ of the uncertainty principle. On this scale quantization is
  almost multiplicative, $q_\eps(a)q_\eps(b)=q_\eps(ab)+O(\eps)$. Hence an
  invertible symbol is quantized to an almost invertible operator, which can be
  corrected to an invertible one, and an idempotent symbol to an almost
  idempotent operator, whose spectral projection is an honest projection. This
  is all that the two quantization maps below use.
\end{discussion}

\subsection{The odd quantization map}

\begin{proposition}\label{bott:prop:odd-sl}
  Let $X$ be a manifold and $g\in\CI(\R^{2n}\times X;\GL(N,\C))$ with $g=\Id$ outside a
  compact set. There is $\eps_0>0$ such that for $0<\eps\le\eps_0$
  \begin{equation*}
    G_\eps(x)=\Id+q_\eps\bigl(g(\cdot,x)-\Id\bigr)\in\Gres(\R^n;\C^N)\quad\text{for all }x\in X;
  \end{equation*}
  $G_\eps$ is a smooth map, equal to $\Id$ outside a compact subset of $X$. Its class
  $\Indsl{odd}[g]=[G_\eps]\in\Kc^{-1}(X)$ is independent of $\eps\in(0,\eps_0]$, depends only
  on $[g]\in\Kc^{-1}(\R^{2n}\times X)$, and defines a homomorphism
  \begin{equation*}
    \Indsl{odd}\colon\Kc^{-1}(\R^{2n}\times X)\longrightarrow\Kc^{-1}(X),
  \end{equation*}
  natural under proper maps and under extension from open subsets of $X$.
\end{proposition}
\begin{proof}
  Let $\alpha=g-\Id$ and $\beta=g^{-1}-\Id$, both compactly supported, with
  $\alpha+\beta+\alpha\beta=0=\alpha+\beta+\beta\alpha$. Then
  \begin{equation*}
    \bigl(\Id+q_\eps(\alpha)\bigr)\bigl(\Id+q_\eps(\beta)\bigr)
    =\Id+q_\eps(\alpha)q_\eps(\beta)-q_\eps(\alpha\beta)=\Id+O(\eps),
  \end{equation*}
  uniformly in $x$ by \cref{bott:lem:sl-product}(ii), and similarly in the other order. So
  for small $\eps$, $G_\eps(x)$ has a left and a right inverse on $L^2$, i.e.\ is invertible,
  and since it lies in $\Id+\Psiiso{-\infty}$ its inverse does too (\cref{ch:isotropic}), so
  $G_\eps(x)\in\Gres(\R^n;\C^N)$, with $\norm{G_\eps^{-1}}$ uniformly bounded. Smoothness
  and the support property follow from \cref{bott:lem:sl-product}. A homotopy $g_t$ as in
  \cref{bott:def:Kc}, or the path $\eps\in(0,\eps_0]$, gives a homotopy of the $G_\eps$, and
  stabilization of $g$ stabilizes $G_\eps$. For two maps $g,h$,
  $G_\eps(g)G_\eps(h)-G_\eps(gh)=q_\eps(\alpha_g)q_\eps(\alpha_h)-q_\eps(\alpha_g\alpha_h)=O(\eps)$,
  so for small $\eps$ the linear path between $G_\eps(g)G_\eps(h)$ and $G_\eps(gh)$ stays in
  $\Gres$; thus $\Indsl{odd}$ is a homomorphism. Naturality is clear.
\end{proof}

\subsection{Projections}
Even classes are conveniently represented by projections. By \emph{projection data} on
a manifold $M$ we mean a smooth family of orthogonal projections
$p\in\CI(M;M(N,\C))$ which is equal to a constant projection $p_\infty$ outside a
compact set. It defines
\begin{equation}\label{bott:eq:proj-class}
  [p]-[p_\infty]=\bigl[(\Ran p,\ M\times\Ran p_\infty,\ p_\infty|_{\Ran p})\bigr]\in\Kc^0(M);
\end{equation}
$p_\infty|_{\Ran p}$ is the identity outside a compact set. For compact $M$ this is the
difference of the classes of the bundles $\Ran p$ and $\Ran p_\infty$.

\begin{lemma}\label{bott:lem:proj-model}
  \begin{enumerate}[label=(\roman*)]
  \item The class \eqref{bott:eq:proj-class} is unchanged under homotopies $p_t$ of
    projection data with $p_t=p_\infty$ outside one compact set, under stabilization
    $(p,p_\infty)\mapsto(p\oplus r,p_\infty\oplus r)$ by a constant projection $r$, and under
    conjugation by a constant unitary. The same holds for smooth families of idempotents
    (not necessarily self-adjoint).
  \item If $M$ is compact, or is the total space of a real vector bundle over a compact
    manifold (for instance $\R^k\times X$ or $T^*Z$ with $X$, $Z$ compact), every element of
    $\Kc^0(M)$ is of the form \eqref{bott:eq:proj-class}.
  \end{enumerate}
\end{lemma}
\begin{proof}
  (i) For a homotopy, $\Ran p_t$ is a bundle over $[0,1]\times M$; parallel transport in
  $t$ for the connection $p_t\,\partial_t$ (\cref{bott:lem:bundle-homotopy}) is an
  isomorphism $\Ran p_0\to\Ran p_1$ which is the identity where $p_t$ does not depend on
  $t$, hence outside a compact set, and relations (a), (b) of \cref{bott:def:Kc} apply.
  Stabilization is (d) and conjugation is (a). For idempotents use the connection
  $p_t\,\partial_t$ in the same way.

  (ii) Let $(E_+,E_-,a)$ be a triple, $a$ invertible outside a compact set $K$. In the
  second case bundles on $M$ are pulled back from the compact base
  (\cref{bott:lem:bundle-homotopy}); in either case there is a bundle $F$ with
  $E_-\oplus F\cong M\times\C^N$ and an embedding $E_+\oplus F\subset M\times\C^{N'}$
  (\cref{bott:lem:bundles}). By (d) and (a) we may assume $E_-=\C^N$ and
  $E_+\subset\C^{N'}$. Choose $\chi\in\CI(M;[0,\pi/2])$ with $\chi=0$ near $K$ and
  $\chi=\pi/2$ outside a compact set, and let $p$ be the orthogonal projection onto
  $\{(\cos\chi\,e,\sin\chi\,a(e));\ e\in E_+\}\subset\C^{N'}\oplus\C^N$, a bundle isomorphic
  to $E_+$, and $p_\infty$ the projection onto $0\oplus\C^N$. Outside a compact set
  $p=p_\infty$, and in terms of the isomorphism with $E_+$ the map $p_\infty|_{\Ran p}$ is
  $e\mapsto\sin\chi\,a(e)$, which equals $a$ outside a compact set. By (a) and (b) the
  class \eqref{bott:eq:proj-class} is that of $(E_+,\C^N,a)$.
\end{proof}

Conversely, two projection data define the same class if and only if they are related
by the moves in (i) (\cite{Atiyah-K}, Ch.~II); we use this only to regard the even
quantization map below as defined on $\Kc^0$.

In the quantized setting the projections differ from a constant one by a smoothing
operator. For such families the relative class is defined by finite rank
approximation. Here $p_\infty\in M(N,\C)$ acts on $L^2(\R^n;\C^N)$ pointwise, and
commutes with the Hermite projections $\Pi_k$ of $L^2(\R^n;\C^N)$ (onto the span of the
$u_\alpha\otimes\C^N$, $\abs\alpha\le k$).

\begin{lemma}\label{bott:lem:relative-class}
  Let $X$ be a manifold and $Q\in\CI(X;p_\infty+\Psiiso{-\infty}(\R^n;\C^N))$ a family of
  idempotents with $Q=p_\infty$ outside a compact subset of $X$. There is a well-defined
  class $\delta(Q)\in\Kc^0(X)$ with the following properties.
  \begin{enumerate}[label=(\roman*)]
  \item If $V\subset\Ran\Pi_k$ is a $p_\infty$-invariant subspace such that $Q$ commutes
    with the orthogonal projection $\Pi_V$ onto $V$ and $Q=p_\infty$ on $V^\perp$, then
    $\delta(Q)$ is the class of the triple $(\Ran(Q\Pi_V),X\times\Ran(p_\infty\Pi_V),p_\infty)$.
    In particular, for $X$ a point, $\delta(Q)=\rank(Q|_V)-\rank(p_\infty|_V)$.
  \item $\delta(Q_0)=\delta(Q_1)$ if $Q_t$, $t\in[0,1]$, is a smooth family of such idempotents
    with $Q_t=p_\infty$ outside one compact set; $\delta$ is additive under direct sums and
    natural under proper maps.
  \end{enumerate}
\end{lemma}
\begin{proof}
  On the circle $\Gamma=\{\abs{\lambda-1}=\frac12\}$ the resolvent of an idempotent $Q$ is
  $(\lambda-Q)^{-1}=\lambda^{-1}(\Id-Q)+(\lambda-1)^{-1}Q$, of norm at most
  $2(1+2\norm Q)$. Hence if $T$ is an operator with $\norm{T-Q}<(4+8\norm Q)^{-1}$, then
  $\Gamma$ lies in the resolvent set of $T$ and the Riesz projection
  $\mathcal R(T)=(2\pi i)^{-1}\oint_\Gamma(\lambda-T)^{-1}d\lambda$ is an idempotent, equal to
  $Q$ if $T=Q$. If moreover $T=p_\infty+R$ with $R\in\Psiiso{-\infty}$, two applications of
  the resolvent identity give
  \begin{equation*}
    (\lambda-T)^{-1}=(\lambda-p_\infty)^{-1}+(\lambda-p_\infty)^{-1}R(\lambda-p_\infty)^{-1}
    +(\lambda-p_\infty)^{-1}R(\lambda-T)^{-1}R(\lambda-p_\infty)^{-1},
  \end{equation*}
  and $(\lambda-p_\infty)^{-1}$ is a matrix, so by \ref{bott:I:corner} (in $n$ dimensions,
  \cref{ch:isotropic}) $\mathcal R(T)-p_\infty\in\Psiiso{-\infty}$, depending smoothly on
  parameters on which $R$ depends smoothly.

  Now write $Q=p_\infty+R$, $R=0$ outside the compact set $C\subset X$. By
  \cref{bott:rem:iso-model}, $\Pi_kR\Pi_k\to R$ uniformly on $C$, so for $k$ large the
  operators $T_s=p_\infty+(1-s)R+s\Pi_kR\Pi_k$ are within the bound above of $Q$, and
  $Q_s=\mathcal R(T_s)$ is a homotopy of idempotents of the stated kind from $Q$ to
  $Q^{(k)}=\mathcal R(p_\infty+\Pi_kR\Pi_k)$. Since $p_\infty+\Pi_kR\Pi_k$ commutes with
  $\Pi_k$ and equals $p_\infty$ on $\Ran(\Id-\Pi_k)$, so does $Q^{(k)}$, and we define
  $\delta(Q)$ by the triple in (i) for $Q^{(k)}$ and $V=\Ran\Pi_k$. Here $\Ran(Q^{(k)}\Pi_k)$ is
  a smooth subbundle of $X\times\Ran\Pi_k$ (\cref{bott:lem:bundles}) equal to
  $\Ran(p_\infty\Pi_k)$ outside $C$, where $p_\infty$ restricts to the identity.

  We first show: (*) if $Q_t$ is a smooth family as in (ii), and all $Q_t$ commute with
  $\Pi_V$ and equal $p_\infty$ on $V^\perp$, then the triples of (i) for $Q_0$ and $Q_1$ are
  equivalent. Indeed, transport along $t$ for the connection $Q_t\partial_t$ on the bundle
  $\Ran(Q_t\Pi_V)$ over $[0,1]\times X$ is an isomorphism which is the identity outside a
  compact set, and (a), (b) of \cref{bott:def:Kc} apply. Moreover, if $V\subset V'$ with
  both $p_\infty$-invariant, then the triple for $V'$ is the triple for $V$ plus
  $(\Ran(p_\infty\Pi_{V'\ominus V}),\Ran(p_\infty\Pi_{V'\ominus V}),\Id)$, a stabilization.

  Independence of $k$: for $k'\ge k$ both $\Pi_kR\Pi_k$ and $\Pi_{k'}R\Pi_{k'}$ are close to
  $R$, so the Riesz projections of the operators on the segment between
  $p_\infty+\Pi_kR\Pi_k$ and $p_\infty+\Pi_{k'}R\Pi_{k'}$ form a family as in (*) with
  $V=\Ran\Pi_{k'}$, joining $Q^{(k)}$ to $Q^{(k')}$. With the stabilization remark this
  shows that $\delta(Q)$ does not depend on $k$. The same argument, with one $k$ chosen
  uniformly on $[0,1]\times C$, gives (ii). For (i), $V\subset\Ran\Pi_k$, and $Q$ commutes
  with $\Pi_k$ and equals $p_\infty$ on $\Ran\Pi_k\ominus V$; so $Q=Q^{(k)}$ and the triples
  for $V$ and $\Ran\Pi_k$ differ by a stabilization. Additivity and naturality are clear.
\end{proof}

For $X$ a point, $\delta(Q)$ is the relative dimension of the two projections, i.e.\
the index of $p_\infty\colon\Ran Q\to\Ran p_\infty$ when $Q$ is self-adjoint.

\subsection{The even quantization map}

\begin{proposition}\label{bott:prop:even-sl}
  Let $X$ be a manifold and $p$ projection data on $\R^{2n}\times X$ with values in
  $M(N,\C)$ and constant value $p_\infty$. There is $\eps_0>0$ such that for
  $0<\eps\le\eps_0$ the self-adjoint operators
  \begin{equation*}
    P_\eps(x)=p_\infty+q_\eps\bigl(p(\cdot,x)-p_\infty\bigr)
  \end{equation*}
  satisfy $\norm{P_\eps^2-P_\eps}\le\frac18$; their spectral projections
  $Q_\eps=\mathcal R(P_\eps)=\chi_{(1/2,\infty)}(P_\eps)$ form a family as in
  \cref{bott:lem:relative-class}, and
  \begin{equation*}
    \Indsl{even}([p]-[p_\infty])=\delta(Q_\eps)\in\Kc^0(X)
  \end{equation*}
  is independent of $\eps\in(0,\eps_0]$ and invariant under the moves of
  \cref{bott:lem:proj-model}(i). It is additive under direct sums and natural. Via
  \cref{bott:lem:proj-model} it defines a homomorphism
  $\Indsl{even}\colon\Kc^0(\R^{2n}\times X)\to\Kc^0(X)$.
\end{proposition}
\begin{proof}
  With $r=p-p_\infty$, $p^2=p$ means $p_\infty r+rp_\infty+r^2=r$, so
  $P_\eps^2-P_\eps=q_\eps(r)q_\eps(r)-q_\eps(r^2)=O(\eps)$ uniformly in $x$
  (\cref{bott:lem:sl-product}). For a self-adjoint $P$ with $\norm{P^2-P}\le\frac18$ the
  spectrum lies in $\{\abs{\lambda^2-\lambda}\le\frac18\}$, which does not meet
  $\{\abs{\lambda-\frac12}\le\frac12\sqrt{1/2}\}$, so $\Gamma$ is in the resolvent set and
  $\mathcal R(P)$ is the spectral projection of $(1/2,\infty)$. The proof of
  \cref{bott:lem:relative-class} shows that $Q_\eps-p_\infty\in\Psiiso{-\infty}$ depends
  smoothly on $(\eps,x)$, and $Q_\eps=p_\infty$ where $r(\cdot,x)=0$. Homotopies of $p$,
  and the parameter $\eps$, give homotopies of $Q_\eps$, so (ii) of
  \cref{bott:lem:relative-class} gives the invariance; stabilization and constant unitary
  conjugation pass through the construction.
\end{proof}

\subsection{Clutching}
Fix $\vartheta\in\CI(\R;[0,\pi/2])$ with
$\vartheta=0$ on $(-\infty,-1]$ and $\vartheta=\pi/2$ on $[1,\infty)$. For an invertible
element $g$ of a unital algebra and $\tau\in\R$ let
\begin{equation*}
  S(\tau,g)=\begin{pmatrix}\cos\tau&-\sin\tau\,g^{-1}\\\sin\tau\,g&\cos\tau\end{pmatrix},
  \qquad S(\tau,g)^{-1}=S(-\tau,g),
\end{equation*}
and, with $\Pi_\infty=\diag(\Id,0)$,
\begin{equation}\label{bott:eq:clutch-proj}
  \Pi_g(t)=S(\vartheta(t),\Id)^{-1}S(\vartheta(t),g)\,\Pi_\infty\,S(\vartheta(t),g)^{-1}S(\vartheta(t),\Id),
  \qquad t\in\R.
\end{equation}
This is an idempotent; since
$S(\tau,g)\Pi_\infty S(\tau,g)^{-1}=\begin{pmatrix}\cos^2\tau&\cos\tau\sin\tau\,g^{-1}\\
\cos\tau\sin\tau\,g&\sin^2\tau\end{pmatrix}$, its entries are affine functions of $g$ and
$g^{-1}$ with coefficients depending on $t$, and a direct computation gives
$\Pi_g(t)=\Pi_\infty$ for $\abs t\ge1$ and for $g=\Id$. If $g$ is unitary, $S(\tau,g)$ is
unitary and $\Pi_g(t)$ is an orthogonal projection.

\begin{lemma}\label{bott:lem:sl-clutching}
  \begin{enumerate}[label=(\roman*)]
  \item For a manifold $M$ and $g\in\CI(M;\mathrm U(N))$ equal to $\Id$ outside a compact set,
    $\Pi_g$ is projection data on $\R\times M$, and $[\Pi_g]-[\Pi_\infty]$ is the class of
    the triple $(\C^N,\C^N,\tilde g)$, $\tilde g(t,m)=\Id$ for $t\le-1$, $\tilde g(t,m)=g(m)$
    for $t\ge1$. For compact $M$ the map
    $\operatorname{clu}\colon K^{-1}(M)\to\Kc^0(\R\times M)$, $[g]\mapsto[\Pi_g]-[\Pi_\infty]$, is
    the isomorphism of \cref{bott:lem:clutching} ($n=1$, base point $e=-1$).
  \item For $G\in\CI(M;\Gres(\R^n;\C^N))$ equal to $\Id$ outside a compact set,
    \eqref{bott:eq:clutch-proj} with $G$ in place of $g$ is a family of idempotents as in
    \cref{bott:lem:relative-class} over $\R\times M$, and
    $\operatorname{clu}[G]=\delta(\Pi_G)$ extends $\operatorname{clu}$ to the isotropic model
    of $\Kc^{-1}(M)$.
  \item Quantization commutes with clutching: for $g\in\CI(\R^{2n}\times M;\mathrm U(N))$
    equal to $\Id$ outside a compact set,
    \begin{equation*}
      \Indsl{even}\bigl([\Pi_g]-[\Pi_\infty]\bigr)=\operatorname{clu}\bigl(\Indsl{odd}[g]\bigr)
      \in\Kc^0(\R\times M),
    \end{equation*}
    where on the left $\Pi_g$ is regarded as projection data on $\R^{2n}\times(\R\times M)$.
  \end{enumerate}
\end{lemma}
\begin{proof}
  (i) The map $F(t)v=S(\vartheta,\Id)^{-1}S(\vartheta,g)(v,0)$ is a global isomorphism
  $\C^N\to\Ran\Pi_g$, and $\Pi_\infty F(t)v$ is the first component of
  $S(-\vartheta,\Id)(\cos\vartheta\,v,\sin\vartheta\,gv)$, i.e.\
  $(\cos^2\vartheta+\sin^2\vartheta\,g)v$. This is $v$ for $t\le-1$ and $gv$ for $t\ge1$; by
  (a), (b) of \cref{bott:def:Kc} the triple of \eqref{bott:eq:proj-class} is equivalent to
  $(\C^N,\C^N,\tilde g)$, which is the image of $[g]$ under the inverse of the
  isomorphism of \cref{bott:lem:clutching} (every class of $K^{-1}(M)$ has a unitary
  representative, by the polar decomposition $g\mapsto g(g^*g)^{-1/2}$ and the homotopy
  $g(g^*g)^{-s/2}$).

  (ii) $\Pi_G-\Pi_\infty$ is affine in $G-\Id$ and $G^{-1}-\Id$, so smoothing. If $G$ takes
  values in a subgroup $\GL(V)$ as in \cref{bott:rem:iso-model}, $\Pi_G$ is reduced by
  $V\oplus V$ and \cref{bott:lem:relative-class}(i) identifies $\delta(\Pi_G)$ with the
  class in (i); as every class has such a representative and $\delta$ is homotopy
  invariant, $\delta(\Pi_G)$ depends only on $[G]$.

  (iii) Write $\Pi_g=\Pi_\infty+L_t(g-\Id)+L'_t(g^{-1}-\Id)$, where $L_t$, $L'_t$ place their
  argument, multiplied by functions of $t$, into blocks of a $2\times2$ matrix and then
  conjugate by the constant (in $w$) matrix $S(\vartheta(t),\Id)$. The quantization $q_\eps$
  acts in the variable $w\in\R^{2n}$ and commutes with $L_t$, $L'_t$. With
  $\alpha=g-\Id$, $\beta=g^{-1}-\Id$ and $G_\eps=\Id+q_\eps(\alpha)$,
  \begin{equation*}
    P_\eps=\Pi_\infty+L_t(q_\eps(\alpha))+L'_t(q_\eps(\beta)),\qquad
    \Pi_{G_\eps}=\Pi_\infty+L_t(q_\eps(\alpha))+L'_t(G_\eps^{-1}-\Id).
  \end{equation*}
  As in the proof of \cref{bott:prop:odd-sl},
  $G_\eps^{-1}-\Id-q_\eps(\beta)=-G_\eps^{-1}\bigl(q_\eps(\alpha)q_\eps(\beta)-q_\eps(\alpha\beta)\bigr)=O(\eps)$,
  so $\norm{P_\eps-\Pi_{G_\eps}}=O(\eps)$ uniformly. For small $\eps$ the Riesz projections
  of the segment from $P_\eps$ to the idempotent $\Pi_{G_\eps}$ form a homotopy as in
  \cref{bott:lem:relative-class}(ii) from $Q_\eps$ to $\Pi_{G_\eps}$. Hence
  $\Indsl{even}([\Pi_g]-[\Pi_\infty])=\delta(Q_\eps)=\delta(\Pi_{G_\eps})=\operatorname{clu}[G_\eps]$.
\end{proof}

\subsection{Bott periodicity}
We identify the phase plane $\R^2$ (for $n=1$) with $\C$ by $w=(x,\xi)\mapsto x+i\xi$, and
write $z=(x+i\xi)/\abs w$ as in \cref{bott:sec:isotropic}. Let
\begin{equation}\label{bott:eq:biso}
  b_{\mathrm{iso}}=\bigl[(\C,\C,z)\bigr]\in\Kc^0(\R^2),\qquad
  b_{\mathrm{iso}}\cdot[E]=\bigl[(\varpi^*E,\varpi^*E,z)\bigr]\in\Kc^0(\R^2\times X)
\end{equation}
for a bundle $E$ over $X$. Since $(\varpi^*E,\varpi^*E,z)\oplus(\varpi^*E,\varpi^*E,\bar z)$
is equivalent to $0$ (by the rotation homotopy of \cref{bott:lem:rotation} joining
$\diag(z,\bar z)$ to $\Id$), $b_{\mathrm{iso}}\cdot[E]$ is minus the image of $[E]$ under the
isomorphism of \cref{bott:thm:bott}; so $[E]\mapsto b_{\mathrm{iso}}\cdot[E]$ extends to an
isomorphism $K^0(X)\to\Kc^0(\R^2\times X)$. Note that $b_{\mathrm{iso}}$ is the clutching of the
symbol $z$ of the operator $\ann$ of index one.

\begin{proposition}\label{bott:prop:sl-bott-one}
  Let $X$ be compact and $n=1$. For every vector bundle $E$ over $X$,
  $\Indsl{even}(b_{\mathrm{iso}}\cdot[E])=[E]$. Consequently
  $\Indsl{even}\colon\Kc^0(\R^2\times X)\to K^0(X)$ is the inverse of the isomorphism
  $[E]\mapsto b_{\mathrm{iso}}\cdot[E]$.
\end{proposition}
\begin{proof}
  \emph{Projection data for $b_{\mathrm{iso}}$.} Let $\varphi\in\CI([0,\infty);[0,\pi/2])$
  with $\varphi=0$ on $[0,1]$ and $\varphi=\pi/2$ on $[2,\infty)$, and let $P_B(w)$ be the
  orthogonal projection onto the line spanned by the unit vector
  $v(w)=(\cos\varphi(\abs w),\sin\varphi(\abs w)\,z)\in\C^2$, smooth since $\varphi=0$ near
  $0$; $P_B=p_\infty=\diag(0,1)$ for $\abs w\ge2$. The section $v$ trivializes $\Ran P_B$, and
  $p_\infty v=(0,\sin\varphi\,z)$, which is $z$ times the basis vector of $\Ran p_\infty$ for
  $\abs w\ge2$. So $[P_B]-[p_\infty]=b_{\mathrm{iso}}$. If $E=\Ran e$ for a smooth family of
  orthogonal projections $e$ on $\C^L$, then $\tilde P=P_B\otimes e+p_\infty\otimes(\Id-e)$
  represents $b_{\mathrm{iso}}\cdot[E]$, as one sees by splitting $\C^L=E\oplus E^\perp$.

  \emph{Charge.} Let $H=D_x^2+x^2$ on $\R$. For $\zeta=(p,q)$ one computes
  $[D_x,\mathcal W(\zeta)]=p\mathcal W(\zeta)$, $[x,\mathcal W(\zeta)]=-q\mathcal W(\zeta)$, and
  from the kernels $\OpW(\xi e^{i\pair\zeta\cdot})=\frac12(D_x\mathcal W+\mathcal WD_x)$,
  $\OpW(xe^{i\pair\zeta\cdot})=\frac12(x\mathcal W+\mathcal Wx)$. Hence
  $[H,\mathcal W(\zeta)]=\OpW\bigl(2(p\xi-qx)e^{i\pair\zeta\cdot}\bigr)$, and by superposition
  \begin{equation*}
    [H,\OpW(c)]=\OpW(2i\,\partial_\theta c),\qquad\partial_\theta=x\partial_\xi-\xi\partial_x,\qquad
    c\in\cS(\R^2).
  \end{equation*}
  If $c=f(\abs w)z^m$ (with $f$ vanishing near $0$ if $m\ne0$), then $2i\partial_\theta c=-2mc$, so
  $H\OpW(c)=\OpW(c)(H-2m)$ and $\OpW(c)$ maps $u_k$ into $\C u_{k-m}$ ($u_j=0$ for $j<0$). The
  entries of $P_B-p_\infty$ are of this form, with $m=0$ on the diagonal, $m=1$ in the
  lower left and $m=-1$ in the upper right entry, and dilation does not change $m$.
  Hence $P_\eps$, $p_\infty$ and therefore $Q_\eps$ are reduced by the subspaces
  \begin{equation*}
    B_0=\C\,u_0\otimes e_1,\qquad B_k=\operatorname{span}\{u_k\otimes e_1,u_{k-1}\otimes e_2\},\ k\ge1,
  \end{equation*}
  whose orthogonal sum is $L^2(\R;\C^2)$.

  \emph{Counting.} For $\eps\le\eps_0$ the Riesz projection $Q_\eps$ depends continuously on
  $\eps$ in norm, so $\rank(Q_\eps|_{B_k})$ is independent of $\eps$. By
  \cref{bott:lem:sl-product}(iii), $P_\eps\to p_\infty+(P_B(0)-p_\infty)=\diag(1,0)$ strongly,
  hence in norm on each $B_k$, so $Q_\eps|_{B_k}$ tends to the projection onto
  $u_k\otimes e_1$, of rank one, for every $k\ge0$. On the other hand
  $\rank(p_\infty|_{B_0})=0$ and $\rank(p_\infty|_{B_k})=1$ for $k\ge1$. Since $Q_\eps-p_\infty$
  is compact, $\norm{(Q_\eps-p_\infty)|_{V_K^\perp}}<\frac12$ for $V_K=B_0\oplus\dots\oplus B_K$ and
  $K$ large; the Riesz projections of $p_\infty+(Q_\eps-p_\infty)(\Pi_{V_K}+s(\Id-\Pi_{V_K}))$,
  $s\in[0,1]$, join $Q_\eps$ to an idempotent which equals $Q_\eps$ on $V_K$ and $p_\infty$ on
  $V_K^\perp$. As $V_K\subset\Ran\Pi_K$ is $p_\infty$-invariant,
  \cref{bott:lem:relative-class} gives $\delta(Q_\eps)=(K+1)-K=1$.

  \emph{Twisting.} $\tilde P_\eps=P_\eps\otimes e+p_\infty\otimes(\Id-e)$ commutes with
  $\Id\otimes e$, so its Riesz projection is $Q_\eps\otimes e+p_\infty\otimes(\Id-e)$. Applying
  the same homotopy in the first factor and \cref{bott:lem:relative-class}(i) with
  $V=V_K\otimes\C^L$, $\delta$ is the class of the bundle
  $\Ran(Q_\eps|_{V_K})\otimes E\oplus\Ran(p_\infty|_{V_K})\otimes E^\perp$ minus that of
  $\Ran(p_\infty|_{V_K})\otimes\C^L$, that is $(K+1)[E]-K[E]=[E]$.

  The last statement follows from \cref{bott:thm:bott}, by which every class of
  $\Kc^0(\R^2\times X)$ is $b_{\mathrm{iso}}\cdot y$ for a unique $y\in K^0(X)$, and additivity.
\end{proof}

\begin{theorem}\label{bott:thm:sl-iso}
  For every manifold $X$ and every $n\ge1$ the maps
  \begin{equation*}
    \Indsl{odd}\colon\Kc^{-1}(\R^{2n}\times X)\to\Kc^{-1}(X),\qquad
    \Indsl{even}\colon\Kc^0(\R^{2n}\times X)\to\Kc^0(X)
  \end{equation*}
  are isomorphisms. The inverse of $\Indsl{even}$ is exterior multiplication (as in
  \cref{ch:families}) by $b_{\mathrm{iso}}^{\boxtimes n}\in\Kc^0(\R^{2n})$, the coordinates of
  $\R^{2n}$ being paired as $(x_j,\xi_j)$; the maps commute with clutching (\cref{bott:lem:sl-clutching}), and
  quantizing in stages, $\R^{2n}=\R^{2n'}\times\R^{2n''}$, gives the same maps.
\end{theorem}
\begin{proof}[Sketch of proof]
  For $n=1$ and compact $X$ the even statement is \cref{bott:prop:sl-bott-one}. The
  extension to non-compact $X$ uses the index isomorphism with compact supports
  (\cref{bott:ex:noncompact}) and the same computation, which is local in $X$. Given the
  even statement for all $X$, \cref{bott:lem:sl-clutching}(iii) and the clutching
  isomorphism $\Kc^{-1}(M)\cong\Kc^0(\R\times M)$ (\cref{bott:lem:clutching}, extended to
  non-compact $M$ in the same way) give the odd statement. For the stages, the Weyl
  quantization on $\R^{n'+n''}$ of a product $a'(w')a''(w'')$ is
  $\OpW(a')\otimes\OpW(a'')$, and the estimates of \cref{bott:lem:sl-product} show that
  quantizing first in $w''$ (with values in matrices) and then, after a finite rank
  approximation of the result as in \cref{bott:lem:relative-class}, in $w'$, gives a
  family homotopic to the one obtained by quantizing in $w$ at once. Iterating the case
  $n=1$ gives the general case. Complete proofs of the topological ingredients are in
  \cite{Atiyah-K}.
\end{proof}

\section{Complex and symplectic bundles; the Thom isomorphism}\label{bott:sec:thom}

The quantization maps of \cref{bott:sec:quantization} were defined for the trivial
bundle $\R^{2n}\times X$ with its constant symplectic form. In the index theorem they are
needed for a bundle of symplectic vector spaces over $X$, such as $N\oplus N^*$ for a
real vector bundle $N$, and in the classical formulation for complex vector bundles. We
first recall that the two notions are equivalent up to homotopy.

\begin{lemma}\label{bott:lem:complex-symplectic}
  Let $V\to X$ be a real vector bundle of even rank. Hermitian forms are taken linear in
  the first and antilinear in the second variable.
  \begin{enumerate}[label=(\roman*)]
  \item If $J$ is a complex structure on $V$ ($J^2=-\Id$) and $h$ is a Hermitian metric for
    $J$, then $\omega=\Im h$ is a symplectic form on $V$, $\omega(Jv,Jw)=\omega(v,w)$, and
    $h(v,w)=\omega(Jv,w)+i\,\omega(v,w)$.
  \item If $\omega$ is a symplectic form on $V$ and $g$ is a Euclidean metric, define $J'$ by
    $\omega(J'v,w)=g(v,w)$. Then $J'$ is $g$-skew-adjoint, $-J'^2$ is positive definite and
    commutes with $J'$, and $J=(-J'^2)^{-1/2}J'$ is a complex structure with
    $\omega(Jv,Jw)=\omega(v,w)$ for which $h(v,w)=\omega(Jv,w)+i\,\omega(v,w)$ is a Hermitian metric.
  \end{enumerate}
  Both constructions depend smoothly on parameters, and each followed by the other (with
  $g=\Re h$ in (ii)) returns the structure one started with. Since the choices of $h$ and of
  $g$ range over convex sets, homotopy classes of complex structures and of symplectic
  structures on $V$ correspond bijectively; if the structures are constant outside a
  compact set (in a trivialization), so is the correspondence.
\end{lemma}
\begin{proof}
  (i) From $h(Jv,w)=ih(v,w)$ we get $\omega(Jv,w)=\Im(ih(v,w))=\Re h(v,w)$, which is symmetric
  and positive definite; so $h=\omega(J\cdot,\cdot)+i\omega$. The form $\omega$ is real and
  antisymmetric because $h(w,v)=\overline{h(v,w)}$, non-degenerate because
  $\omega(Jv,v)=h(v,v)>0$ for $v\ne0$, and $J$-invariant because $h(Jv,Jw)=h(v,w)$.

  (ii) $J'$ is invertible since $\omega$ and $g$ are non-degenerate. Using the definition twice,
  $g(v,J'w)=\omega(J'v,J'w)=-\omega(J'w,J'v)=-g(w,J'v)=-g(J'v,w)$, so $J'$ is skew-adjoint and
  $g(-J'^2v,v)=g(J'v,J'v)>0$. Let $A=(-J'^2)^{1/2}$, the positive square root, which commutes
  with $J'$ and is self-adjoint; then $J=A^{-1}J'$ satisfies $J^2=A^{-2}J'^2=-\Id$, and
  $\omega(Jv,w)=\omega(J'A^{-1}v,w)=g(A^{-1}v,w)$ is symmetric and positive definite. Replacing $w$
  by $Jw$ and using the skew-adjointness of $J'$ one checks $\omega(Jv,Jw)=\omega(v,w)$, so $h$ is
  Hermitian with $h(Jv,w)=ih(v,w)$. The square root of a positive definite matrix is a
  smooth (indeed real analytic) function of the matrix, which gives smoothness in
  parameters. If $(J,h)$ is given and $\omega=\Im h$, $g=\Re h$, then $\omega(J'v,w)=g(v,w)=\omega(Jv,w)$
  forces $J'=J$, $A=\Id$; conversely $\Im h=\omega$ in (ii).
\end{proof}

On $\C$ with $h(v,w)=v\bar w$ and $v=x+i\xi$, $\omega(v,w)=\Im(v\bar w)$ is the form
$d\xi\wedge dx$ used for the phase plane in \cref{ch:isotropic}; thus the complex structure
compatible with the standard symplectic form of $\R^{2n}=\{(x,\xi)\}$ is the
identification $(x,\xi)\mapsto x+i\xi$ used in \cref{bott:sec:quantization}.

\subsection{The Bott class of a complex bundle}
Let $V\to X$ be a complex vector bundle with Hermitian metric and $\varpi\colon V\to X$ the
projection. On $\Lambda^*V$ let $\iota_v$ be the adjoint of exterior multiplication by
$v$, so that $v\wedge\iota_v+\iota_v v\wedge=\abs v^2$. Then
$c(v)=v\wedge\cdot-\iota_v$ maps $\Lambda^{\mathrm{ev}}V_x$ to $\Lambda^{\mathrm{odd}}V_x$ and
$c(v)^*c(v)=\abs v^2$.

\begin{definition}\label{bott:def:bott-class}
  The \emph{Bott class} (or Thom class) of $V$ is
  \begin{equation*}
    \beta_V=\bigl[(\varpi^*\Lambda^{\mathrm{ev}}V,\ \varpi^*\Lambda^{\mathrm{odd}}V,\ c)\bigr]\in\Kc^0(V),
  \end{equation*}
  for compact $X$; for a bundle $E$ over $X$ we write
  $\beta_V\cdot[E]=[(\varpi^*(\Lambda^{\mathrm{ev}}V\otimes E),\varpi^*(\Lambda^{\mathrm{odd}}V\otimes E),c\otimes\Id)]$.
\end{definition}

For $V=X\times\C$, $c(v)1=v$ and $\beta_V\cdot[E]=[(\varpi^*E,\varpi^*E,v)]=b_{\mathrm{iso}}\cdot[E]$,
since $v$ and $v/\abs v$ are homotopic outside the zero section. For $V=V_1\oplus V_2$,
$\Lambda V=\Lambda V_1\otimes\Lambda V_2$ as graded bundles and $c(v_1+v_2)$ is the graded
(Clifford) combination of $c(v_1)$ and $c(v_2)$; so $\beta_{V_1\oplus V_2}$ is the product of
$\varpi_1^*\beta_{V_1}$ and $\varpi_2^*\beta_{V_2}$ in the sense of \cref{ch:families}, and
$\beta_{X\times\C^n}=b_{\mathrm{iso}}^{\boxtimes n}$. The class $\lambda_{\mathcal V}$ used in
\cref{ch:families} for a real bundle $\mathcal V$ is $\beta_{\mathcal V\otimes\C}$.

\begin{theorem}[Thom isomorphism]\label{bott:thm:thom}
  Let $X$ be compact and $V\to X$ a Hermitian complex vector bundle of rank $n$. Then
  \begin{equation*}
    K^0(X)\ni[E]-[F]\longmapsto\beta_V\cdot[E]-\beta_V\cdot[F]\in\Kc^0(V)
  \end{equation*}
  is an isomorphism; denote its inverse by $\operatorname{Thom}_V\colon\Kc^0(V)\to K^0(X)$. If
  $V=X\times\C^n$ with the standard structure, then
  $\operatorname{Thom}_V=\Indsl{even}$, the even quantization map of
  \cref{bott:prop:even-sl} for $\C^n=\R^{2n}$.
\end{theorem}
\begin{proof}[Proof and references]
  For $V=X\times\C$ this is \cref{bott:prop:sl-bott-one} together with \cref{bott:thm:bott}.
  For trivial $V$ of rank $n$ it follows from $\beta_V=b_{\mathrm{iso}}^{\boxtimes n}$ and
  \cref{bott:thm:sl-iso}. For general $V$ choose a complement $V'$, $V\oplus V'\cong X\times\C^m$
  (\cref{bott:lem:bundles}); multiplicativity of the Bott class and the trivial case,
  applied over the total space of $V$ (a non-compact base, as in
  \cref{bott:thm:sl-iso}), show that $\beta_V\cdot$ has a left and a right inverse. We do not
  give the details, which are standard; see \cite{Atiyah-K}, \S2.7.
\end{proof}

\begin{remark}[Thom isomorphism by quantization]\label{bott:rem:thom-quantization}
  For a symplectic bundle $(V,\omega)$ over $X$, $\operatorname{Thom}_V$ is again given by
  fibrewise quantization. Choose a compatible complex structure and Hermitian metric
  (\cref{bott:lem:complex-symplectic}). Unitary frames identify the fibres symplectically with
  $\R^{2n}$, and two such frames differ by an element $U$ of the unitary group
  $\mathrm U(n)\subset\mathrm{Sp}(2n,\R)$. By the symplectic invariance of the Weyl
  quantization (\cref{ch:isotropic}), $\OpW(a\circ U)=M_U^{-1}\OpW(a)M_U$ for a unitary
  operator $M_U$ determined up to a constant factor; since $\abs w^2\circ U=\abs w^2$, $M_U$
  commutes with the harmonic oscillator and preserves its eigenspaces. Hence the
  constructions of \cref{bott:lem:relative-class,bott:prop:even-sl}, carried out in local unitary
  frames, are compatible on overlaps up to conjugation by such operators, which does not
  change the finite rank bundles up to isomorphism. They define
  $\operatorname{Ind}^{\mathrm{even}}_{\mathrm{sl},V}\colon\Kc^0(V)\to K^0(X)$ (and similarly an odd map
  $\Kc^{-1}(V)\to K^{-1}(X)$), and $\operatorname{Ind}^{\mathrm{even}}_{\mathrm{sl},V}(\beta_V\cdot[E])=[E]$ by
  the computation of \cref{bott:prop:sl-bott-one} in each factor; thus
  $\operatorname{Ind}^{\mathrm{even}}_{\mathrm{sl},V}=\operatorname{Thom}_V$. To make the patching
  canonical one uses the Fock space of each fibre, on which $\mathrm U(n)$ acts by a genuine
  representation; we omit the details.
\end{remark}

\section{The analytic index map}\label{bott:sec:analytic-index}

Let $Z$ be a compact manifold and $\varpi\colon T^*Z\to Z$. An elliptic operator
$P\in\Psido{0}(Z;E_+,E_-)$ is Fredholm, and its index
\begin{equation*}
  \ind P=\dim\Nul P-\dim\big(\CI(Z;E_-)/P\CI(Z;E_+)\big)
\end{equation*}
is a homotopy invariant (\cref{ch:index}), hence depends only on $E_\pm$ and the homotopy class
of $\sigma_0(P)\in\CI(S^*Z;\Iso(E_+,E_-))$. The principal symbol, extended homogeneously of degree
$0$ to $T^*Z\setminus0$, defines a triple $(\varpi^*E_+,\varpi^*E_-,\sigma_0(P))$ and so a class in
$\Kc^0(T^*Z)$.

\begin{proposition}\label{bott:prop:analytic-index}
  There is a unique homomorphism $\ind_a\colon\Kc^0(T^*Z)\to\Z$, the \emph{analytic index},
  with $\ind P=\ind_a[(\varpi^*E_+,\varpi^*E_-,\sigma_0(P))]$ for every elliptic
  $P\in\Psido{0}(Z;E_+,E_-)$.
\end{proposition}
\begin{proof}
  Every class is of this form. By \cref{bott:lem:bundle-homotopy}(ii) and (a) a class is
  represented by $(\varpi^*E_+,\varpi^*E_-,a)$ with $E_\pm$ bundles over $Z$ and $a$ an isomorphism
  for $|\zeta|\ge R$ (with respect to some metric). The homotopy (c)
  $a_s(\zeta)=a\big(((1-s)|\zeta|+sR)\zeta/|\zeta|\big)$, $|\zeta|\ge R$, ends at a map homogeneous
  of degree $0$, i.e.\ at the symbol $\sigma$ of some elliptic $P\in\Psido{0}(Z;E_+,E_-)$. Define
  $\ind_a$ of the class to be $\ind P$.

  This does not depend on the choices: by homotopy invariance $\ind P$ depends only on the
  homotopy class of $\sigma$, which takes care of $R$, of (b) and of (c). Under (d), $P$ is
  replaced by $P\oplus\Id_H$, with the same index. Under (a), bundle isomorphisms $\phi_\pm$ of
  $\varpi^*E_\pm$ are homotopic, through $\phi_\pm(t\zeta)$, to the pull-backs of their
  restrictions $\phi^0_\pm$ to the zero section; so $\sigma$ is replaced by a symbol homotopic to
  $\phi^0_-\sigma(\phi^0_+)^{-1}$, which is the symbol of $\phi^0_-P(\phi^0_+)^{-1}$, an operator with
  the same index since $\phi_\pm^0$ are invertible. Additivity follows from
  $\ind(P\oplus P')=\ind P+\ind P'$, and uniqueness from the fact that every class arises from an
  elliptic operator.
\end{proof}

The problem posed by Gelfand around 1960 was to compute this map topologically. The answer of
Atiyah and Singer is to define a second, \emph{topological}, index map by embedding $Z$ into a
Euclidean space (or a sphere), constructing on the larger space an operator with the same index
and with symbol derived from that of $P$, and using the fact, essentially Bott periodicity, that
for suitable operators on Euclidean space the index map is an isomorphism; the two maps are then
shown to coincide. This is taken up in \cref{ch:families}.

\section{Exercises}

\begin{exercise}\label{bott:ex:noncompact}
  Show that the index isomorphism of \cref{bott:thm:index-iso} carries over to non-compact
  manifolds: for every manifold $M$ there is an isomorphism $\Kc^{-2}(M)\cong\Kc^0(M)$.
  (Hint: represent elements of $\Kc^{-2}(M)$ by loops equal to $\Id$ outside $K\times\Sphere$, choose
  a compact manifold with boundary $M'\Subset M$ containing a neighbourhood of $K$, and apply the
  families index over $M'$, observing that the index class comes with a canonical trivialization
  over $M\setminus K$.)
\end{exercise}

\begin{exercise}
  Show directly that the Toeplitz operator $T_{z^k}$ has index $-k$. For a smooth family of
  projections $p$ on $\C^N$ over $X$ show that $pL+(\Id-p)$ is surjective with null space
  the constant sections of $\Ran p$, and compute $\ind(pL+\Id-p)=[\Ran p]$ directly from
  \cref{bott:def:families-index} with $k=0$.
\end{exercise}

\begin{exercise}
  Show that for pointed loops $a,b$ over $X$ the families index satisfies
  $\ind(T_aT_b)=\ind(T_{ab})=\ind(a)+\ind(b)$. (Use \cref{bott:lem:rotation} and the fact that
  $T_aT_b-T_{ab}$ has vanishing symbol.)
\end{exercise}

\begin{exercise}
  Let $a(\theta)=\diag(e^{2i\theta},e^{-i\theta})$. Carry out the steps of the proof of
  \cref{bott:thm:linearization} for $a$ and find the resulting normal form; check that its index
  is $-1$ in accordance with \cref{bott:cor:index-circle}.
\end{exercise}

\begin{exercise}
  Show that $\Kc^0(\R)=0$ and $\Kc^0(\R^2)\cong\Z$, generated by the Bott class
  $[(\C,\C,\bar w/|w|)]$.
\end{exercise}

\begin{exercise}
  Let $H\to\Sphere^2=\C\cup\{\infty\}$ be the line bundle obtained by clutching the trivial bundles
  over the two hemispheres with the map $z\mapsto z$ on the equator (this is the tautological
  bundle over $\mathbb{CP}^1$, up to conjugation). Show that $[H]-[\C]$ generates
  $\tilde K^0(\Sphere^2)\cong\Z$, and prove that $(H\otimes H)\oplus\C\cong H\oplus H$ by comparing
  clutching functions and using \cref{bott:lem:rotation}.
\end{exercise}

\begin{exercise}
  Show that if $A\colon X\to\PsiT{0}(\Sphere;\C^N)$ is a smooth family of elliptic Toeplitz
  operators whose null spaces have constant dimension, then $\Nul A$ and $\Nul A^*$ are smooth
  vector bundles and $\ind(A)=[\Nul A]-[\Nul A^*]$. (Use the generalized inverse to construct
  local frames, then \cref{bott:lem:null-complement}.) Prove the same for isotropic families.
\end{exercise}

\begin{exercise}\label{bott:ex:unitary}
  Let $\mathrm U^{-\infty}_{\mathrm{iso}}(\R^n)=\{\Id+A\in\Gres(\R^n);\ (\Id+A)^{-1}=\Id+A^*\}$.
  Show that for $g=\Id+A\in\Gres(\R^n)$ the operator $(g^*g)^{-1/2}-\Id$ is smoothing
  (write $g^*g=\Id+B$ with $B$ smoothing and self-adjoint, and use the Cauchy integral for
  $\lambda^{-1/2}$ together with the resolvent identities, as in the proof of
  \cref{bott:lem:relative-class}), and deduce that $g\mapsto g(g^*g)^{-s/2}$, $s\in[0,1]$, deforms
  compactly supported maps into $\Gres$ to maps into $\mathrm U^{-\infty}_{\mathrm{iso}}$. Conclude
  that $\Kc^{-1}(X)$ is the set of homotopy classes of compactly supported maps into
  $\mathrm U^{-\infty}_{\mathrm{iso}}(\R^n)$.
\end{exercise}

\begin{exercise}\label{bott:ex:GL-connected}
  Prove directly that $\GL(N,\C)$ is connected, by induction on $N$: an invertible $a$ is joined
  to one with $ae_1=e_1$ by moving $ae_1$ to $e_1$ in $\C^N\setminus\{0\}$ while keeping the other
  columns linearly independent of it. Show also that a loop $a$ in $\GL(N,\C)$ with
  $\wind\det a=0$ is null-homotopic: approximate $a$ by a trigonometric polynomial, deform
  it so that $e_1^*a(\theta)e_1$ never vanishes, and reduce to $\GL(N-1,\C)$. Compare with
  \cref{bott:cor:reduced}, where the stable version is deduced from the index.
\end{exercise}

\begin{exercise}\label{bott:ex:radial-weyl}
  Use the commutation relation $[H,\OpW(c)]=\OpW(2i\partial_\theta c)$ from the proof of
  \cref{bott:prop:sl-bott-one} to show that for $f\in\cS(\R)$ the operator $\OpW(f(x^2+\xi^2))$
  on $L^2(\R)$ is diagonal in the Hermite basis. Compute its eigenvalues for
  $f(s)=e^{-s}$, and show that the projection onto $\C u_0$ has Weyl symbol
  $2e^{-(x^2+\xi^2)}$.
\end{exercise}

\begin{exercise}\label{bott:ex:relative-dimension}
  Let $Q$, $P$ be orthogonal projections on a Hilbert space with $Q-P$ compact. Show that
  $P|_{\Ran Q}\colon\Ran Q\to\Ran P$ is Fredholm, that its index is unchanged under norm
  continuous deformations of $Q$ with $Q-P$ compact, and that it equals
  $\rank(Q|_V)-\rank(P|_V)$ whenever $V$ is a finite dimensional subspace reducing both $P$
  and $Q$ with $Q=P$ on $V^\perp$. Deduce that $\delta(Q)$ in \cref{bott:lem:relative-class} is
  this index when $X$ is a point.
\end{exercise}

\begin{exercise}\label{bott:ex:clutched-bundle}
  For $a\in\CI(X;\GL(N,\C))$, $X$ compact, let $E_a$ be the bundle over $\Sphere\times X$ whose
  smooth sections are the $s\in\CI(\R\times X;\C^N)$ with $s(t+2\pi,x)=a(x)s(t,x)$. Show that
  $E_a$ is a smooth vector bundle, that $E_a\oplus E_{a^{-1}}$ is trivial, and that
  $[E_a]-[\C^N]\in K^0(\Sphere\times X)$ is the image of $\operatorname{clu}[a]\in\Kc^0(\R\times X)$
  (\cref{bott:lem:sl-clutching}) under extension from $(0,2\pi)\times X$, up to sign.
\end{exercise}

\section*{Notes}
The linearization of loops (\cref{bott:sec:linearization}) is the heart of Atiyah's
proof of Bott periodicity via clutching functions \cite{Atiyah-K}; the use of the families index
of Toeplitz operators to invert the Bott map follows the source lectures and is closely related to
the treatment of Bott periodicity through the Grassmannian of a polarized Hilbert space in
\cite{Pressley-Segal}. Compared to the source we have replaced the residue argument for linear
loops by the Riesz spectral projection (\cref{bott:lem:linear}), proved that a normal form loop is
trivial when the ranges of $\pi_\pm$ are isomorphic in a single lemma
(\cref{bott:lem:normal-trivial}), defined the families index through arbitrary smooth complements
of the range, corrected the sign of the Bott map (which, with $\ind U=-1$, sends
$[E_+]-[E_-]$ to the class with $p_+z^{-1}$), corrected the second curve in the deformation of
$\Id-P_0+P_0hP_0$ and normalized it at $1\in\Sphere$, and derived the homotopy groups of $\Ginf$
from the index isomorphism for compact spheres together with the clutching construction, instead
of from the (unproved) non-compact index isomorphism.

\Cref{bott:sec:isotropic,bott:sec:quantization,bott:sec:thom} and
\cref{bott:lem:circle-split} come from the K-theory chapter of the 2008 revision of the
lectures (\emph{From microlocal to global analysis}), which starts from the residual
isotropic group as the model for odd K-theory (\cref{bott:rem:iso-model}). There the injectivity
of the isotropic index map is proved by Atiyah's rotation argument, which we only sketch
(\cref{bott:rem:rotation-trick}); instead we combine the isotropic index with the linearization
of loops, so that the only analytic input is the normalized annihilation operator
(\cref{bott:lem:annihilation}). We orient the circle at infinity counterclockwise in the
$(x,\xi)$-plane, so that $\ind A=+\wind\sigma_0(A)$; the source writes $-\wind$, with the
opposite orientation. The semiclassical quantization maps are only stated in the source;
we have defined them through the rescaled Weyl quantization, proved the basic estimate with
the Weyl operators (\cref{bott:lem:sl-product}), given the relative class of projections
(\cref{bott:lem:relative-class}), proved that quantization commutes with clutching using a
simplified clutching projection \eqref{bott:eq:clutch-proj}, and computed the even map for $n=1$
explicitly by means of the rotation invariance of the harmonic oscillator
(\cref{bott:prop:sl-bott-one}); the general case and the Thom isomorphism for non-trivial
bundles are only sketched, with reference to \cite{Atiyah-K}. The source's sections on Chern
forms, the Chern character and the Todd class are incorporated in \cref{ch:families}; its
sections on stabilization, delooping and looping sequences, $C^*$-algebras and the norm closure
of $\Psido{0}$ are empty and have been omitted. The isotropic proof of Bott periodicity via the
index of the harmonic oscillator is the approach of M.~F. Atiyah, \emph{Bott periodicity and the
index of elliptic operators}, Quart. J. Math. Oxford \textbf{19} (1968), 113--140.

The source conjectures that $\GT(\Sphere,Y;E)$ is actually contractible. Since it is an open subset
of an affine Fr\'echet space (\cref{bott:lem:open}), it has the homotopy type of a CW complex by a
theorem of R.~Palais (Homotopy theory of infinite dimensional manifolds, Topology 5 (1966)), and a
weakly contractible CW complex is contractible by Whitehead's theorem; we have not checked the
details of this argument in the present setting.
